% (demande de stage) — Français 1re Tech — Cours signé OnBuch+
% Programme officiel MINESEC. Compiler avec : tectonic N10-demande-stage.tex
\newcommand{\DOCMATIERE}{Français}
\newcommand{\DOCNIVEAU}{1re Tech}
\newcommand{\DOCMODULE}{Français – classe de Première technique}
\newcommand{\DOCLECON}{N10}
\newcommand{\DOCTITRE}{(demande de stage)}
% OnBuch+ — préambule « cours signés » ENRICHI (compilé avec tectonic / XeLaTeX)
% Système visuel « Pop » d'OnBuch+ : fond crème (jamais blanc), bordures encre
% épaisses, ombres « dures » décalées, titres Archivo Black en capitales.
% Version MATHÉMATIQUES : graphes (pgfplots/tikz), boîtes pédagogiques
% (définition, propriété, méthode, attention, le savais-tu, exercice résolu,
% exercice, TP, bilan), page de garde et signature OnBuch+ ancrée en bas.
%
% Macros attendues dans main.tex AVANT \input{preamble} :
%   \DOCMATIERE  \DOCNIVEAU  \DOCMODULE  \DOCLECON (numéro)  \DOCTITRE
\documentclass[11pt]{article}

\usepackage{fontspec}
\usepackage[french]{babel}
\usepackage[a4paper,margin=2cm,top=3.4cm,bottom=2.8cm,headheight=1.4cm,footskip=1.3cm]{geometry}
\usepackage{amsmath,amssymb,amsfonts}
\usepackage{mathtools} % \xmapsto, \coloneqq... souvent utilisés par les modèles
\usepackage{cancel} % simplifications barrées (bilans de forces)
% \abs{z} (module, valeur absolue) et \norm{u} (norme), utilisés par les modèles
\providecommand{\abs}{}\let\abs\relax\DeclarePairedDelimiter{\abs}{\lvert}{\rvert}
\providecommand{\norm}{}\let\norm\relax\DeclarePairedDelimiter{\norm}{\lVert}{\rVert}
\usepackage[table]{xcolor} % [table] : fournit aussi \rowcolors (alternance de lignes)
\usepackage{pagecolor}
\usepackage{tikz}
\usetikzlibrary{arrows.meta,positioning,calc,shapes.geometric,shapes.misc,shapes.arrows,decorations.pathmorphing,decorations.pathreplacing,decorations.markings,patterns,shadows,mindmap,trees,fit,backgrounds,matrix,intersections,angles,quotes}
\usepackage{pgfplots}
\usepackage[version=4]{mhchem} % équations nucléaires \ce{^{235}_{92}U}
\usepackage[european,siunitx]{circuitikz} % schémas électriques
\pgfplotsset{compat=1.18}
\usepgfplotslibrary{fillbetween,groupplots} % aires entre courbes (calcul intégral)
\usepackage{enumitem}
\usepackage{fancyhdr}
% [most] charge toutes les bibliothèques tcolorbox (listings, minted, raster,
% documentation...) et gonfle inutilement la mémoire TeX sur un document long
% (~300 boîtes) : on ne charge que ce qui est réellement utilisé.
\usepackage[skins,breakable]{tcolorbox}
\usepackage{graphicx}
\usepackage{colortbl}
\usepackage{booktabs}
\usepackage{tabularx}
\usepackage{multirow}
\usepackage{diagbox} % case d'en-tête diagonale des tableaux à double entrée
% Colonnes extensibles centrée / gauche / droite (C, L, R), souvent utilisées par les modèles
\newcolumntype{C}{>{\centering\arraybackslash}X}
\newcolumntype{L}{>{\raggedright\arraybackslash}X}
\newcolumntype{R}{>{\raggedleft\arraybackslash}X}
\usepackage{array}
\usepackage{multicol}
\usepackage{siunitx}
% parse-numbers=false : accepte aussi \SI{2,5 \times 10^{-3}}{...}
\sisetup{locale=FR,output-decimal-marker={,},group-minimum-digits=4,parse-numbers=false}
\DeclareSIUnit\ohm{\ensuremath{\symup{\Omega}}} % U+2126 absent de la police
\DeclareSIUnit\division{div} % réglages d'oscilloscope (ms/div, V/div)
\DeclareSIUnit\tr{tr} % tours (vitesse de rotation, ex. \SI{1500}{\tr\per\minute})
\DeclareSIUnit\molar{M} % concentration molaire (ex. \SI{3}{\molar}, \SI{1}{\micro\molar})
\DeclareSIUnit\an{an} % année (le modèle confond parfois avec \year, un compteur TeX) — ex. \SI{5}{\kilo\meter\per\an}
\DeclareSIUnit\FCFA{FCFA} % franc CFA (ex. \SI{500}{\FCFA\per\kilo\gram})
\DeclareSIUnit\degreeBrix{\textdegree Brix} % teneur en sucre d'un sirop/jus (réfractométrie)
\DeclareSIUnit\month{mois} % le modèle utilise parfois \month, un compteur TeX, comme unité de durée
\DeclareSIUnit\microliter{\micro\liter} % le modèle écrit parfois \microliter en un seul token
\DeclareSIUnit\million{millions} % dénombrement (ex. \SI{15}{\million\per\mL} spermatozoïdes)
\usepackage{qrcode}
% \degree : commande fréquemment utilisée par les modèles pour un angle ou une
% température en dehors de \SI{}{} (non fournie par les paquets chargés).
\providecommand{\degree}{^{\circ}}
% \qed (fin de démonstration), utilisé par les modèles sans amsthm
\providecommand{\qed}{\hfill$\square$}
% \barre : double barre des valeurs interdites dans un tableau de variations (tkz-tab)
\providecommand{\barre}{\ensuremath{\|}}
% environnement proof (amsthm non chargé) utilisé par les modèles
\ifdefined\proof\else\newenvironment{proof}[1][Démonstration]{\par\noindent\textit{#1.}\ }{\qed\par}\fi
\usepackage{lastpage}
\usepackage{hyperref}
\hypersetup{hidelinks}

% ============================================================
% Palette « Pop » OnBuch+ — mêmes hex que la classe P (plus_ui.dart)
% ============================================================
\definecolor{poporange}{HTML}{F59321}
\definecolor{popdark}{HTML}{E07A0C}
\definecolor{popcream}{HTML}{FFF6E8}
\definecolor{popink}{HTML}{1C1714}
\definecolor{poppurple}{HTML}{7A5AE0}
\definecolor{popgreen}{HTML}{2E9E6A}
\definecolor{poppink}{HTML}{E0457B}
\definecolor{popblue}{HTML}{2F7DE1}
\definecolor{popgold}{HTML}{C98A12}
\definecolor{popmuted}{HTML}{8A7F76}
\definecolor{popline}{HTML}{EADFD2}
\definecolor{popgray}{HTML}{8A7F76}
\colorlet{popgrey}{popgray}
% Alias tolérants (le modèle nomme parfois les couleurs différemment)
\colorlet{popwhite}{white}
\colorlet{popblack}{popink}
\colorlet{popred}{poppink}
\colorlet{popyellow}{popgold}
% Teintes claires (fonds de tableaux/encadrés) — le modèle nomme parfois
% les couleurs avec un suffixe L (ex. popblueL) sans les avoir définies.
\colorlet{poporangeL}{poporange!20}
\colorlet{popdarkL}{popdark!20}
\colorlet{poppurpleL}{poppurple!20}
\colorlet{popgreenL}{popgreen!20}
\colorlet{poppinkL}{poppink!20}
\colorlet{popblueL}{popblue!20}
\colorlet{popgoldL}{popgold!20}
\colorlet{popredL}{popred!20}
\colorlet{popgrayL}{popgray!20}
% Teintes claires pour fonds de tableaux / zones
\colorlet{popblueL}{popblue!10!white}
\colorlet{poporangeL}{poporange!14!white}
\colorlet{popgreenL}{popgreen!12!white}
\colorlet{poppurpleL}{poppurple!12!white}
\colorlet{poppinkL}{poppink!12!white}
\colorlet{popgoldL}{popgold!14!white}

\pagecolor{popcream}

% ============================================================
% Typographie
% ============================================================
\defaultfontfeatures{Ligatures=TeX}
\setmainfont{PlusJakartaSans-Medium.ttf}[Path=../../../fonts/,BoldFont=PlusJakartaSans-Bold.ttf]
% Maths en sans-serif (Fira Math) avec chiffres et lettres droites de Plus
% Jakarta, pour que les formules se fondent dans le texte.
\usepackage[math-style=ISO,bold-style=ISO]{unicode-math}
\setmathfont{FiraMath-Regular.otf}[Path=../../../fonts/]
\setmathfont{PlusJakartaSans-Medium.ttf}[Path=../../../fonts/,range={up/num,up/{Latin,latin},\mathup}]
\usepackage{icomma} % virgule décimale sans espace en mode math
% Caractères mathématiques Unicode absents des polices de texte (Plus Jakarta,
% Archivo Black) : rendus par leur équivalent mathématique, en texte comme en
% formule — sinon ils s'affichent en carré vide (ex. « Arithmétique dans ℤ »).
\usepackage{newunicodechar}
\newunicodechar{ℝ}{\ensuremath{\mathbb{R}}}
\newunicodechar{ℤ}{\ensuremath{\mathbb{Z}}}
\newunicodechar{ℂ}{\ensuremath{\mathbb{C}}}
\newunicodechar{ℕ}{\ensuremath{\mathbb{N}}}
\newunicodechar{ℚ}{\ensuremath{\mathbb{Q}}}
\newunicodechar{𝔻}{\ensuremath{\mathbb{D}}}
\newunicodechar{α}{\ensuremath{\alpha}}
\newunicodechar{β}{\ensuremath{\beta}}
\newunicodechar{γ}{\ensuremath{\gamma}}
\newunicodechar{δ}{\ensuremath{\delta}}
\newunicodechar{ε}{\ensuremath{\varepsilon}}
\newunicodechar{θ}{\ensuremath{\theta}}
\newunicodechar{λ}{\ensuremath{\lambda}}
\newunicodechar{μ}{\ensuremath{\mu}}
\newunicodechar{σ}{\ensuremath{\sigma}}
\newunicodechar{τ}{\ensuremath{\tau}}
\newunicodechar{φ}{\ensuremath{\varphi}}
\newunicodechar{ψ}{\ensuremath{\psi}}
\newunicodechar{ω}{\ensuremath{\omega}}
\newunicodechar{Σ}{\ensuremath{\Sigma}}
% majuscules produites par \MakeUppercase dans les titres (α-aminés -> Α)
\newunicodechar{Α}{\ensuremath{\alpha}}
\newunicodechar{Β}{\ensuremath{\beta}}
\newunicodechar{Γ}{\ensuremath{\Gamma}}
\newunicodechar{Δ}{\ensuremath{\Delta}}
\newunicodechar{Ε}{\ensuremath{\varepsilon}}
\newunicodechar{Θ}{\ensuremath{\Theta}}
\newunicodechar{Λ}{\ensuremath{\Lambda}}
\newunicodechar{Μ}{\ensuremath{\mu}}
\newunicodechar{Τ}{\ensuremath{\tau}}
\newunicodechar{Φ}{\ensuremath{\Phi}}
\newunicodechar{Ψ}{\ensuremath{\Psi}}
\newunicodechar{Ω}{\ensuremath{\Omega}}
\newunicodechar{↦}{\ensuremath{\mapsto}}
\newunicodechar{∈}{\ensuremath{\in}}
\newunicodechar{∉}{\ensuremath{\notin}}
\newunicodechar{∧}{\ensuremath{\wedge}}
\newunicodechar{⇔}{\ensuremath{\Leftrightarrow}}
\newunicodechar{⟺}{\ensuremath{\Longleftrightarrow}}
\newunicodechar{⇒}{\ensuremath{\Rightarrow}}
\newunicodechar{⟹}{\ensuremath{\Longrightarrow}}
\newunicodechar{∘}{\ensuremath{\circ}}
\newunicodechar{∪}{\ensuremath{\cup}}
\newunicodechar{∩}{\ensuremath{\cap}}
\newunicodechar{≡}{\ensuremath{\equiv}}
\newunicodechar{⊂}{\ensuremath{\subset}}
\newunicodechar{⊥}{\ensuremath{\perp}}
\newunicodechar{∃}{\ensuremath{\exists}}
\newunicodechar{∀}{\ensuremath{\forall}}
\newunicodechar{∛}{\ensuremath{\sqrt[3]{\,}}}
\newunicodechar{∜}{\ensuremath{\sqrt[4]{\,}}}
\newunicodechar{⃗}{\ensuremath{{}^{\to}}}
\newunicodechar{ⁿ}{\ifmmode^{n}\else\textsuperscript{\ensuremath{n}}\fi}
\newunicodechar{ˣ}{\ifmmode^{x}\else\textsuperscript{\ensuremath{x}}\fi}
\newunicodechar{ᵇ}{\ifmmode^{b}\else\textsuperscript{\ensuremath{b}}\fi}
\newunicodechar{ʳ}{\ifmmode^{r}\else\textsuperscript{\ensuremath{r}}\fi}
\newunicodechar{ᵘ}{\ifmmode^{u}\else\textsuperscript{\ensuremath{u}}\fi}
\newunicodechar{ʸ}{\ifmmode^{y}\else\textsuperscript{\ensuremath{y}}\fi}
\newunicodechar{ᵃ}{\ifmmode^{a}\else\textsuperscript{\ensuremath{a}}\fi}
\newunicodechar{ᵉ}{\ifmmode^{e}\else\textsuperscript{\ensuremath{e}}\fi}
\newunicodechar{ᵏ}{\ifmmode^{k}\else\textsuperscript{\ensuremath{k}}\fi}
\newunicodechar{ᵖ}{\ifmmode^{p}\else\textsuperscript{\ensuremath{p}}\fi}
\newunicodechar{ᴺ}{\ifmmode^{N}\else\textsuperscript{\ensuremath{N}}\fi}
\newunicodechar{ᶜ}{\ifmmode^{c}\else\textsuperscript{\ensuremath{c}}\fi}
\newunicodechar{ʰ}{\ifmmode^{h}\else\textsuperscript{\ensuremath{h}}\fi}
\newunicodechar{ᵠ}{\ifmmode^{q}\else\textsuperscript{\ensuremath{q}}\fi}
\newunicodechar{ₙ}{\ifmmode_{n}\else\textsubscript{\ensuremath{n}}\fi}
\newunicodechar{ₐ}{\ifmmode_{a}\else\textsubscript{\ensuremath{a}}\fi}
\newunicodechar{ᵢ}{\ifmmode_{i}\else\textsubscript{\ensuremath{i}}\fi}
\newunicodechar{ᵣ}{\ifmmode_{r}\else\textsubscript{\ensuremath{r}}\fi}
\newunicodechar{ₖ}{\ifmmode_{k}\else\textsubscript{\ensuremath{k}}\fi}
\newunicodechar{ᵦ}{\ifmmode_{\beta}\else\textsubscript{\ensuremath{\beta}}\fi}
\newunicodechar{ₓ}{\ifmmode_{x}\else\textsubscript{\ensuremath{x}}\fi}
\newfontfamily\popdisplay{ArchivoBlack-Regular.ttf}[Path=../../../fonts/]
\newfontfamily\poplogo{SpaceGrotesk-Bold.ttf}[Path=../../../fonts/]
\newfontfamily\popsemi{PlusJakartaSans-SemiBold.ttf}[Path=../../../fonts/]
\newfontfamily\popxbold{PlusJakartaSans-ExtraBold.ttf}[Path=../../../fonts/]
\newfontfamily\popxboldls{PlusJakartaSans-ExtraBold.ttf}[Path=../../../fonts/,LetterSpace=3.0]

\setlist[enumerate]{leftmargin=1.7em,itemsep=5pt,topsep=4pt}
\setlist[itemize]{leftmargin=1.7em,itemsep=4pt,topsep=4pt,label={\color{poporange}\textbullet}}
\setlength{\parindent}{0pt}
\setlength{\parskip}{5pt}
\renewcommand{\arraystretch}{1.25}

% pgfplots : style Pop par défaut
\pgfplotsset{
  every axis/.append style={axis line style={popink,line width=1pt},tick style={popink},
    grid style={popline,line width=0.5pt},label style={font=\small},tick label style={font=\footnotesize}},
  popcurve/.style={poporange,line width=1.8pt,no markers,smooth},
}
% Même style disponible dans un \draw TikZ simple (hors pgfplots)
\tikzset{popcurve/.style={poporange,line width=1.8pt}}
\tikzset{popgrid/.style={popline,very thin}}
\colorlet{popcurve}{poporange} % le nom du style est parfois utilisé comme couleur

% ============================================================
% Pill / squiggle
% ============================================================
\newcommand{\poppill}[3]{%
  \tikz[baseline=-0.55ex]\node[draw=popink,line width=1.2pt,rounded corners=2.6pt,%
    fill=#2,inner xsep=6.5pt,inner ysep=3pt]{%
    \popxboldls\fontsize{7}{7}\selectfont\color{#3}\MakeUppercase{#1}};%
}
\newcommand{\popsquiggle}[1]{%
  \tikz[baseline=-0.4ex]{
    \draw[#1,line width=2.6pt,line cap=round]
      plot [smooth] coordinates {(0,0.09) (0.34,0.01) (0.68,0.21) (1.02,0.01) (1.36,0.10)};
  }%
}
% Mot-clé surligné (marqueur orange sous le texte)
\newcommand{\cle}[1]{{\popxbold\color{popdark}#1}}

% ============================================================
% En-tête / pied
% ============================================================
\pagestyle{fancy}
\fancyhf{}
\renewcommand{\headrulewidth}{0pt}
\renewcommand{\footrulewidth}{0pt}
\lhead{\raisebox{-0.15ex}{\poplogo\fontsize{13}{13}\selectfont\color{popink}OnBuch\color{poporange}+}}
\rhead{\poppill{\DOCMATIERE\ \textbullet\ \DOCNIVEAU\ \textbullet\ Leçon \DOCLECON}{white}{popink}}
\lfoot{\popxbold\fontsize{7.6}{7.6}\selectfont\color{popmuted}\MakeUppercase{Cours signé OnBuch+}\ \textbullet\ {\popsemi\DOCTITRE}}
\rfoot{\tikz[baseline=-0.6ex]\node[rounded corners=8.5pt,fill=popink,minimum height=17pt,minimum width=17pt,inner xsep=5pt,inner sep=0]{\popxbold\fontsize{8}{8}\selectfont\color{popcream}\thepage\,/\,\pageref*{LastPage}};}

% ============================================================
% Titres
% ============================================================
\newcommand{\chaptitre}[2]{%
  \begin{tcolorbox}[enhanced,colback=poporange,colframe=popink,boxrule=2.6pt,arc=11pt,
      drop shadow=popink,left=15pt,right=15pt,top=12pt,bottom=11pt,before skip=3pt,after skip=18pt]
    \poppill{#2}{popink}{popcream}\par\vspace{9pt}
    {\raggedright\hyphenpenalty=10000\exhyphenpenalty=10000\popdisplay\fontsize{22}{24}\selectfont\color{popcream}\MakeUppercase{#1}\par}
  \end{tcolorbox}
}
% Section numérotée : pastille orange avec le numéro + titre Archivo
\newcounter{popsec}
\newcommand{\coursec}[1]{%
  \stepcounter{popsec}\setcounter{popsub}{0}%
  \par\vspace{16pt}\noindent
  \begin{minipage}{\linewidth}
  \tikz[baseline=-0.7ex]\node[draw=popink,line width=1.6pt,fill=poporange,rounded corners=4pt,minimum size=22pt,inner sep=0]{\popdisplay\fontsize{12}{12}\selectfont\color{popcream}\thepopsec};%
  \hspace{8pt}{\popdisplay\fontsize{15}{17}\selectfont\color{popink}\MakeUppercase{#1}}\par
  \vspace{4pt}\noindent{\color{poporange}\rule{36pt}{3.6pt}}
  \end{minipage}\par\nopagebreak\vspace{8pt}
}
\newcounter{popsub}
\newcommand{\courssub}[1]{%
  \stepcounter{popsub}%
  \par\vspace{9pt}\noindent{\popxbold\fontsize{11.5}{13}\selectfont\color{popink}{\color{poporange}\thepopsec.\thepopsub}\hspace{6pt}#1}\par\nopagebreak\vspace{4pt}
}

% ============================================================
% Boîtes pédagogiques — même recette : fond blanc, bordure épaisse colorée,
% ombre dure encre, badge plein. Titre optionnel [#1].
% ============================================================
\tcbset{popbase/.style={breakable,enhanced,colback=white,boxrule=2.3pt,arc=9pt,
  shadow={3pt}{-3pt}{0pt}{popink},left=14pt,right=14pt,top=11pt,bottom=11pt,before skip=11pt,after skip=15pt,
  fontupper=\popsemi\fontsize{10.6}{14.5}\selectfont\color{popink},
  fontlower=\popsemi\fontsize{10.6}{14.5}\selectfont\color{popink}}}
% Coche dessinée (le glyphe ✓ n'existe pas dans les polices du document)
\newcommand{\popcheck}[1]{\tikz[baseline=-0.5ex]\draw[#1,line width=1.6pt,line cap=round,line join=round](-0.13,0.0)--(-0.04,-0.09)--(0.13,0.1);}
\providecommand{\coche}{\popcheck{popgreen}}
\providecommand{\resultat}[1]{\par\smallskip\noindent\fcolorbox{popgreen}{popgreenL}{\parbox{\dimexpr\linewidth-2\fboxsep-2\fboxrule\relax}{\textbf{Résultat.} #1}}\par\smallskip}
\newcommand{\popboxhead}[3]{\poppill{#1}{#2}{white}\ifx\relax#3\relax\else\hspace{6pt}{\popxbold\fontsize{10.6}{12}\selectfont\color{popink}#3}\fi\par\vspace{7pt}}

\newtcolorbox{aretenir}[1][]{popbase,colframe=popgreen,before upper={\popboxhead{À retenir}{popgreen}{#1}}}
\newtcolorbox{exemplebox}[1][]{popbase,colframe=poporange,before upper={\popboxhead{Exemple}{poporange}{#1}}}
\newtcolorbox{definition}[1][]{popbase,colframe=popblue,before upper={\popboxhead{Définition}{popblue}{#1}}}
\newtcolorbox{propriete}[1][]{popbase,colframe=poppurple,before upper={\popboxhead{Propriété}{poppurple}{#1}}}
\newtcolorbox{methode}[1][]{popbase,colframe=popgold,before upper={\popboxhead{Méthode}{popgold}{#1}}}
\newtcolorbox{attention}[1][]{popbase,colframe=poppink,before upper={\popboxhead{Attention}{poppink}{#1}}}
\newtcolorbox{experience}[1][]{popbase,colframe=popblue,colback=popblueL,before upper={\popboxhead{Expérience}{popblue}{#1}}}
% « Le savais-tu ? » : carte violette pleine (contexte, culture, Cameroun)
\newtcolorbox{savaistu}[1][]{popbase,colback=poppurple,colframe=popink,
  fontupper=\popsemi\fontsize{10.4}{14.2}\selectfont\color{white},
  before upper={\poppill{Le savais-tu\,?}{white}{poppurple}\ifx\relax#1\relax\else\hspace{6pt}{\popxbold\color{white}#1}\fi\par\vspace{7pt}}}
% Exercice résolu : énoncé en haut, \tcblower puis la solution sur fond crème
\newcounter{popexo}\newcounter{popexr}
\newtcolorbox{exoresolu}[1][]{popbase,colframe=poporange,colbacklower=popcream,
  segmentation style={popink,line width=1.2pt,dashed},
  before upper={\stepcounter{popexr}\popboxhead{Exercice résolu \thepopexr}{poporange}{#1}},
  before lower={\poppill{Solution}{popink}{popcream}\par\vspace{6pt}}}
% Exercice (non résolu en place) — la correction est dans la partie Corrigés
\newtcolorbox{exercice}[1][]{popbase,colframe=popink,
  before upper={\stepcounter{popexo}\popboxhead{Exercice \thepopexo}{popink}{#1}}}
% Corrigé
\newtcolorbox{corrige}[1][]{popbase,colframe=popgreen,colback=popgreenL,
  before upper={\popboxhead{Corrigé}{popgreen}{#1}}}
% Objectifs / prérequis / bilan
\newtcolorbox{objectifs}{popbase,colframe=popink,colback=poporangeL,
  before upper={\popboxhead{Objectifs de la leçon}{popink}{}}}
\newtcolorbox{prerequis}{popbase,colframe=popink,colback=popblueL,
  before upper={\popboxhead{Prérequis}{popblue}{}}}
\newtcolorbox{bilan}{popbase,colframe=popink,colback=white,boxrule=2.8pt,
  before upper={\popboxhead{Fiche bilan}{poporange}{L'essentiel à connaître pour l'examen}}}
% Figure encadrée avec légende
\newtcolorbox{popfigure}[1][]{enhanced,breakable=false,colback=white,colframe=popline,boxrule=1.4pt,arc=8pt,
  left=10pt,right=10pt,top=10pt,bottom=8pt,before skip=10pt,after skip=12pt,halign=center,
  lower separated=false,halign lower=center,
  fontlower=\popsemi\fontsize{9}{11.5}\selectfont\color{popmuted},
  #1}
\newcommand{\legende}[1]{\par\vspace{6pt}{\popsemi\fontsize{9}{11.5}\selectfont\color{popmuted}#1}}

% ============================================================
% Signature OnBuch+
% ============================================================
\newcommand{\poplogoinline}[1]{{\poplogo\fontsize{#1}{#1}\selectfont\color{popink}OnBuch\color{poporange}+}}
\newcommand{\signatureonbuch}{%
  \begin{tcolorbox}[enhanced,colback=popink,colframe=popink,boxrule=2.2pt,arc=10pt,
      drop shadow=poporange,left=14pt,right=14pt,top=10pt,bottom=10pt,before skip=0pt,after skip=0pt]
    \tikz[baseline=-0.7ex]\node[circle,fill=popgreen,draw=popcream,line width=1.3pt,minimum size=22pt,inner sep=0]{\popcheck{white}};%
    \hspace{9pt}\begin{minipage}[c]{0.62\linewidth}
      {\popxbold\fontsize{12}{13}\selectfont\color{popcream}Signé\ \textbullet\ {\poplogo OnBuch\color{poporange}+}}\par\vspace{2pt}
      {\raggedright\popsemi\fontsize{8}{10.5}\selectfont\color{popline}\MakeUppercase{Équipe pédagogique OnBuch+}\\\MakeUppercase{Conforme au programme officiel MINESEC}\par}
    \end{minipage}\hfill
    {\poplogo\fontsize{22}{22}\selectfont\color{popcream}OnBuch\color{poporange}+}
  \end{tcolorbox}%
}

% ============================================================
% Page de garde
% #1 titre · #2 matière · #3 niveau · #4 module · #5 n° leçon · #6 durée ·
% #7 description / objectifs (texte ou itemize)
% ============================================================
\newcommand{\pagegarde}[7]{%
  \thispagestyle{empty}
  \newgeometry{a4paper,margin=1.7cm,top=1.6cm,bottom=1.4cm}%
  \begin{tikzpicture}[remember picture,overlay]
    \fill[poporange!18] ([xshift=-3.2cm,yshift=-3.6cm]current page.north east) circle (4.6cm);
    \fill[poppurple!14] ([xshift=2.4cm,yshift=7.6cm]current page.south west) circle (3.8cm);
  \end{tikzpicture}%
  {\poplogo\fontsize{30}{30}\selectfont\color{popink}OnBuch\color{poporange}+}\hfill
  \raisebox{0.6ex}{\poppill{Cours signé \textbullet\ Premium}{popink}{popcream}}\par
  \vspace{1pt}\popsquiggle{poppurple}\par
  \vspace{8pt}
  \begin{tcolorbox}[enhanced,colback=poporange,colframe=popink,boxrule=2.8pt,arc=13pt,
      drop shadow=popink,left=18pt,right=18pt,top=12pt,bottom=13pt,after skip=10pt]
    \poppill{#2 \textbullet\ #3}{popink}{popcream}\hspace{5pt}\poppill{#4}{white}{popink}\par\vspace{8pt}
    {\popdisplay\fontsize{14}{14}\selectfont\color{popink}LEÇON #5}\par\vspace{4pt}
    {\raggedright\hyphenpenalty=10000\exhyphenpenalty=10000\popdisplay\fontsize{25}{27}\selectfont\color{popcream}\MakeUppercase{#1}\par}
    \vspace{8pt}
    {\popxbold\fontsize{9.5}{9.5}\selectfont\color{popink}\MakeUppercase{Durée conseillée : #6}}
  \end{tcolorbox}
  \begin{tcolorbox}[enhanced,colback=white,colframe=popink,boxrule=2.2pt,arc=10pt,
      drop shadow=popink,left=16pt,right=16pt,top=10pt,bottom=10pt,after skip=8pt]
    \poppill{Dans cette leçon}{poppurple}{white}\par\vspace{5pt}
    \popsemi\fontsize{9.3}{12.6}\selectfont\color{popink}#7
  \end{tcolorbox}
  \noindent\begin{minipage}[t]{0.487\linewidth}
    \begin{tcolorbox}[enhanced,colback=popgreen,colframe=popink,boxrule=2.2pt,arc=10pt,
        drop shadow=popink,left=11pt,right=11pt,top=10pt,bottom=10pt,height=3.1cm,valign=center]
      \tcbox[colback=white,colframe=popink,boxrule=1.3pt,arc=5pt,boxsep=0pt,nobeforeafter,
        left=3pt,right=3pt,top=3pt,bottom=3pt,valign=center]{\qrcode[height=1.9cm]{https://play.google.com/store/apps/details?id=com.luvvix.onbuch&hl=fr}}%
      \hspace{8pt}\begin{minipage}[c]{0.5\linewidth}
        {\popdisplay\fontsize{11}{12.5}\selectfont\color{white}\MakeUppercase{Télécharge OnBuch}}\par\vspace{4pt}
        {\raggedright\popsemi\fontsize{8.4}{11}\selectfont\color{white}Gratuit \textbullet\ Google Play\\Tuteur IA, annales, concours, résultats\par}
      \end{minipage}
    \end{tcolorbox}
  \end{minipage}\hfill
  \begin{minipage}[t]{0.487\linewidth}
    \begin{tcolorbox}[enhanced,colback=poppurple,colframe=popink,boxrule=2.2pt,arc=10pt,
        drop shadow=popink,left=11pt,right=11pt,top=10pt,bottom=10pt,height=3.1cm,valign=center]
      \tikz[baseline=-0.7ex]\node[circle,fill=white,draw=popink,line width=1.3pt,minimum size=1.6cm,inner sep=0]{\popdisplay\fontsize{24}{24}\selectfont\color{poppurple}+};%
      \hspace{8pt}\begin{minipage}[c]{0.6\linewidth}
        {\popdisplay\fontsize{11}{12.5}\selectfont\color{white}\MakeUppercase{Passe à OnBuch+}}\par\vspace{4pt}
        {\raggedright\popsemi\fontsize{8.4}{11}\selectfont\color{white}Tous les cours signés, Tuteur IA illimité, fiches PDF — onglet OnBuch+ de l'app\par}
      \end{minipage}
    \end{tcolorbox}
  \end{minipage}
  \vspace{8pt}
  \popcontenu
  \vfill
  \signatureonbuch
  \restoregeometry
  \newpage
}

% Rangée « Ce cours contient » (page de garde)
\newcommand{\poptile}[3]{% #1 couleur #2 gros texte #3 légende
  \begin{tcolorbox}[enhanced,colback=#1,colframe=popink,boxrule=2pt,arc=9pt,shadow={2.5pt}{-2.5pt}{0pt}{popink},
      left=6pt,right=6pt,top=9pt,bottom=9pt,halign=center,valign=center,height=1.75cm,width=\linewidth,nobeforeafter]
    {\popdisplay\fontsize{12.5}{13}\selectfont\color{white}\MakeUppercase{#2}}\par\vspace{3pt}
    {\centering\popsemi\fontsize{7.8}{9.5}\selectfont\color{white}#3\par}
  \end{tcolorbox}}
\newcommand{\popcontenu}{%
  \par\noindent{\popxboldls\fontsize{7.5}{7.5}\selectfont\color{popmuted}CE COURS SIGNÉ CONTIENT}\par\vspace{7pt}
  \noindent\begin{minipage}[t]{0.235\linewidth}\poptile{popblue}{Cours}{complet, illustré, pas à pas}\end{minipage}\hfill
  \begin{minipage}[t]{0.235\linewidth}\poptile{poporange}{Méthodes}{et exercices résolus}\end{minipage}\hfill
  \begin{minipage}[t]{0.235\linewidth}\poptile{poppink}{Exercices}{type examen + corrigés}\end{minipage}\hfill
  \begin{minipage}[t]{0.235\linewidth}\poptile{popgreen}{Fiche bilan}{l'essentiel à retenir}\end{minipage}\par}

% Sommaire
\newtcolorbox{sommaire}{enhanced,colback=white,colframe=popink,boxrule=2.2pt,arc=10pt,
  drop shadow=popink,left=18pt,right=18pt,top=13pt,bottom=13pt,before skip=6pt,after skip=20pt,
  before upper={\poppill{Sommaire}{poppurple}{white}\par\vspace{9pt}\popsemi\fontsize{10.6}{14.5}\selectfont\color{popink}}}

% Signature de fin de cours — poussée tout en bas de la dernière page
\newcommand{\findecours}{%
  \par\vfill\nopagebreak
  \begin{center}\popsquiggle{poporange}\end{center}\vspace{4pt}
  \signatureonbuch
}
\AtBeginDocument{\def\llbracket{\mathopen{[\mkern-3.5mu[}}\def\rrbracket{\mathclose{]\mkern-3.5mu]}}} % intervalles d'entiers (glyphe ⟦ absent de la police)
% Glyphes absents de Fira Math : redéfinis à partir de glyphes disponibles
\AtBeginDocument{%
  \def\setminus{\mathbin{\backslash}}\let\smallsetminus\setminus
  \def\vdots{\mathord{\vbox{\baselineskip3.2pt\lineskiplimit0pt\kern4pt\hbox{.}\hbox{.}\hbox{.}}}}%
  \def\ddots{\mathinner{\mkern1mu\raise6.5pt\hbox{.}\mkern2mu\raise3.6pt\hbox{.}\mkern2mu\raise0.7pt\hbox{.}\mkern1mu}}%
  \def\triangle{\mathord{\scalebox{0.9}{$\symup{\Delta}$}}}%
  \def\nabla{\mathord{\raisebox{\depth}{\scalebox{1}[-1]{$\symup{\Delta}$}}}}%
  \def\checkmark{\popcheck{popdark}}%
  \def\star{\mathbin{\ast}}\def\bigstar{\mathord{\ast}}%
  \def\diamond{\mathbin{\scalebox{0.6}{$\Diamond$}}}%
  \def\bigcup{\mathop{\mathchoice{\raisebox{-0.3em}{\scalebox{1.7}{$\cup$}}}{\scalebox{1.25}{$\cup$}}{\cup}{\cup}}\displaylimits}%
  \def\bigcap{\mathop{\mathchoice{\raisebox{-0.3em}{\scalebox{1.7}{$\cap$}}}{\scalebox{1.25}{$\cap$}}{\cap}{\cap}}\displaylimits}%
  \def\lesseqqgtr{\mathrel{\lessgtr}}\def\bigoplus{\mathop{\oplus}\displaylimits}%
}
\providecommand{\ssi}{si et seulement si} % abréviation fréquente
\AtBeginDocument{\providecommand{\wideparen}{\overparen}} % arc de cercle

%% FIGFIX-BEGIN v4 — correctifs globaux des figures (docs/onbuch-generation/tools/figfix)
\usepackage{adjustbox}\usepackage{etoolbox}
% toute figure TikZ/pgfplots plus large que la ligne est réduite à la largeur disponible
\BeforeBeginEnvironment{tikzpicture}{\begin{adjustbox}{max width=\linewidth}}
\AfterEndEnvironment{tikzpicture}{\end{adjustbox}}
% légendes pgfplots lisibles par-dessus les courbes
\pgfplotsset{every axis/.append style={legend style={fill=white,fill opacity=0.92,text opacity=1,draw=black!15,inner sep=3pt}}}
% étiquettes d'arêtes (syntaxe edge["..."]) avec fond blanc
\tikzset{every edge quotes/.append style={fill=white,inner sep=1.2pt}}
%% FIGFIX-END
\begin{document}

\pagegarde{(demande de stage)}{Français}{1re Tech}{Français – classe de Première technique}{N10}{14 séances (fiche de progression harmonisée)}{Cette leçon conduit l'élève de Première technique à maîtriser la rédaction complète d'une lettre officielle de demande de stage : structure externe et interne, introduction, corps et conclusion. Elle s'appuie sur l'analyse de textes fonctionnels authentiques et sur les notions de cohésion et de cohérence du texte, avec des prolongements en stylistique (texte explicatif) et en communication par l'image.
\vspace{4pt}
\begin{multicols}{2}\small
\begin{itemize}
\item À la fin de cette leçon, je sais identifier les marques de cohésion et de cohérence dans un texte fonctionnel.
\item À la fin de cette leçon, je sais décrire la structure externe et interne d'une lettre officielle de demande de stage.
\item À la fin de cette leçon, je sais lire et interpréter un texte fonctionnel (lettre, annonce, affiche) en dégageant ses informations essentielles.
\item À la fin de cette leçon, je sais reconnaître les procédés explicatifs (définition, exemple, comparaison) dans un texte.
\item À la fin de cette leçon, je sais analyser une situation de communication (émetteur, destinataire, objet, contexte) avant de rédiger.
\item À la fin de cette leçon, je sais rédiger l'introduction, le corps et la conclusion d'une demande de stage.
\end{itemize}\end{multicols}}

\begin{sommaire}
\begin{multicols}{2}
\begin{enumerate}
\item Cohésion et cohérence du texte
\item La lettre officielle de demande de stage : structure externe et interne
\item Lecture de textes fonctionnels et texte explicatif
\item La communication par l'image
\item Rédaction de la demande de stage : introduction, corps, conclusion
\item Production, confrontation, amélioration et remédiation
\item Activité d'intégration
\item Méthodes et exercices résolus
\item Exercices
\item Corrigés des exercices
\item Fiche bilan
\end{enumerate}
\end{multicols}
\end{sommaire}

\begin{objectifs}
\begin{itemize}
\item À la fin de cette leçon, je sais identifier les marques de cohésion et de cohérence dans un texte fonctionnel.
\item À la fin de cette leçon, je sais décrire la structure externe et interne d'une lettre officielle de demande de stage.
\item À la fin de cette leçon, je sais lire et interpréter un texte fonctionnel (lettre, annonce, affiche) en dégageant ses informations essentielles.
\item À la fin de cette leçon, je sais reconnaître les procédés explicatifs (définition, exemple, comparaison) dans un texte.
\item À la fin de cette leçon, je sais analyser une situation de communication (émetteur, destinataire, objet, contexte) avant de rédiger.
\item À la fin de cette leçon, je sais rédiger l'introduction, le corps et la conclusion d'une demande de stage.
\item À la fin de cette leçon, je sais interpréter le message d'une image (affiche, photographie, logo) dans une communication professionnelle.
\item À la fin de cette leçon, je sais confronter ma production à celle d'autrui, l'évaluer et l'améliorer.
\item À la fin de cette leçon, je sais rédiger et justifier une démarche, une réponse ou une conclusion dans une situation concrète de la vie courante.
\end{itemize}
\end{objectifs}

\begin{prerequis}
\begin{itemize}
\item Les types et formes de phrases (classe de Seconde)
\item Les éléments de la correspondance administrative vus en Seconde (en-tête, formules de politesse)
\item Les connecteurs logiques simples (d'abord, ensuite, enfin, parce que)
\item La notion de situation de communication (qui parle, à qui, de quoi, dans quel contexte)
\item La lecture méthodique d'un texte (repérage du thème et de la thèse)
\end{itemize}
\end{prerequis}

\newpage

\coursec{Cohésion et cohérence du texte}

Les grandes vacances approchent et le lycée technique de Nkongsamba demande à chaque élève de Première de trouver un stage en entreprise. Amina veut postuler à la SABC de Douala, mais son camarade Rodrigue lui affirme qu'une simple visite au bureau suffit. Le directeur des ressources humaines, lui, exige une lettre officielle écrite. Comment rédiger une demande de stage claire, cohérente et convaincante qui donnera envie à l'entreprise de l'accepter ? C'est tout l'enjeu de cette leçon. Pour y parvenir, nous devons d'abord maîtriser les deux piliers invisibles de tout texte qui « tient la route » : la \cle{cohésion} et la \cle{cohérence}.

\courssub{La cohésion textuelle : les liens grammaticaux et lexicaux}

\begin{definition}[Cohésion textuelle]
La \textbf{cohésion} est l'ensemble des procédés \emph{linguistiques} (grammaticaux et lexicaux) qui assurent la liaison explicite entre les énoncés successifs d'un texte. Elle rend le texte « tissé » : chaque phrase accroche la précédente et prépare la suivante par des marques repérables à la surface du discours.
\end{definition}

La cohésion ne relève pas du hasard ; elle repose sur une boîte à outils précise que tout rédacteur technique doit connaître et maîtriser.

\begin{methode}[Repérer les marques de cohésion dans un texte]
\textbf{Étape 1 : Repérer les reprises pronominales.} Soulignez tous les pronoms (personnels, relatifs, démonstratifs, possessifs) et fléchez-les vers leur antécédent (le nom qu'ils remplacent).\\
\textbf{Étape 2 : Identifier les substitutions lexicales.} Repérez les synonymes, hyperonymes (termes plus généraux), antonymes ou paraphrases qui reprennent une idée sans répéter le mot exact.\\
\textbf{Étape 3 : Délimiter les champs lexicaux.} Listez les mots appartenant à un même domaine sémantique (ex. : stage, entreprise, convention, tuteur, mission, rapport). Ils tissent une toile de fond thématique.\\
\textbf{Étape 4 : Isoler les connecteurs logiques et temporels.} Encadrez les mots de liaison (donc, cependant, ensuite, par conséquent, en effet, de plus) et précisez la relation qu'ils instaurent (cause, conséquence, opposition, chronologie, explication).
\end{methode}

\textbf{Les quatre piliers de la cohésion.}\par

\begin{itemize}
    \item \textbf{Les reprises pronominales (anaphore / cataphore).} Elles évitent la répétition lourde et créent une chaîne de référence.
    \begin{exemplebox}[Chaîne de référence pronominale]
    \emph{« Monsieur le Directeur, \textbf{votre} entreprise jouit d'une excellente réputation. \textbf{Elle} forme régulièrement des stagiaires. \textbf{Cette} politique nous intéresse. \textbf{Elle} prouve votre engagement. »}\\
    \textbf{Chaîne :} \cle{votre entreprise} $\rightarrow$ \cle{Elle} (pronom personnel) $\rightarrow$ \cle{Cette politique} (démonstratif + hyperonyme) $\rightarrow$ \cle{Elle} (reprise du sujet).
    \end{exemplebox}

    \item \textbf{Les substitutions lexicales (variation lexicale).} Synonymie, hyperonymie, métonymie, ellipse. Elles enrichissent le vocabulaire tout en maintenant la référence.
    \begin{exemplebox}[Substitution]
    \emph{« Je sollicite un \textbf{stage} de fin d'études. Cette \textbf{période de formation} en milieu professionnel me permettra d'acquérir des compétences. \textbf{Cette immersion} est obligatoire. »}\\
    Stage = Période de formation = Immersion (chaîne synonymique / hyperonymique).
    \end{exemplebox}

    \item \textbf{Les champs lexicaux (isotopies).} La répétition de mots d'un même domaine sémantique assure la continuité thématique.
    \begin{exemplebox}[Champ lexical « Stage / Entreprise »]
    \emph{Convention, tuteur, mission, rapport, évaluation, encadrement, service, département, ressources humaines, stage, stage, stage.}
    \end{exemplebox}

    \item \textbf{Les connecteurs (opérateurs argumentatifs).} Ils explicite la logique du discours : \cle{donc} (conséquence), \cle{cependant} (concession/opposition), \cle{en outre} (addition), \cle{en effet} (justification), \cle{enfin} (marque de fin/achèvement).
\end{itemize}

\begin{popfigure}
\begin{tikzpicture}[
    node distance=1.2cm and 1.5cm,
    box/.style={rectangle, draw=popblue, thick, fill=popblue!10, rounded corners, minimum width=2.2cm, minimum height=1cm, align=center, font=\small},
    arrow/.style={-Stealth, thick, popdark},
    label/.style={font=\tiny, text=popmuted, above, sloped}
]
% Nodes
\node[box, align=center] (nom){\textbf{Nom}\\\emph{Monsieur le Directeur}};
\node[box, right=of nom, align=center] (pro){\textbf{Pronom}\\\emph{Vous / Votre}};
\node[box, right=of pro, align=center] (syn){\textbf{Synonyme / Hyperonyme}\\\emph{Votre structure / L'établissement}};
\node[box, below right=of syn, align=center] (conn){\textbf{Connecteur}\\\emph{Par conséquent}};
\node[box, below=of conn, align=center] (nom2){\textbf{Nom (nouveau thème)}\\\emph{Ma candidature}};
% Arrows
\draw[arrow] (nom) -- node[label] {renvoie à} (pro);
\draw[arrow] (pro) -- node[label] {varie vers} (syn);
\draw[arrow] (syn) -- node[label] {logique :} (conn);
\draw[arrow] (conn) -- node[label] {introduit} (nom2);
% Global label
\node[above=0.3cm of nom, font=\bfseries\small, text=popblue] {Chaîne de cohésion dans l'incipit d'une lettre};
\end{tikzpicture}
\legende{Schéma de la chaîne de reprises (nom $\rightarrow$ pronom $\rightarrow$ synonyme) articulée par un connecteur logique. Chaque flèche représente un lien de cohésion explicite.}
\end{popfigure}

\begin{exoresolu}[Analyse de cohésion : extrait de lettre officielle camerounaise]
\textbf{Énoncé :} Repérez et classez les marques de cohésion dans l'extrait suivant :\\
\emph{« Objet : Demande de stage de fin d'études.\\
Monsieur le Directeur Général,\\
Actuellement élève en classe de Première Technique au Lycée de Nkongsamba, \textbf{je} dois effectuer un stage de six semaines. \textbf{Cette} immersion en milieu professionnel est prévue du 15 juillet au 30 août. \textbf{Elle} me permettra de confronter mes acquis théoriques à la réalité du terrain. \textbf{Par ailleurs}, votre entreprise, leader dans le secteur agro-industriel, constitue un cadre idéal. \textbf{C'est pourquoi} \textbf{je} sollicite votre bienveillance. »}

\textbf{Solution détaillée :}
\begin{center}
\begin{tabularx}{\linewidth}{|l|X|X|}\hline
\textbf{Type} & \textbf{Marque repérée} & \textbf{Fonction / Antécédent} \\ \hline
Reprise pronominale & \emph{je} (x2) & Locuteur (élève), sujet constant. \\ \hline
Démonstratif + Nom & \emph{Cette immersion} & Reprend \emph{un stage} (hyperonyme + précision). \\ \hline
Pronom personnel & \emph{Elle} & Reprend \emph{Cette immersion} (sujet). \\ \hline
Connecteur (addition) & \emph{Par ailleurs} & Passe de la description du stage à la description de l'entreprise. \\ \hline
Substitution (hyperonyme) & \emph{votre entreprise} $\rightarrow$ \emph{un cadre idéal} & Évaluation positive de l'antécédent. \\ \hline
Connecteur (conséquence) & \emph{C'est pourquoi} & Lie l'argument (cadre idéal) à la demande (sollicite). \\ \hline
Champ lexical & stage, immersion, milieu professionnel, terrain, entreprise, secteur, cadre, stage & Isotopie « Formation / Entreprise ». \\ \hline
\end{tabularx}
\end{center}
\textbf{Conclusion :} Le texte est fortement cohésif : les chaînes de référence sont continues, les connecteurs balisent l'argumentation, le champ lexical ancre le propos.
\end{exoresolu}

\courssub{La cohérence textuelle : l'unité de sens et la logique des idées}

\begin{definition}[Cohérence textuelle]
La \textbf{cohérence} est la propriété sémantique et pragmatique qui fait qu'un texte forme un \emph{tout uni de sens}. Elle repose sur : (1) l'unité thématique (un seul sujet principal), (2) la progression logique des idées (thème $\rightarrow$ thèse $\rightarrow$ arguments $\rightarrow$ conclusion), (3) la conformité aux attentes du destinataire et aux conventions du genre textual (ici : la lettre administrative).
\end{definition}

La cohérence, c'est le « fil conducteur » invisible que le lecteur suit sans effort. Un texte cohérent répond à la question : \emph{« Où veut-on en venir et pourquoi cet ordre d'idées ? »}

\begin{attention}[Piège majeur : Confondre cohésion (forme) et cohérence (sens)]
C'est l'erreur la plus fréquente en examen.
\begin{itemize}
    \item \textbf{Cohésion} = \emph{Forme de surface}. On la voit : pronoms, connecteurs, champs lexicaux. Un texte peut être \textbf{très cohésif} (plein de « donc », « cependant », « il », « cette » bien employés) mais \textbf{totalement incohérent} (idées contradictoires, ordre absurde, hors-sujet).
    \item \textbf{Cohérence} = \emph{Profondeur du sens}. Elle se juge à la lecture globale. Un texte peut être \textbf{peu cohésif} (peu de connecteurs, répétitions de noms) mais \textbf{parfaitement cohérent} (idée claire, progression naturelle, but atteint).
\end{itemize}
\textbf{Au Probatoire :} On vous demandera souvent de \emph{justifier} la cohérence (organisation des idées) et non seulement de \emph{lister} des connecteurs (cohésion).
\end{attention}

\begin{center}
\begin{tabularx}{\linewidth}{|>{\bfseries}l|X|X|}\hline
\rowcolor{popblueL}
\textbf{Critère} & \textbf{Cohésion} & \textbf{Cohérence} \\ \hline
\textbf{Nature} & Linguistique (forme) & Sémantique / Pragmatique (sens) \\ \hline
\textbf{Question clé} & \emph{Comment les phrases s'accrochent-elles ?} & \emph{Le texte a-t-il un sens global logique ?} \\ \hline
\textbf{Outils} & Pronoms, connecteurs, champs lexicaux, substitutions & Thème unique, plan logique, genre textual, intention communicative \\ \hline
\textbf{Visibilité} & Repérable au niveau de la phrase / interphrastique & Repérable à la lecture globale (macrostructure) \\ \hline
\textbf{Exemple de défaillance} & « Je veux un stage. Il est difficile. Donc je pars. » (Qui ? Quoi ? Liens flous) & « Je veux un stage à la SABC. Les éléphants sont gris. Donc j'aime le chocolat. » (Phrases liées mais sens absent) \\ \hline
\textbf{Exemple de réussite} & « Je sollicite un stage. \textbf{Ce} stage durera 6 semaines. \textbf{Il} débutera en juillet. » & « Je suis élève en 1ère Tech. \textbf{Le programme exige} un stage. \textbf{Je postule} chez vous \textbf{car} vous êtes leader. \textbf{J'espère} une réponse favorable. » \\ \hline
\end{tabularx}
\end{center}

\courssub{Distinction et articulation : un texte cohésif mais incohérent ?}

Oui, c'est possible. La cohésion est \emph{nécessaire} mais non \emph{suffisante} pour la cohérence.

\begin{exemplebox}[Paragraphe cohésif mais incohérent — Analyse et correction pas à pas]
\textbf{Texte original (incohérent) :}
\emph{« Monsieur le Directeur, \textbf{je} vous écris pour un stage. \textbf{Votre} entreprise est loin. \textbf{Cependant}, \textbf{je} l'aime bien. \textbf{Donc}, \textbf{je} ne viendrai pas. \textbf{En effet}, mes notes sont bonnes. \textbf{C'est pourquoi}, acceptez ma candidature. »}

\textbf{Analyse de la cohésion (apparence correcte) :}
\begin{itemize}
    \item Reprises : \emph{je} (x4), \emph{Votre entreprise} $\rightarrow$ \emph{l'} (l'aime), \emph{je} $\rightarrow$ \emph{ma} (candidature).
    \item Connecteurs : \emph{Cependant} (opposition), \emph{Donc} (conséquence), \emph{En effet} (justification), \emph{C'est pourquoi} (conclusion).
    \item Champ lexical : stage, entreprise, notes, candidature.
\end{itemize}
$\rightarrow$ \textbf{Le texte est cohésif.}

\textbf{Analyse de la cohérence (échec total) :}
\begin{itemize}
    \item \textbf{Contradiction logique :} « Je l'aime bien \textbf{cependant} je ne viendrai pas » (l'opposition est absurde : l'éloignement justifierait le refus, pas l'appréciation).
    \item \textbf{Non-sens argumentatif :} « Mes notes sont bonnes \textbf{en effet} je ne viendrai pas » (aucune relation de cause à effet).
    \item \textbf{Incohérence pragmatique :} La conclusion « acceptez ma candidature » contredit l'annonce « je ne viendrai pas ».
    \item \textbf{Thème instable :} Éloignement $\rightarrow$ Appréciation $\rightarrow$ Refus $\rightarrow$ Notes $\rightarrow$ Acceptation.
\end{itemize}
$\rightarrow$ \textbf{Le texte est incohérent.}

\textbf{Correction (rétablissement de la cohérence) :}
\emph{« Monsieur le Directeur, \textbf{je} vous écris pour solliciter un stage. \textbf{Votre} entreprise est \textbf{éloignée} de mon domicile, \textbf{cependant} sa réputation d'excellence \textbf{me motive vivement}. \textbf{Par conséquent}, \textbf{je} m'engage à trouver un hébergement sur place. \textbf{En effet}, mes résultats scolaires sont bons et prouvent ma rigueur. \textbf{C'est pourquoi}, \textbf{je} vous prie d'accepter ma candidature. »}

\textbf{Verbalisation de la correction :}
\begin{enumerate}
    \item Remplacer « loin » (constat neutre/négatif) par « éloignée » + concession « cependant » + argument positif « réputation ».
    \item Remplacer « Donc je ne viendrai pas » (non-sens) par « Par conséquent, je m'engage à... » (solution constructive).
    \item Rétablir la justification « En effet » sur un argument pertinent (rigueur prouvée par notes).
    \item Aligner la conclusion sur l'objectif unique : obtenir le stage.
\end{enumerate}
\end{exemplebox}

\courssub{Application : repérage des marques de cohésion dans une lettre officielle camerounaise}

\begin{exoresolu}[Exercice de repérage guidé]
\textbf{Énoncé :} Voici le corps d'une demande de stage adressée au Directeur Général de la SABC. Identifiez \textbf{toutes} les marques de cohésion (grammaticales et lexicales) et expliquez leur rôle dans la construction de l'argumentation.

\emph{« Douala, le 10 juin 2025\\
Monsieur le Directeur Général,\\
Par la présente, \textbf{je}, soussigné \textbf{Ngo Bitjek Amina}, élève en classe de Première Technique (Série F4) au Lycée Technique de Nkongsamba, \textbf{vous} adresse ma candidature pour un stage de fin d'études.\\
\textbf{Cette} formation technique m'impose un stage de six semaines, du 14 juillet au 25 août 2025. \textbf{Il} s'agit pour moi de mettre en pratique les connaissances acquises en technologie brassicole et en maintenance industrielle. \textbf{En outre}, \textbf{votre} société, fleuron de l'industrie camerounaise, dispose d'un parc machine moderne que \textbf{je} souhaite découvrir.\\
\textbf{Par conséquent}, \textbf{je} \textbf{vous} saurais gré de bien vouloir accorder une suite favorable à \textbf{ma} demande. \textbf{Dans l'attente}, \textbf{je} \textbf{vous} prie d'agréer, Monsieur le Directeur Général, l'expression de \textbf{mes} salutations distinguées.\\
\textbf{Ngo Bitjek Amina} »}

\textbf{Solution détaillée :}
\begin{center}
\begin{tabularx}{\linewidth}{|l|X|X|}\hline
\rowcolor{popblueL}
\textbf{Marque} & \textbf{Type / Catégorie} & \textbf{Rôle cohésif / Argumentatif} \\ \hline
\emph{je} (x4) & Pronom personnel (1ère pers. sg) & Anaphore constante : ancre le locuteur (Amina) comme sujet agentif. \\ \hline
\emph{soussigné Ngo Bitjek Amina} & Nom propre + apposition & Identification formelle (exigence administrative). \\ \hline
\emph{vous} (x3) & Pronom personnel (2ème pers. pl. politesse) & Adressage au destinataire (DG) ; marque la relation institutionnelle. \\ \hline
\emph{Cette formation} & Démonstratif + hyperonyme & Reprend « Première Technique » ; ancre l'obligation dans le cursus. \\ \hline
\emph{Il s'agit pour moi de} & Présentatif + infinitif & Opérateur de reformulation / explicitation de l'objectif. \\ \hline
\emph{En outre} & Connecteur additif & Ajoute un argument positif (parc machine) à l'argument obligation. \\ \hline
\emph{votre société / fleuron / parc machine} & Chaîne nominale (hyperonyme $\rightarrow$ métonymie $\rightarrow$ métonymie) & Valorisation de l'entreprise (argument séductif / flatterie mesurée). \\ \hline
\emph{que je souhaite découvrir} & Relatif + verbe volitif & Lie le stagiaire à l'objet technique (projection active). \\ \hline
\emph{Par conséquent} & Connecteur consécutif & Point d'orgue logique : Obligation + Intérêt $\rightarrow$ Demande formelle. \\ \hline
\emph{ma demande / ma candidature} & Possessif + synonymie & Reprise anaphorique de l'acte de langage principal. \\ \hline
\emph{Dans l'attente} & Locution prépositive (ellipse de « Je reste dans l'attente ») & Marque de politesse finale / clôture protocolaire. \\ \hline
\emph{mes salutations} & Possessif + nom figé & Appropriation de la formule de politesse codifiée. \\ \hline
\end{tabularx}
\end{center}
\textbf{Synthèse :} La cohésion sert ici une \textbf{stratégie argumentative} en trois temps : 1. \textbf{Identité \& Obligation} (Qui ? Quoi ? Pourquoi ?), 2. \textbf{Intérêt \& Adéquation} (Pourquoi vous ?), 3. \textbf{Demande \& Politesse} (Action attendue). Les connecteurs (\emph{En outre}, \emph{Par conséquent}) sont les charnières logiques de cette progression.
\end{exoresolu}

\courssub{Le texte explicatif : réception et analyse}

\begin{definition}[Texte explicatif]
Le \textbf{texte explicatif} a pour fonction principale d'informer, de faire comprendre un phénomène, un processus, une règle ou une organisation. Dans le cadre professionnel (lettre de stage, rapport, notice), il répond à la question \emph{« Comment ? »} ou \emph{« Pourquoi ? »}. Il se caractérise par l'objectivité, la clarté, la précision terminologique et une structure logique (définition, description, analyse causes/effets, illustration).
\end{definition}

\begin{propriete}[Marqueurs du texte explicatif]
\begin{itemize}
    \item \textbf{Temps :} Présent de vérité générale (ex. : « Le stage \textbf{permet} d'acquérir... »), futur (procédure), passé composé (récit d'expérience).
    \item \textbf{Connecteurs :} \cle{car}, \cle{parce que}, \cle{donc}, \cle{par conséquent} (causalité) ; \cle{ainsi}, \cle{autrement dit} (reformulation) ; \cle{par exemple}, \cle{notamment} (illustration) ; \cle{premièrment}, \cle{ensuite}, \cle{enfin} (énumération/chronologie).
    \item \textbf{Structures :} Définitions (\emph{... est un ... qui ...}), classifications, comparaisons, schémas verbaux.
    \item \textbf{Vocabulaire :} Termes techniques précis (pas de familier), noms abstraits, verbes d'action ou d'état.
\end{itemize}
\end{propriete}

\begin{exemplebox}[Extrait explicatif dans une lettre de stage (justification du choix de l'entreprise)]
\emph{« Votre unité de production de la SABC Douala-Bassa intègre une chaîne d'embouteillage automatisée Krones de dernière génération. \textbf{Cette technologie} impose une maintenance prédictive rigoureuse basée sur l'analyse vibratoire. \textbf{En effet}, la surveillance continue des roulements \textbf{permet} d'anticiper les pannes et \textbf{garantit} la cadence de 60 000 bouteilles/heure. \textbf{Ainsi}, intégrer votre service maintenance m'offrirait une mise en situation réelle sur des équipements industriels de pointe. »}
\textbf{Analyse :} Précision technique (Krones, analyse vibratoire, 60 000 b/h) $\rightarrow$ Causalité (\emph{En effet}, \textbf{permet/garantit}) $\rightarrow$ Conclusion logique (\emph{Ainsi}) liant l'entreprise au projet du stagiaire.
\end{exemplebox}

\courssub{La lettre officielle de demande de stage : structure externe et interne}

\begin{definition}[Lettre officielle (administrative / professionnelle)]
La \textbf{lettre officielle} est un écrit codifié, à visée fonctionnelle (demander, informer, réclamer, postuler), adressé à une autorité ou une structure (entreprise, administration). Elle obéit à des normes strictes de présentation (structure externe) et d'organisation du contenu (structure interne) qui garantissent sa recevabilité et son efficacité.
\end{definition}

\begin{methode}[Structure externe : les 7 blocs obligatoires (norme administrative camerounaise)]
\textbf{1. En-tête expéditeur (haut gauche) :} Nom, Prénom, Adresse complète, Téléphone, Email, Classe/Série, Établissement.\\
\textbf{2. En-tête destinataire (haut droite, sous l'expéditeur) :} Qualité, Nom/Structure, Adresse (ex. : Monsieur le Directeur Général, SABC, BP 4070 Douala).\\
\textbf{3. Lieu et Date (ligne séparée, alignée à droite) :} Douala, le 10 juin 2025.\\
\textbf{4. Objet (gras, centré ou gauche, sans « Objet : » selon usages) :} \textbf{Demande de stage de fin d'études (Première Technique, Série F4)}.\\
\textbf{5. Formule d'appel (salutation) :} \emph{Monsieur le Directeur Général,} (suivi de retour à la ligne).\\
\textbf{6. Corps de la lettre (développé ci-après).}\\
\textbf{7. Formule de politesse (fermeture codifiée) + Signature manuscrite + Nom imprimé.}\\
\textbf{8. Mentions finales (bas de page) :} \emph{Pièces jointes :} CV, Bulletin de notes, Attestation de scolarité, Convention de stage (vierge).
\end{methode}

\begin{methode}[Structure interne du corps : la progression en 3 temps (standard Probatoire)]
\begin{enumerate}
    \item \textbf{Introduction (Le « Qui / Quoi / Pourquoi scolaire ») :} Identité précise $\rightarrow$ Contexte formation $\rightarrow$ Obligation curriculaire (durée, dates) $\rightarrow$ Objectif pédagogique (mettre en pratique, valider module).
    \item \textbf{Développement / Corps (Le « Pourquoi vous / Adéquation / Engagements ») :} Argument 1 : Notoriété / Spécialité de l'entreprise (flatterie mesurée + preuve technique). Argument 2 : Adéquation profil/poste (compétences acquises $\leftrightarrow$ besoins entreprise). Argument 3 : Engagement personnel (règlement, horaires, discrétion, assiduité).
    \item \textbf{Conclusion (L'Appel à l'action) :} Synthèse implicite $\rightarrow$ Demande formelle (\emph{sollicite votre bienveillance / prie d'agréer}) $\rightarrow$ Disponibilité pour entretien $\rightarrow$ Formule de politesse complète.
\end{enumerate}
\end{methode}

\courssub{La communication par l'image : un atout pour la candidature}

\begin{definition}[Communication par l'image]
Dans la recherche de stage, l'image (photo d'identité sur CV, logo de l'entreprise visée, infographie des compétences, mise en page de la lettre) n'est pas décorative : elle \textbf{communique}. Elle renforce la crédibilité (éthos), attire l'attention (pathos) et structure l'information (logos). Savoir la lire et la produire est une compétence technique.
\end{definition}

\begin{propriete}[Lecture d'un visuel professionnel (méthode O.P.T.I.C.)]
\begin{itemize}
    \item \textbf{O} - \textbf{O}bserver : Type d'image (photo, logo, graphique, organigramme, mise en page), couleurs, formes, texte intégré.
    \item \textbf{P} - \textbf{P}résenter le sujet : De quoi s'agit-il ? (Ex. : Logo SABC = étoile rouge + typographie bleue = dynamisme, ancrage national).
    \item \textbf{T} - \textbf{T}raiter l'information : Quel message ? Quelle cible ? Quel effet recherché ? (Ex. : Organigramme du service maintenance $\rightarrow$ montre la hiérarchie, identifie le tuteur potentiel).
    \item \textbf{I} - \textbf{I}nterpréter : Intention de l'émetteur (entreprise : sérieux, modernité ; candidat : rigueur, clarté).
    \item \textbf{C} - \textbf{C}onclure : Pertinence pour ma candidature (ex. : j'utilise les codes couleurs de l'entreprise dans mon CV pour montrer mon intégration culturelle).
\end{itemize}
\end{propriete}

\begin{exemplebox}[Analyse d'une mise en page de lettre (communication visuelle du texte)]
\begin{center}
\begin{tabularx}{\linewidth}{|l|X|X|}\hline
\rowcolor{popblueL}
\textbf{Élément visuel} & \textbf{Code / Norme} & \textbf{Effet sur le destinataire (DRH)} \\ \hline
Police \emph{Times New Roman} / \emph{Arial} 12 & Lisibilité, standard administratif & Professionnalisme, respect des codes. \\ \hline
Marges 2,5 cm / Interligne 1,15 & Aération, place pour annotations & Respect, facilité de lecture rapide. \\ \hline
Alignement justifié (corps) / Gauche (en-têtes) & Structure visuelle claire & Rigueur, organisation. \\ \hline
Gras sur \textbf{Objet}, \textbf{Dates}, \textbf{Durée} & Hiérarchisation de l'info clé & Repérage immédiat des infos critiques. \\ \hline
Espace blanc avant/après formules politesse & Respiration visuelle & Solennité, distinction. \\ \hline
\end{tabularx}
\end{center}
\textbf{Conseil :} Votre lettre \emph{est} une image avant d'être lue. Une mise en page soignée vaut « 50\% de l'admissibilité » aux dires de nombreux DRH camerounais.
\end{exemplebox}

\courssub{Rédaction de la demande de stage : introduction, corps, conclusion}

\begin{methode}[Rédaction pas à pas de l'introduction (Accroche + Contexte)]
\begin{enumerate}
    \item \textbf{Formule d'appel} : \emph{Monsieur le Directeur Général,} (ou \emph{Monsieur le Responsable des Ressources Humaines,}).
    \item \textbf{Phrase d'amorce (Objet rappelé) :} \emph{Par la présente, je, soussigné [Prénom NOM], élève en classe de Première Technique (Série [XX]) au Lycée Technique de [Ville], vous adresse ma candidature pour un stage de fin d'études.}
    \item \textbf{Ancrage curriculaire} : \emph{Dans le cadre de ma formation en [Spécialité, ex. Maintenance Industrielle], le programme prévoit un stage obligatoire de [X] semaines, du [Date début] au [Date fin].}
    \item \textbf{Objectif pédagogique} : \emph{Cette immersion vise à confronter mes acquis théoriques (notamment en [Module clé 1] et [Module clé 2]) à la réalité industrielle.}
\end{enumerate}
\end{methode}

\begin{methode}[Rédaction pas à pas du corps (Argumentation + Preuve + Engagement)]
\begin{enumerate}
    \item \textbf{Transition (Connecteur additif) :} \emph{En outre / Par ailleurs / De surcroît,}
    \item \textbf{Argument « Entreprise » (Flatterie argumentée) :} \emph{Votre entreprise, [référence / leader / fleuron] dans le secteur [Secteur précis], attire mon attention par [réalisations concrètes : parc machine, certification ISO, projet d'extension, réputation sociale].}
    \item \textbf{Preuve technique (Texte explicatif) :} \emph{Vos installations [détail technique] constituent un terrain d'apprentissage idéal pour approfondir mes compétences en [Compétence ciblée].}
    \item \textbf{Argument « Candidat » (Adéquation) :} \emph{Mon parcours scolaire (mention [Mention] au Probatoire blanc / notes en [Matières clés]) et mon projet professionnel ([Métier visé]) sont en parfaite adéquation avec vos activités.}
    \item \textbf{Engagements (Rassurer) :} \emph{Je m'engage à respecter scrupuleusement le règlement intérieur, les horaires, les consignes de sécurité et la confidentialité industrielle. Je ferai preuve d'assiduité, de discrétion et d'initiative.}
\end{enumerate}
\end{methode}

\begin{methode}[Rédaction pas à pas de la conclusion (Clôture + Politesse)]
\begin{enumerate}
    \item \textbf{Connecteur conclusif :} \emph{Par conséquent / C'est pourquoi / Au vu de ce qui précède,}
    \item \textbf{Demande formelle :} \emph{je vous saurais gré de bien vouloir accorder une suite favorable à ma demande / je sollicite votre bienveillance pour intégrer votre service [Nom du service].}
    \item \textbf{Disponibilité :} \emph{Je me tiens à votre entière disposition pour tout entretien à votre convenance.}
    \item \textbf{Formule de politesse complète (invariable, sans point final) :} \\
    \emph{Dans l'attente d'une réponse favorable, je vous prie d'agréer, Monsieur le Directeur Général, l'expression de mes salutations distinguées.}
    \item \textbf{Signature + Nom imprimé + Pièces jointes.}
\end{enumerate}
\end{methode}

\courssub{Production, confrontation, amélioration et remédiation}

\begin{exercice}[Remise en ordre et justification — Cohérence macrostructurelle]
Les phrases suivantes constituent le paragraphe central (corps de lettre) d'une demande de stage. Elles ont été mélangées.\\
\textbf{1.} Réorganisez-les pour obtenir un texte \textbf{cohérent} (progression logique : Contexte $\rightarrow$ Contraintes $\rightarrow$ Choix de l'entreprise $\rightarrow$ Engagements $\rightarrow$ Demande).\\
\textbf{2.} Soulignez les \textbf{mots-outils de cohésion} qui justifient votre ordonnancement.\\
\textbf{3.} Ajoutez, si nécessaire, \textbf{un connecteur manquant} pour fluidifier la transition entre deux phrases.

\medskip
\textbf{Phrases mélangées :}
\begin{enumerate}
    \item[A.] Je suis prêt à respecter le règlement intérieur et les horaires de votre usine.
    \item[B.] Actuellement en Première Maintenance Industrielle au Lycée Technique de Douala, je dois effectuer un stage obligatoire de huit semaines.
    \item[C.] C'est pourquoi je sollicite votre bienveillance pour intégrer votre service maintenance.
    \item[D.] Votre entreprise, référence dans le secteur de l'aluminium, correspond parfaitement à ma spécialité.
    \item[E.] De plus, ce stage me permettra de valider mon module de soudure TIG.
    \item[F.] En effet, vos installations récentes offrent un terrain d'apprentissage idéal.
\end{enumerate}
\end{exercice}

\begin{corrige}[Exercice n°1 — Remise en ordre et justification]
\textbf{Ordre logique : B $\rightarrow$ E $\rightarrow$ D $\rightarrow$ F $\rightarrow$ A $\rightarrow$ C}

\textbf{Texte rétabli :}
\emph{« \textbf{B} Actuellement en Première Maintenance Industrielle au Lycée Technique de Douala, je dois effectuer un stage obligatoire de huit semaines. \textbf{E} \textbf{De plus}, ce stage me permettra de valider mon module de soudure TIG. \textbf{D} Votre entreprise, référence dans le secteur de l'aluminium, correspond parfaitement à ma spécialité. \textbf{F} \textbf{En effet}, vos installations récentes offrent un terrain d'apprentissage idéal. \textbf{A} Je suis prêt à respecter le règlement intérieur et les horaires de votre usine. \textbf{C} \textbf{C'est pourquoi} je sollicite votre bienveillance pour intégrer votre service maintenance. »}

\textbf{Justification de la cohérence (macrostructure) :}
\begin{enumerate}
    \item \textbf{B (Situation initiale / Contexte) :} Présentation du locuteur, formation, contrainte curriculaire (stage obligatoire). Thème posé.
    \item \textbf{E (Objectif pédagogique personnel) :} « De plus » ajoute un bénéfice propre (validation module). Progression : contrainte $\rightarrow$ opportunité.
    \item \textbf{D (Choix ciblé de l'entreprise) :} Bascule vers le destinataire. Argument de pertinence (spécialité / secteur).
    \item \textbf{F (Preuve / Justification du choix) :} « En effet » explicite D : installations = terrain idéal. Preuve concrète.
    \item \textbf{A (Engagement / Garantie) :} Passage du « je veux » au « je m'engage ». Rassure l'employeur (règlement, horaires).
    \item \textbf{C (Conclusion / Appel à l'action) :} « C'est pourquoi » synthétise B+E+D+F+A $\rightarrow$ Demande formelle.
\end{enumerate}

\textbf{Mots-outils de cohésion repérés (microstructure) :}
\begin{itemize}
    \item \emph{De plus} (B$\rightarrow$E) : Addition d'arguments subjectifs.
    \item \emph{Votre entreprise / ma spécialité} (E$\rightarrow$D) : Chaîne lexicale « formation / métier / secteur ».
    \item \emph{En effet} (D$\rightarrow$F) : Relation explication / preuve.
    \item \emph{Vos installations / terrain d'apprentissage} (F) : Champ lexical technique (maintenance, soudure, installations).
    \item \emph{Je suis prêt / respecter} (A) : Engagement modal (volonté + capacité).
    \item \emph{C'est pourquoi} (A$\rightarrow$C) : Conséquence finale, marque l'aboutissement de l'argumentation.
\end{itemize}

\textbf{Connecteur ajouté :} \emph{De plus} au début de E (absent dans l'énoncé mélangé, indispensable pour marquer l'addition d'argument personnel après la contrainte scolaire). Sans lui, la juxtaposition B/E parait brute.
\end{corrige}

\begin{exercice}[Rédaction complète : Mise en situation Probatoire]
\textbf{Sujet :} Vous êtes \emph{Kouamé Franck}, élève en \textbf{Première Technique, Série F3 (Électrotechnique)} au \textbf{Lycée Technique de Bafoussam}. Vous devez effectuer un stage de \textbf{6 semaines} (du \textbf{7 juillet au 18 août 2025}) à \textbf{ENEO Cameroon} (Siège : Bonanjo, Douala), service \textbf{Distribution Réseau HTA/BT}.\\
Rédigez la \textbf{lettre officielle de demande de stage} complète (structure externe + interne). Votre lettre doit faire preuve de \textbf{cohésion} (chaînes de référence, connecteurs, champ lexical) et de \textbf{cohérence} (progression logique en 3 temps, adéquation profil/poste). Intégrez un argument technique précis (ex. : poste de transformation, protection différentielle, logiciel de supervision SCADA).
\end{exercice}

\begin{corrige}[Exercice n°2 — Proposition de rédaction type (Modèle Premium OnBuch+)]
\begin{popfigure}
\begin{tikzpicture}[node distance=0.2cm]
\node[align=left, font=\small\ttfamily, text=popdark, inner sep=0.5cm, fill=popcream, draw=poporange, rounded corners, thick, text width=\linewidth-1cm] (letter) {
Kouamé Franck\\
BP 124 Bafoussam\\
Tél : 6XX XX XX XX\\
Email : franck.kouame@email.cm\\
Élève 1ère Tech F3 (Électrotechnique) — Lycée Technique de Bafoussam
\hfill Douala, le 15 mai 2025
Monsieur le Directeur Général\\
ENEO Cameroon\\
Service Distribution — Direction Réseau\\
BP 2360 Douala (Bonanjo)
\textbf{Objet : Demande de stage de fin d'études (Première Technique F3 — Électrotechnique)}
Monsieur le Directeur Général,
Par la présente, \textbf{je}, soussigné \textbf{Kouamé Franck}, élève en classe de Première Technique, Série F3 (Électrotechnique) au Lycée Technique de Bafoussam, \textbf{vous} adresse ma candidature pour un stage de fin d'études.
\textbf{Cette} formation technique m'impose un stage de six semaines, du 7 juillet au 18 août 2025. \textbf{Il} s'agit pour moi de mettre en pratique les connaissances acquises en machines tournantes, installations électriques BT/HTA et automatismes programmables. \textbf{En outre}, \textbf{votre} société, opérateur historique et leader de la distribution d'énergie au Cameroun, dispose d'un réseau HTA/BT étendu et de postes de transformation modernes que \textbf{je} souhaite découvrir.
\textbf{En effet}, la supervision du réseau via le système SCADA au niveau du Dispatching national, ainsi que la maintenance préventive des protections différentielles et magnéto-thermiques sur les postes source, constituent des domaines d'excellence où \textbf{je} pourrai confronter mes acquis théoriques (notamment le dimensionnement de câbles et le calcul de court-circuit) à la réalité du terrain. \textbf{De plus}, la politique RSE d'ENEO en matière de sécurité électrique (norme NF C 15-100) rejoint ma rigueur scolaire.
\textbf{Par conséquent}, \textbf{je} m'engage à respecter scrupuleusement le règlement intérieur, les consignes de sécurité (port des EPI, consignation) et les horaires de service. \textbf{Je} \textbf{vous} saurais gré de bien vouloir accorder une suite favorable à \textbf{ma} demande. \textbf{Dans l'attente}, \textbf{je} \textbf{vous} prie d'agréer, Monsieur le Directeur Général, l'expression de \textbf{mes} salutations distinguées.
Kouamé Franck
\textbf{Pièces jointes :} CV — Bulletins 1er et 2e trimestres — Attestation de scolarité — Convention de stage (vierge) — Copie CNI.
};
\end{tikzpicture}
\legende{Modèle de lettre complète respectant la structure externe (en-têtes, objet, date, PJ) et interne (Intro : B+E ; Corps : D+F+A ; Conclusion : C). Cohésion assurée par chaînes \emph{je/vous/cette/il}, connecteurs \emph{En outre/En effet/De plus/Par conséquent}, champ lexical Électrotechnique/ENEO.}
\end{popfigure}
\end{corrige}

\begin{aretenir}[Bilan : Cohésion et Cohérence, les deux jambes du texte]
\begin{itemize}
    \item \textbf{Cohésion} = \emph{Couture} visible (pronoms, connecteurs, champs lexicaux). On la \textbf{repère} (analyse linguistique).
    \item \textbf{Cohérence} = \emph{Architecture} invisible (thème, plan, intention, genre). On la \textbf{évalue} (lecture globale).
    \item \textbf{Lien :} La cohésion \textbf{sert} la cohérence. Les connecteurs ne sont pas de la décoration : ils \textbf{balisent} la logique (cause, conséquence, concession, chronologie).
    \item \textbf{Structure lettre officielle :} 7 blocs externes codifiés + 3 temps internes (Contexte scolaire $\rightarrow$ Adéquation entreprise/technique $\rightarrow$ Engagement/Demande).
    \item \textbf{Texte explicatif :} Objectivité, terminologie précise, connecteurs causaux/illustratifs, présent de vérité. Sert à justifier le choix de l'entreprise.
    \item \textbf{Communication par l'image :} Mise en page, logo, photo, organigramme = signaux de professionnalisme (O.P.T.I.C.).
    \item \textbf{Au Probatoire :} Vérifiez : (1) Chaînes de référence claires (Je / Mon stage / Ma candidature), (2) Connecteurs logiques adaptés à l'argumentation (En outre, Par conséquent, C'est pourquoi), (3) Progression en 3 temps, (4) Précision technique (spécialité), (5) Formule de politesse exacte, (6) Pièces jointes listées.
\end{itemize}
\end{aretenir}

\coursec{La lettre officielle de demande de stage : structure externe et interne}

Dans la section précédente, nous avons vu qu'un texte n'est pas une simple suite de phrases : il doit être \cle{cohérent} (les idées s'enchaînent logiquement) et \cle{cohésif} (des outils linguistiques relient les énoncés entre eux). Ces deux exigences trouvent une application immédiate dans un genre de texte que tout élève de Première technique devra maîtriser pour sa vie professionnelle : la \cle{lettre officielle}, et plus particulièrement la \cle{demande de stage}. En effet, une lettre de demande de stage mal organisée, sans objet, sans formule de politesse, sera écartée par le destinataire avant même d'être lue. La cohérence et la cohésion ne sont donc pas des notions abstraites : elles conditionnent l'efficacité pratique de votre écrit.

\courssub{Nature et fonction de la lettre officielle}

Commençons par une situation concrète. Vous êtes élève en Première F3 (génie civil) ou en Première G (techniques administratives) et votre établissement vous demande d'effectuer un stage d'observation en entreprise pendant les vacances. Comment entrer en contact avec le directeur d'une entreprise que vous ne connaissez pas ? Vous ne pouvez pas lui téléphoner familièrement ni lui envoyer un message oral. Il faut un écrit \cle{formel}, qui respecte des conventions précises : c'est la lettre officielle.

\begin{definition}[La lettre officielle]
La \cle{lettre officielle} est un écrit formel, rédigé dans un \cle{registre soutenu}, adressé à une autorité administrative (préfet, maire, proviseur, ministre) ou à une personne morale (entreprise, banque, ONG) en vue d'obtenir quelque chose : un stage, un emploi, une autorisation, une information, une réparation. Elle obéit à des règles fixes de présentation (structure externe) et d'organisation des idées (structure interne).
\end{definition}

La \cle{demande de stage} est donc une variété de lettre officielle dont la fonction est précise : solliciter auprès d'une entreprise ou d'une administration l'autorisation d'y effectuer un stage, c'est-à-dire une période de formation pratique. Elle engage l'image de l'élève et de son établissement : c'est pourquoi sa rédaction exige soin, correction et respect des conventions.

\begin{attention}[Une lettre officielle n'est pas une lettre amicale]
Piège fréquent : rédiger sa demande de stage comme une lettre à un camarade. Sont donc interdits : les abréviations (« bjr », « svp », « jtm »), les tutoiements, les formules familières (« Salut », « Gros bisous »), les fautes d'orthographe et le papier déchiré ou froissé. Le destinataire juge le candidat sur la forme autant que sur le fond.
\end{attention}

\courssub{La structure externe : le squelette visible de la lettre}

La \cle{structure externe} désigne l'ensemble des éléments matériels qui encadrent le contenu de la lettre et qui permettent de l'identifier, de la dater et de l'acheminer. Ces éléments sont au nombre de huit et occupent des places fixes sur la page.

\begin{propriete}[Les huit éléments de la structure externe]
\begin{enumerate}
\item \textbf{L'en-tête} : nom, prénom, adresse et coordonnées (téléphone, courriel) de l'\cle{expéditeur}, placés en haut à gauche.
\item \textbf{Le lieu et la date} de rédaction, placés en haut à droite (ex. : « Bafoussam, le 12 mai 2026 »).
\item \textbf{Les coordonnées du destinataire} : fonction et adresse, placées à droite, sous la date (ex. : « Monsieur le Directeur de l'agence Express Union, Bafoussam »).
\item \textbf{L'objet} de la lettre, placé à gauche, avant la formule d'appel (ex. : « Objet : demande de stage académique »).
\item \textbf{La formule d'appel} : « Monsieur le Directeur, », « Madame la Directrice, ».
\item \textbf{Le corps} de la lettre, organisé en paragraphes avec alinéas.
\item \textbf{La formule de politesse} finale, en une phrase complète.
\item \textbf{La signature}, précédée éventuellement du nom, en bas à droite.
\end{enumerate}
\end{propriete}

\begin{attention}[Ne jamais oublier l'objet]
L'\cle{objet} est l'élément le plus souvent oublié par les candidats, et le plus grave à oublier : il annonce en une ligne le contenu de la lettre et permet au destinataire de la classer immédiatement. Il se place \textbf{avant la formule d'appel}, à gauche, et commence par le mot « Objet : » suivi d'un groupe nominal (jamais une phrase complète). On écrit « Objet : demande de stage de vacances » et non « Objet : je vous demande un stage ».
\end{attention}

\begin{methode}[Disposer correctement les éléments de la structure externe]
\begin{enumerate}
\item Trace mentalement la page en deux colonnes : colonne de gauche pour l'expéditeur et l'objet, colonne de droite pour la date et le destinataire.
\item Écris en haut à gauche ton identité complète et tes coordonnées (en-tête).
\item À droite, écris le lieu et la date de rédaction, puis, en dessous, la fonction et l'adresse du destinataire.
\item Revient à gauche : écris l'objet en une ligne, puis la formule d'appel adaptée au statut du destinataire.
\item Rédige le corps en paragraphes aérés, avec un alinéa au début de chaque paragraphe.
\item Termine par la formule de politesse, puis signe en bas à droite.
\item Relis et vérifie : aucun élément ne doit manquer, aucune rature ne doit apparaître.
\end{enumerate}
\end{methode}

\begin{popfigure}
\begin{center}
\begin{tikzpicture}[scale=0.92, every node/.style={font=\small}]
% Page
\draw[thick, popdark, fill=popcream] (0,0) rectangle (9.2,12.4);
% En-tête
\node[anchor=north west, text width=4.2cm] (entete) at (0.5,12.1) {\textbf{FOTSO Brice}\\ Lycée technique de Bafoussam\\ BP 123 Bafoussam\\ Tél. : 6 90 00 00 00};
% Date
\node[anchor=north west, text width=4cm] (date) at (5.6,12.1) {Bafoussam, le 12 mai 2026};
% Destinataire
\node[anchor=north west, text width=4.2cm] (dest) at (5.6,10.6) {Monsieur le Directeur\\ Agence Express Union\\ Bafoussam};
% Objet
\node[anchor=north west, text width=6cm] (objet) at (0.5,9.2) {\textbf{Objet :} demande de stage académique.};
% Appel
\node[anchor=north west] (appel) at (0.5,8.4) {Monsieur le Directeur,};
% Corps
\draw[popline] (1.1,7.6) -- (8.7,7.6);
\draw[popline] (0.5,7.1) -- (8.7,7.1);
\draw[popline] (0.5,6.6) -- (8.7,6.6);
\draw[popline] (1.1,5.9) -- (8.7,5.9);
\draw[popline] (0.5,5.4) -- (8.7,5.4);
\draw[popline] (0.5,4.9) -- (8.7,4.9);
\draw[popline] (1.1,4.2) -- (8.7,4.2);
\draw[popline] (0.5,3.7) -- (8.7,3.7);
% Politesse
\node[anchor=north west, text width=7.8cm] (pol) at (0.5,3.3) {Veuillez agréer, Monsieur le Directeur, l'expression de ma haute considération.};
% Signature
\node[anchor=north west] (sign) at (6.2,2.2) {\emph{Signature}};
% Flèches et étiquettes
\draw[-{Stealth}, popblue, thick] (entete.east) -- (10.1,11.4) node[right, popblue, align=left]{\textbf{1. En-tête}\\ (expéditeur)};
\draw[-{Stealth}, poporange, thick] (date.east) -- (10.1,11.0) node[right, poporange, align=left]{\textbf{2. Lieu et date}};
\draw[-{Stealth}, popgreen, thick] (dest.east) -- (10.1,9.9) node[right, popgreen, align=left]{\textbf{3. Destinataire}};
\draw[-{Stealth}, poppurple, thick] (objet.west) -- (-0.6,8.9) node[left, poppurple, align=right]{\textbf{4. Objet}};
\draw[-{Stealth}, poppink, thick] (appel.west) -- (-0.6,8.2) node[left, poppink, align=right]{\textbf{5. Formule}\\ \textbf{d'appel}};
\draw[-{Stealth}, popgold, thick] (8.7,6.2) -- (10.1,6.2) node[right, popgold, align=left]{\textbf{6. Corps}\\ (alinéas)};
\draw[-{Stealth}, popblue, thick] (pol.west) -- (-0.6,2.9) node[left, popblue, align=right]{\textbf{7. Formule de}\\ \textbf{politesse}};
\draw[-{Stealth}, poporange, thick] (sign.south) -- (7.0,0.7) node[below, poporange, align=center]{\textbf{8. Signature}};
\end{tikzpicture}
\end{center}
\legende{Schéma annoté d'une page de lettre officielle : les huit éléments de la structure externe et leur place respective sur la page.}
\end{popfigure}

\courssub{La présentation matérielle : marges, alinéas, pièces jointes}

Une lettre officielle se juge d'abord à l'œil, avant même la lecture. La \cle{présentation matérielle} obéit à quelques règles simples :

\begin{itemize}
\item \textbf{Les marges} : laisse une marge régulière (environ 2 cm) de chaque côté de la page ; la lettre ne doit jamais toucher les bords du papier.
\item \textbf{Les alinéas} : chaque nouveau paragraphe commence par un retrait (un alinéa) ; on passe à la ligne entre l'introduction, le corps et la conclusion.
\item \textbf{L'aération} : évite le « mur de texte » ; une lettre de demande de stage tient en une seule page, en trois ou quatre paragraphes courts.
\item \textbf{Le support} : papier blanc propre, écriture lisible ou texte dactylographié, aucune rature.
\item \textbf{Les pièces jointes} : la demande de stage est souvent accompagnée de documents annexes qu'on signale en bas de la lettre par la mention « Pièces jointes : » (ou « P.J. : ») suivie de la liste : curriculum vitae (CV), relevé de notes, attestation de scolarité.
\end{itemize}

\begin{savaistu}[La demande de stage au Cameroun]
Au Cameroun, une demande de stage déposée dans une entreprise ou une administration est très souvent accompagnée d'une \cle{photocopie de la carte nationale d'identité} (ou de l'acte de naissance pour les mineurs) et d'une \cle{attestation de scolarité} délivrée par l'établissement. Certaines entreprises exigent en plus une lettre de recommandation du proviseur ou du directeur de l'école. Il est donc prudent de préparer un dossier complet avant de se déplacer : cela montre votre sérieux et évite les allers-retours inutiles.
\end{savaistu}

\courssub{La structure interne : introduction, corps, conclusion}

Si la structure externe est le squelette visible, la \cle{structure interne} est l'organisation des idées dans le corps de la lettre. C'est ici que les notions de cohésion et de cohérence, étudiées dans la section précédente, deviennent opératoires : les trois parties doivent s'enchaîner logiquement grâce à des connecteurs bien choisis.

\begin{propriete}[Les trois mouvements de la demande de stage]
\begin{enumerate}
\item \textbf{L'introduction} : elle présente brièvement l'expéditeur (identité, classe, établissement) et annonce le \cle{motif} de la lettre : la demande de stage, avec sa période et sa durée.
\item \textbf{Le corps} : il développe les \cle{arguments} qui rendent la demande recevable : la formation suivie, les compétences déjà acquises, les qualités personnelles (ponctualité, sens du travail en équipe), les disponibilités exactes.
\item \textbf{La conclusion} : elle exprime les \cle{remerciements} et l'\cle{attente d'une suite favorable}, puis enchaîne naturellement sur la formule de politesse.
\end{enumerate}
\end{propriete}

L'enchaînement logique est assuré par des connecteurs adaptés au registre soutenu : « dans le cadre de », « c'est pourquoi », « en effet », « de plus », « en outre », « dans l'espoir que », « en vous remerciant par avance ». Ces outils de cohésion guident le lecteur d'une idée à l'autre sans rupture.

\courssub{Le registre soutenu et les formules consacrées}

La lettre officielle exige le \cle{registre soutenu} : vocabulaire précis et poli, phrases complètes et correctement ponctuées, emploi du « je » assumé mais discret, et surtout des \cle{formules consacrées}, c'est-à-dire des expressions figées que l'usage impose.

\begin{itemize}
\item \textbf{Pour ouvrir la demande} : « J'ai l'honneur de solliciter auprès de votre haute bienveillance un stage au sein de votre entreprise. » ; « Je me permets de vous adresser ma candidature pour un stage de vacances. »
\item \textbf{Pour exprimer l'attente} : « Dans l'espoir que ma demande retiendra votre attention... » ; « En vous remerciant par avance de la suite que vous voudrez bien réserver à ma requête... »
\item \textbf{Pour conclure (formule de politesse)} : elle s'adapte au statut du destinataire, comme le montre le tableau ci-dessous.
\end{itemize}

\begin{center}
\begin{tabularx}{\linewidth}{|l|X|X|}
\hline
\rowcolor{popblueL} \textbf{Destinataire} & \textbf{Formule d'appel} & \textbf{Formule de politesse} \\
\hline
Directeur d'entreprise & Monsieur le Directeur, & Veuillez agréer, Monsieur le Directeur, l'expression de ma haute considération. \\
\hline
Directrice d'entreprise & Madame la Directrice, & Veuillez agréer, Madame la Directrice, l'expression de ma haute considération. \\
\hline
Préfet, ministre, haute autorité & Monsieur le Préfet, / Monsieur le Ministre, & Veuillez agréer, Monsieur le Préfet, l'expression de ma très haute considération. \\
\hline
Proviseur, principal & Monsieur le Proviseur, & Je vous prie d'agréer, Monsieur le Proviseur, l'expression de mon respectueux dévouement. \\
\hline
Destinataire inconnu & Monsieur, / Madame, & Je vous prie d'agréer, Monsieur, Madame, l'expression de mes salutations distinguées. \\
\hline
\end{tabularx}
\end{center}

\begin{attention}[Accorder la formule de politesse avec la formule d'appel]
La formule de politesse doit \textbf{reprendre exactement} le titre utilisé dans la formule d'appel : si vous commencez par « Monsieur le Directeur, », vous devez terminer par « Veuillez agréer, \textbf{Monsieur le Directeur}, ... ». Mélanger « Madame » au début et « Monsieur le Directeur » à la fin est une faute grave qui trahit un manque d'attention. Autre piège : « haute considération » se réserve aux supérieurs hiérarchiques ; pour un destinataire de rang modeste, « salutations distinguées » ou « considération distinguée » suffisent.
\end{attention}

\courssub{Observation guidée de deux modèles authentiques}

Avant de rédiger, observons deux modèles de demandes de stage adressées à des entreprises camerounaises. L'observation guidée consiste à repérer, dans chaque texte, les éléments de la structure externe et les trois mouvements de la structure interne.

\begin{exemplebox}[Texte fonctionnel N°2 : demande de stage à Express Union, Bafoussam]
\textbf{FOTSO Brice}\\
Élève en classe de Première G2, Lycée technique de Bafoussam\\
BP 123 Bafoussam --- Tél. : 6 90 00 00 00\\[4pt]
\hfill Bafoussam, le 12 mai 2026\\[4pt]
\hfill \textbf{Monsieur le Directeur de l'agence Express Union de Bafoussam}\\[4pt]
\textbf{Objet :} demande de stage académique.\\
\textbf{Pièces jointes :} attestation de scolarité, photocopie de la CNI.\\[4pt]
Monsieur le Directeur,\\[4pt]
\hspace{0.6cm} J'ai l'honneur de solliciter auprès de votre haute bienveillance un stage académique au sein de votre agence, pour la période allant du 1\textsuperscript{er} juillet au 31 août 2026.\\[4pt]
\hspace{0.6cm} Actuellement élève en classe de Première G2 (techniques administratives et comptables) au Lycée technique de Bafoussam, je souhaite, dans le cadre de ma formation, découvrir le fonctionnement concret d'une institution de microfinance. De plus, mes cours de comptabilité et d'informatique m'ont donné des bases solides que je voudrais mettre en pratique. Ponctuel et rigoureux, je suis disponible à temps plein pendant toute la durée du stage.\\[4pt]
\hspace{0.6cm} Dans l'espoir que ma demande retiendra votre attention, je vous remercie par avance de la suite que vous voudrez bien lui réserver.\\[4pt]
\hspace{0.6cm} Veuillez agréer, Monsieur le Directeur, l'expression de ma haute considération.\\[6pt]
\hfill \emph{Signature}
\end{exemplebox}

\begin{exemplebox}[Texte fonctionnel N°3 : demande de stage à la SODECOTON, Garoua (extrait)]
\textbf{ABDOULAYE Mariam}\\
Élève en classe de Première F4, Lycée technique de Garoua\\
BP 45 Garoua\\[4pt]
\hfill Garoua, le 20 avril 2026\\[4pt]
\hfill \textbf{Monsieur le Directeur des Ressources Humaines de la SODECOTON, Garoua}\\[4pt]
\textbf{Objet :} demande de stage ouvrier.\\[4pt]
Monsieur le Directeur,\\[4pt]
\hspace{0.6cm} J'ai l'honneur de venir très respectueusement auprès de vous solliciter un stage d'observation dans votre usine, du 15 juin au 15 juillet 2026. Élève en Première F4 (génie mécanique), je désire observer le fonctionnement et la maintenance des machines textiles. En outre, ma formation pratique en atelier scolaire me permettra de m'intégrer rapidement à vos équipes.\\[4pt]
\hspace{0.6cm} En vous remerciant par avance, je vous prie d'agréer, Monsieur le Directeur, l'expression de ma haute considération.\\[6pt]
\hfill \emph{Signature}
\end{exemplebox}

\textbf{Grille d'observation.} Pour chaque texte, vérifiez : (1) les huit éléments externes sont-ils présents et bien placés ? (2) l'introduction annonce-t-elle clairement le motif et la période du stage ? (3) le corps avance-t-il des arguments liés à la formation et aux disponibilités ? (4) la conclusion exprime-t-elle remerciements et attente ? (5) la formule de politesse correspond-elle à la formule d'appel ? Vous constaterez que les deux modèles, bien que destinés à des entreprises différentes (microfinance à Bafoussam, industrie textile à Garoua), respectent exactement la même architecture : c'est la preuve que la lettre officielle est un genre \cle{codifié}, dont la maîtrise s'acquiert par l'observation puis par l'entraînement.

\begin{exoresolu}[Repérer les éléments de la structure externe]
Énoncé. Voici le début d'une lettre : « Monsieur le Directeur, j'ai l'honneur de vous demander un stage. Objet : demande de stage. Douala, le 3 juin 2026. NGA Marie, élève en Première G1. » Relève trois erreurs de disposition et corrige-les.
\tcblower
Solution détaillée. Première erreur : l'\textbf{objet} est placé après la formule d'appel ; il doit se placer \textbf{avant} « Monsieur le Directeur, ». Deuxième erreur : le \textbf{lieu et la date} doivent figurer en haut à droite, avant le corps, et non après la formule d'appel. Troisième erreur : l'\textbf{en-tête} (identité et coordonnées de l'expéditeur) doit ouvrir la lettre, en haut à gauche, et non apparaître après la date. Ordre correct : en-tête de NGA Marie (haut à gauche) ; « Douala, le 3 juin 2026 » (haut à droite) ; coordonnées du destinataire ; « Objet : demande de stage » ; « Monsieur le Directeur, » ; puis le corps commençant par « J'ai l'honneur de... ».
\end{exoresolu}

\begin{exercice}[À vous de jouer]
Tu es élève en Première technique au Lycée technique de Ngaoundéré. Rédige uniquement la \textbf{structure externe complète} (en-tête, lieu et date, destinataire, objet, formule d'appel, formule de politesse, signature) d'une demande de stage adressée au directeur de la CAMTEL de ta ville, puis rédige l'introduction de la lettre en deux ou trois phrases annonçant le motif et la période du stage.
\end{exercice}

\begin{corrige}[Exercice : structure externe et introduction]
Éléments attendus : en haut à gauche, nom et prénom de l'élève, classe, lycée, adresse et téléphone ; en haut à droite, « Ngaoundéré, le ... 2026 » ; à droite, « Monsieur le Directeur de la CAMTEL, Ngaoundéré » ; à gauche, « Objet : demande de stage académique » ; formule d'appel : « Monsieur le Directeur, ». Introduction possible : « J'ai l'honneur de solliciter auprès de votre haute bienveillance un stage académique au sein de votre entreprise, pour la période allant du 1\textsuperscript{er} juillet au 31 août 2026. Élève en classe de Première technique au Lycée technique de Ngaoundéré, je souhaite, dans le cadre de ma formation, découvrir le métier de technicien des télécommunications. » Formule finale : « Veuillez agréer, Monsieur le Directeur, l'expression de ma haute considération », suivie de la signature en bas à droite.
\end{corrige}

\begin{aretenir}[L'essentiel de la section]
La lettre officielle de demande de stage est un écrit formel en registre soutenu, adressé à une autorité ou une entreprise. Sa \textbf{structure externe} compte huit éléments à places fixes : en-tête, lieu et date, destinataire, objet (avant l'appel !), formule d'appel, corps, formule de politesse, signature. Sa \textbf{structure interne} suit trois mouvements : introduction (motif), corps (arguments : formation, compétences, disponibilités), conclusion (remerciements et attente d'une suite favorable). Les formules consacrées (« J'ai l'honneur de... », « Veuillez agréer... ») et l'accord entre formule d'appel et formule de politesse sont obligatoires.
\end{aretenir}

Cette architecture maîtrisée, nous pourrons, dans la section suivante, lire plus finement des textes fonctionnels et comprendre comment le texte explicatif organise l'information, avant de passer nous-mêmes à la rédaction complète de la demande de stage.

\coursec{Lecture de textes fonctionnels et texte explicatif}

Dans la section précédente, nous avons démonté la lettre officielle de demande de stage comme on démonte un moteur : en-tête, objet, formules, signature. Mais avant de \emph{fabriquer} soi-même une telle lettre, il faut savoir \emph{lire} les documents qui circulent autour d'elle : l'annonce de stage publiée par l'entreprise, l'affiche du service des stages, le formulaire à remplir, la note de service qui explique la procédure. Ces documents ne se lisent pas comme un poème ou un roman. C'est toute la question de cette section : comment lire vite et juste un texte \emph{fonctionnel}, et comment reconnaître les procédés d'un texte \emph{explicatif} qui nous guide pas à pas.

\courssub{Le texte fonctionnel : définition et repères}

Imaginez que vous receviez deux enveloppes. La première contient une nouvelle de Mongo Beti ; la seconde contient une convocation à un entretien de stage. Lisez-vous les deux textes de la même façon ? Évidemment non. Pour la nouvelle, vous prenez votre temps, vous savourez le style, vous vous laissez porter. Pour la convocation, vos yeux courent directement vers trois informations : \emph{où ? quand ? que dois-je apporter ?} Cette différence de comportement de lecture repose sur une différence de nature entre les textes.

\begin{definition}[Le texte fonctionnel]
Un \cle{texte fonctionnel} (ou texte utilitaire, texte pragmatique) est un texte dont la fonction principale est de \textbf{déclencher une action} ou de \textbf{transmettre une information précise} à un destinataire déterminé. Il n'est pas écrit pour être admiré, mais pour être \emph{utilisé}. On y trouve : les lettres administratives et commerciales, les annonces, les affiches, les avis, les convocations, les formulaires, les modes d'emploi, les règlements, les notes de service.
\end{definition}

Le texte fonctionnel possède des caractéristiques facilement observables :

\begin{itemize}
\item une \textbf{situation de communication claire} : on sait qui écrit (l'émetteur), à qui (le destinataire) et pour quoi (l'objectif) ;
\item une \textbf{organisation visible} : titres, rubriques, encadrés, listes, caractères gras qui guident l'œil ;
\item un \textbf{vocabulaire précis et souvent administratif} : « dossier », « candidature », « délai », « pièces à fournir », « sous peine de » ;
\item une \textbf{économie de moyens} : phrases courtes, pas de descriptions inutiles, pas d'ambiguïté.
\end{itemize}

\begin{attention}[Un texte fonctionnel ne se lit pas comme un texte littéraire]
L'erreur fréquente consiste à lire une annonce ou une note de service comme on lit un récit : du début à la fin, lentement, en cherchant l'émotion ou le style. Dans un texte fonctionnel, on cherche \textbf{l'information utile}, pas le plaisir. On lit d'abord les éléments mis en évidence (titre, objet, dates, lieux), puis seulement ensuite le détail. Inversement, au Probatoire technique, si l'on vous demande l'objectif d'un texte fonctionnel, ne répondez jamais « faire rêver le lecteur » ou « exprimer les sentiments de l'auteur » : cherchez l'\emph{action} attendue du destinataire.
\end{attention}

\courssub{Méthode de lecture d'un texte fonctionnel}

\begin{methode}[Lire efficacement un texte fonctionnel en trois étapes]
\begin{enumerate}
\item \textbf{Identifier la situation de communication.} Avant même de lire le contenu, repérez : \emph{qui écrit ?} (une entreprise, un ministère, un proviseur), \emph{à qui ?} (un stagiaire, le public, les parents d'élèves), \emph{dans quel cadre ?} (recrutement, information, réclamation). Les indices se trouvent dans l'en-tête, le logo, la signature.
\item \textbf{Dégager l'objectif du texte.} Posez la question : \emph{que doit faire ou savoir le destinataire après lecture ?} Postuler ? Se présenter à telle date ? Payer ? Respecter une consigne ? L'objectif est souvent formulé par un verbe d'action ou un impératif.
\item \textbf{Relever les informations pratiques.} Dates, lieux, horaires, conditions, pièces à fournir, contacts. Ce sont elles qui rendent le texte \emph{efficace} : si l'une manque, le texte échoue dans sa mission.
\end{enumerate}
\end{methode}

Appliquons cette méthode à un cas concret.

\begin{exemplebox}[Analyse guidée d'une annonce de stage]
Voici une annonce publiée par une entreprise de Yaoundé :

\emph{« SOCIÉTÉ CAMTEL — DIRECTION RÉGIONALE DU CENTRE. AVIS DE RECRUTEMENT DE STAGIAIRES. La CAMTEL offre aux élèves des classes de Première et Terminale techniques des stages académiques du 1er juillet au 31 août. Conditions : lettre de motivation adressée au Directeur régional, photocopie de la carte d'identité scolaire, attestation de scolarité. Dépôt des dossiers au secrétariat, immeuble CAMTEL, quartier Elig-Essono, avant le 15 juin à 15 h 30, délai de rigueur. »}

Appliquons les trois étapes :
\begin{enumerate}
\item \textbf{Situation} : émetteur = la CAMTEL (Direction régionale du Centre) ; destinataires = les élèves de Première et Terminale techniques ; cadre = recrutement de stagiaires.
\item \textbf{Objectif} : pousser les élèves à \emph{déposer un dossier de candidature}.
\item \textbf{Informations pratiques} : période du stage (1er juillet – 31 août), pièces à fournir (3 documents), lieu de dépôt (secrétariat, Elig-Essono), date limite (15 juin, 15 h 30), sanction implicite (« délai de rigueur » = aucun dossier ne sera accepté après).
\end{enumerate}
En moins d'une minute, le lecteur sait exactement quoi faire : c'est la marque d'un texte fonctionnel réussi.
\end{exemplebox}

\courssub{Étude du Texte fonctionnel N°2 : structure et vocabulaire administratif}

Le Texte fonctionnel N°2 de votre manuel est une \textbf{lettre administrative} (par exemple une réponse d'une entreprise à une demande de stage). Lisons-le méthodiquement.

\textbf{Repérage des éléments de structure externe.} On y retrouve tous les repères étudiés dans la section précédente :

\begin{itemize}
\item l'\textbf{en-tête} : nom et coordonnées de l'entreprise (gauche), nom et coordonnées du destinataire (droite) ;
\item le \textbf{lieu et la date} : « Yaoundé, le... » ;
\item les \textbf{références} : « N° 245/DR/CAM/26 » et « V/réf. : votre lettre du 10 mai » ;
\item l'\textbf{objet} : « Réponse à votre demande de stage » ;
\item la \textbf{formule d'appel} : « Monsieur, » ;
\item le \textbf{corps} en deux ou trois paragraphes brefs ;
\item la \textbf{formule de politesse} et la \textbf{signature}.
\end{itemize}

\textbf{Repérage du vocabulaire administratif.} Relevons les expressions typiques : « \emph{nous accusons réception de} », « \emph{nous avons l'honneur de vous informer que} », « \emph{nous vous prions de bien vouloir} », « \emph{sous toutes réserves} », « \emph{pièces jointes} ». Ce vocabulaire codifié a deux fonctions : il garantit la \textbf{politesse} (distance respectueuse entre administration et usager) et la \textbf{précision juridique} (chaque formule a un sens exact).

\begin{aretenir}[Le tableau de lecture méthodique]
Pour tout texte fonctionnel, remplissez mentalement (ou sur votre copie) ce tableau à quatre cases :

\begin{center}
\begin{tabularx}{\linewidth}{|l|X|}
\hline
\rowcolor{popblueL} \textbf{Question} & \textbf{Réponse relevée dans le texte} \\
\hline
Qui écrit ? (émetteur) & Le Directeur régional de la CAMTEL \\
\hline
À qui ? (destinataire) & L'élève auteur de la demande de stage \\
\hline
Pour quoi ? (objectif) & L'informer que sa demande est acceptée et lui fixer la date de prise de fonction \\
\hline
Informations clés & Date de début (1er juillet), service d'affectation (maintenance des réseaux), documents à présenter le premier jour \\
\hline
\end{tabularx}
\end{center}

Ce tableau est un outil d'examen : il vous permet de répondre en quelques lignes à la question « Quel est l'objectif de ce texte ? ».
\end{aretenir}

\courssub{Étude du Texte fonctionnel N°3 : comparaison et efficacité}

Le Texte fonctionnel N°3 poursuit le même objectif général (communiquer une information utile), mais dans une \textbf{forme différente} : par exemple une affiche ou un avis affiché au tableau d'information du lycée, là où le N°2 était une lettre nominative. Comparons-les.

\begin{center}
\begin{tabularx}{\linewidth}{|l|X|X|}
\hline
\rowcolor{popblueL} \textbf{Critère} & \textbf{Texte N°2 (lettre)} & \textbf{Texte N°3 (affiche/avis)} \\
\hline
Destinataire & Une seule personne, nommée & Un collectif (tous les élèves) \\
\hline
Registre & Soutenu, formules de politesse & Direct, impératifs, phrases nominales \\
\hline
Mise en page & Paragraphes rédigés & Gros caractères, rubriques, listes \\
\hline
Informations & Détaillées, personnalisées & Réduites à l'essentiel (quoi, qui, quand, où) \\
\hline
Efficacité & Réponse individuelle précise & Diffusion rapide à un large public \\
\hline
\end{tabularx}
\end{center}

\textbf{Conclusion de la comparaison.} Les deux textes sont fonctionnels, mais leur \emph{registre} et leur \emph{mise en page} s'adaptent au destinataire et au canal de diffusion. Un bon lecteur repère immédiatement ces choix ; un bon rédacteur saura les reproduire. La \textbf{clarté} reste le critère commun : dans les deux cas, une information ambiguë (une date floue, un lieu imprécis) rendrait le texte inutile.

\courssub{Le texte explicatif : visée et procédés}

Certains textes fonctionnels ne se contentent pas d'informer : ils doivent \emph{faire comprendre} quelque chose de nouveau ou de compliqué. C'est le cas de la note de service qui explique la procédure de demande de stage, ou du mode d'emploi d'une machine de l'atelier.

\begin{definition}[Le texte explicatif]
Un \cle{texte explicatif} est un texte dont la visée est de \textbf{faire comprendre} au lecteur une notion, un phénomène ou une procédure. Pour cela, il mobilise des \textbf{procédés explicatifs} : la \emph{définition}, l'\emph{exemple}, la \emph{comparaison}, la \emph{reformulation} et l'enchaînement \emph{cause-conséquence}. On le reconnaît à des marques linguistiques : « c'est-à-dire », « par exemple », « autrement dit », « en effet », « c'est pourquoi », « de sorte que ».
\end{definition}

\begin{popfigure}
\begin{tikzpicture}[grow=right, level distance=3.6cm,
  level 1/.style={sibling distance=3.4cm},
  level 2/.style={sibling distance=1.1cm, level distance=4.6cm},
  every node/.style={font=\small},
  edge from parent/.style={draw=popmuted, thick}]
\node[draw=popblue, fill=popblueL, rounded corners, thick, align=center] {\textbf{Procédés}\\\textbf{explicatifs}}
  child { node[draw=popgreen, fill=popgreenL, rounded corners, align=center] {Cause-\\conséquence}
    child { node[align=left, font=\scriptsize] {« Le dossier est incomplet,\\ \emph{de sorte que} la demande\\ est rejetée. »} } }
  child { node[draw=poppink, fill=poppinkL, rounded corners] {Reformulation}
    child { node[align=left, font=\scriptsize] {« Délai de rigueur,\\ \emph{c'est-à-dire} aucune\\ prolongation possible. »} } }
  child { node[draw=poporange, fill=poporangeL, rounded corners] {Comparaison}
    child { node[align=left, font=\scriptsize] {« La lettre de motivation\\ est \emph{comme} une carte\\ de visite écrite. »} } }
  child { node[draw=popgold, fill=popgoldL, rounded corners] {Exemple}
    child { node[align=left, font=\scriptsize] {« Joignez vos pièces,\\ \emph{par exemple} l'attestation\\ de scolarité. »} } }
  child { node[draw=poppurple, fill=poppurpleL, rounded corners] {Définition}
    child { node[align=left, font=\scriptsize] {« Un stage \emph{est} une période\\ de formation pratique\\ en entreprise. »} } };
\end{tikzpicture}
\legende{Arbre de classification des cinq procédés explicatifs, avec un exemple pour chacun. Les marques linguistiques (en italique) permettent de repérer chaque procédé dans un texte.}
\end{popfigure}

\courssub{Repérage des procédés explicatifs dans un texte de procédure}

Lisons maintenant ce court texte explicatif, du type de celui qu'un proviseur peut afficher pour guider les élèves :

\emph{« Pour obtenir un stage, l'élève doit d'abord rédiger une demande de stage, c'est-à-dire une lettre officielle par laquelle il sollicite un accueil en entreprise. Cette lettre est comme une carte de visite : elle donne la première impression du candidat. Elle doit comporter, par exemple, l'objet de la demande et la période souhaitée. Comme les entreprises reçoivent de nombreuses demandes, tout dossier incomplet est écarté ; c'est pourquoi il faut vérifier chaque pièce avant le dépôt. Autrement dit, la rigueur est la condition du succès. »}

\begin{exoresolu}[Repérer les procédés explicatifs]
\textbf{Énoncé.} Relevez dans le texte ci-dessus cinq procédés explicatifs différents, citez le passage correspondant et nommez la marque linguistique qui le signale.
\tcblower
\textbf{Solution détaillée.}
\begin{enumerate}
\item \textbf{Définition} : « une demande de stage, \emph{c'est-à-dire} une lettre officielle par laquelle il sollicite un accueil en entreprise » — la marque « c'est-à-dire » introduit la définition du terme technique.
\item \textbf{Comparaison} : « cette lettre est \emph{comme} une carte de visite » — la marque « comme » rapproche deux réalités pour éclairer la première.
\item \textbf{Exemple} : « elle doit comporter, \emph{par exemple}, l'objet de la demande » — « par exemple » illustre une règle générale par un cas précis.
\item \textbf{Cause-conséquence} : « \emph{comme} les entreprises reçoivent de nombreuses demandes, tout dossier incomplet est écarté ; \emph{c'est pourquoi} il faut vérifier chaque pièce » — la subordonnée de cause et le connecteur « c'est pourquoi » enchaînent logiquement les faits.
\item \textbf{Reformulation} : « \emph{autrement dit}, la rigueur est la condition du succès » — « autrement dit » reprend l'idée en termes plus simples.
\end{enumerate}
\end{exoresolu}

\begin{attention}[Ne pas confondre expliquer et raconter]
Le texte explicatif répond à la question \emph{« comment ? »} ou \emph{« pourquoi ? »}, pas à la question \emph{« que s'est-il passé ? »}. Si votre analyse parle de personnages, de péripéties ou de sentiments, vous êtes en train de décrire un texte narratif, pas un texte explicatif. Autre piège : la simple énumération (« d'abord, ensuite, enfin ») organise une procédure mais n'explique rien à elle seule ; l'explication suppose un lien logique (cause, définition, illustration).
\end{attention}

\begin{exercice}[À vous de lire]
Voici un avis affiché au lycée : « \emph{Les élèves de Première technique désirant effectuer un stage de vacances sont invités à retirer une fiche de vœux au secrétariat des élèves, bâtiment B, du lundi au vendredi, de 10 h à 12 h. La fiche remplie, signée du parent, sera déposée avant le 30 mai. Aucune demande tardive ne sera examinée.} »
\begin{enumerate}
\item Remplissez le tableau de lecture méthodique (émetteur, destinataire, objectif, informations clés).
\item Ce texte est-il plutôt de type informatif ou explicatif ? Justifiez en une phrase.
\item Quelle marque de mise en page aurait renforcé son efficacité ?
\end{enumerate}
\end{exercice}

\begin{corrige}[Exercice « À vous de lire »]
\begin{enumerate}
\item Émetteur : l'administration du lycée (secrétariat des élèves). Destinataires : les élèves de Première technique. Objectif : les amener à retirer puis déposer une fiche de vœux de stage. Informations clés : lieu (bâtiment B), horaires (10 h – 12 h, jours ouvrables), pièce exigée (signature du parent), date limite (30 mai), sanction (demande tardive rejetée).
\item Texte plutôt \textbf{informatif} : il transmet des consignes pratiques sans définir ni justifier ; il n'emploie aucun connecteur explicatif (« c'est-à-dire », « par exemple », « c'est pourquoi »).
\item Une mise en page en rubriques ou en liste (Quoi / Qui / Où / Quand) et la date limite en caractères gras auraient rendu l'information immédiatement visible.
\end{enumerate}
\end{corrige}

\begin{aretenir}[Bilan de la section]
Un \cle{texte fonctionnel} vise une action ou une information ; on le lit en trois étapes (situation, objectif, informations pratiques). Un \cle{texte explicatif} vise à faire comprendre grâce à cinq procédés repérables par leurs marques linguistiques : définition, exemple, comparaison, reformulation, cause-conséquence. Ces deux compétences de lecture préparent directement la rédaction de votre propre demande de stage : vous savez désormais lire l'annonce à laquelle vous répondrez, et expliquer clairement votre projet dans le corps de votre lettre.
\end{aretenir}

\coursec{La communication par l'image}

Dans la section précédente, nous avons appris à lire des textes fonctionnels et à repérer les procédés du texte explicatif. Or, dans la vie professionnelle, un texte ne circule presque jamais seul : une annonce de recrutement, une plaquette d'entreprise, une demande de stage accompagnée d'un curriculum vitae, tout cela s'accompagne d'\cle{images} — logo, photographie, affiche, pictogramme. Savoir lire une image est donc aussi important que savoir lire un texte. C'est l'objet de cette section : comprendre comment une image \emph{communique}, c'est-à-dire comment elle transmet un message, et apprendre à l'analyser méthodiquement.

\courssub{L'image comme message : dénotation et connotation}

\textbf{Intuition.} Regardons une affiche placardée à l'entrée d'une grande entreprise de Douala : on y voit un jeune technicien en tenue de travail, souriant, devant une machine moderne, avec le slogan « Construis ton avenir avec nous ». Tout le monde voit la même chose : un jeune homme, une machine, des couleurs vives. Mais chacun \emph{comprend} davantage : l'entreprise est moderne, elle accueille les jeunes, le travail y est valorisant. L'image dit donc deux choses à la fois : ce qu'elle \emph{montre} et ce qu'elle \emph{suggère}.

\begin{definition}[Dénotation et connotation]
La \cle{dénotation} est le premier niveau de lecture d'une image : c'est la description \textbf{objective} de ce qui est représenté (les objets, les personnages, les couleurs, les lieux), sans interprétation. La \cle{connotation} est le second niveau de lecture : c'est l'ensemble des \textbf{idées, valeurs et émotions} que l'image suggère au spectateur (modernité, sérieux, bonheur, richesse, confiance...).
\end{definition}

\begin{exemplebox}[Dénotation et connotation sur une affiche]
Sur l'affiche de l'entreprise de Douala décrite plus haut :
\begin{itemize}
\item \textbf{Dénotation} : un jeune homme portant une tenue bleue et un casque, debout devant une machine, souriant ; fond orange ; texte blanc.
\item \textbf{Connotation} : le sourire suggère la satisfaction au travail ; la machine moderne suggère le progrès technique ; la tenue propre suggère le professionnalisme ; le fond orange suggère l'énergie et le dynamisme.
\end{itemize}
\end{exemplebox}

\begin{attention}[Ne pas confondre description et interprétation]
L'erreur la plus fréquente à l'examen est de mélanger les deux niveaux. Dire « le jeune homme est heureux » est déjà une interprétation : objectivement, on voit seulement qu'il \emph{sourit}. Dans une analyse rigoureuse, commencez toujours par décrire ce qui est visible (dénotation), puis seulement proposez une interprétation (connotation), en la justifiant par un élément précis de l'image.
\end{attention}

\begin{popfigure}
\begin{center}
\begin{tikzpicture}[font=\small]
% Cadre de l'affiche
\draw[fill=poporangeL, draw=popdark, thick, rounded corners] (0,0) rectangle (5,6.5);
% Soleil / fond
\draw[fill=popgoldL, draw=popgold] (3.9,5.5) circle (0.5);
% Machine
\draw[fill=popblueL, draw=popblue, thick] (2.6,0.6) rectangle (4.5,2.4);
\draw[fill=popblue, draw=popblue] (3.0,2.4) rectangle (4.1,2.9);
\draw[fill=white, draw=popblue] (2.9,1.2) circle (0.35);
% Personnage
\draw[fill=popcream, draw=popdark] (1.2,3.4) circle (0.45);
\draw[fill=popgold, draw=popdark] (0.7,3.95) -- (1.7,3.95) -- (1.55,4.15) -- (0.85,4.15) -- cycle;
\draw[fill=popblue, draw=popdark, thick] (0.65,1.1) rectangle (1.75,2.95);
\draw[popdark, thick] (1.0,1.1) -- (1.0,0.3);
\draw[popdark, thick] (1.4,1.1) -- (1.4,0.3);
% Sourire
\draw[popdark] (1.05,3.25) arc (200:340:0.2);
% Slogan
\draw[fill=popdark, rounded corners] (0.3,5.0) rectangle (4.7,6.2);
\node[align=center, font=\bfseries, text=white] at (2.5,5.6) {CONSTRUIS TON AVENIR\\AVEC NOUS};
% Flèches d'analyse
\draw[-{Stealth}, very thick, popblue] (5.1,3.0) -- (6.6,3.0);
\node[align=left, anchor=west, text=popblue, font=\bfseries] at (6.7,3.9) {DÉNOTATION};
\node[align=left, anchor=west, text=popdark] at (6.7,3.2) {jeune technicien, casque,\\machine, fond orange,\\slogan blanc};
\draw[-{Stealth}, very thick, poppurple] (5.1,1.6) -- (6.6,1.6);
\node[align=left, anchor=west, text=poppurple, font=\bfseries] at (6.7,2.3) {CONNOTATION};
\node[align=left, anchor=west, text=popdark] at (6.7,1.3) {modernité, dynamisme,\\accueil des jeunes,\\avenir prometteur};
\end{tikzpicture}
\legende{Schéma de lecture d'une affiche de recrutement de stagiaires : la dénotation décrit ce qui est visible (à gauche), la connotation dégage ce que l'image suggère (à droite).}
\end{center}
\end{popfigure}

\courssub{Les éléments d'analyse d'une image}

Pour analyser une image de manière complète, on examine cinq éléments : le \cle{cadrage}, les \cle{couleurs}, la \cle{lumière}, les \cle{personnages} et le \cle{texte associé}.

\begin{itemize}
\item \textbf{Le cadrage} : c'est le choix de ce qui entre dans l'image et de la place occupée par chaque élément. Un gros plan sur un visage insiste sur l'émotion ; un plan large sur un atelier montre l'ampleur de l'activité. Un personnage placé au centre est présenté comme le sujet principal.
\item \textbf{Les couleurs} : elles portent des valeurs symboliques. Le bleu évoque la confiance et la technicité ; le vert, la nature, la santé ou l'agriculture ; le rouge, l'urgence ou l'énergie ; le doré, la réussite. Les couleurs vives attirent le regard, les couleurs sobres suggèrent le sérieux.
\item \textbf{La lumière} : une image lumineuse suggère l'optimisme et la transparence ; une image sombre ou contrastée peut créer de la gravité ou du mystère.
\item \textbf{Les personnages} : leur âge, leur tenue, leur gestuelle, leur regard (vers le spectateur ou ailleurs) construisent l'identification. Un jeune en tenue de technicien permet au lycéen de se projeter.
\item \textbf{Le texte associé} : la \cle{légende} explique ou oriente la lecture ; le \cle{slogan} résume le message en une formule courte et frappante. Texte et image se complètent : l'image attire, le texte fixe le sens.
\end{itemize}

\begin{center}
\begin{tabularx}{\linewidth}{|l|X|X|}
\hline
\rowcolor{popblueL} \textbf{Élément visuel} & \textbf{Question à se poser} & \textbf{Effet produit (exemple)} \\
\hline
Cadrage & Qu'est-ce qui est montré en premier, en grand, au centre ? & Gros plan sur le stagiaire : le jeune est valorisé. \\
\hline
Couleurs & Quelles couleurs dominent ? Que symbolisent-elles ? & Bleu dominant : confiance, sérieux technique. \\
\hline
Lumière & L'image est-elle claire, sombre, contrastée ? & Lumière vive : optimisme, transparence de l'entreprise. \\
\hline
Personnages & Qui est représenté ? Âge, tenue, geste, regard ? & Jeune souriant en tenue de travail : identification du lecteur. \\
\hline
Texte associé & Que disent la légende et le slogan ? & « Construis ton avenir » : promesse adressée au jeune. \\
\hline
\end{tabularx}
\end{center}

\begin{methode}[Analyser une image en quatre questions]
Pour analyser n'importe quelle image (affiche, logo, photographie, couverture de plaquette), pose-toi successivement quatre questions :
\begin{enumerate}
\item \textbf{Que voit-on ?} Décris objectivement l'image : c'est la dénotation (éléments, couleurs, cadrage, texte).
\item \textbf{Qui ?} À qui s'adresse l'image ? Qui l'a produite et pour quel public ?
\item \textbf{Quel effet ?} Quelles émotions ou idées les choix visuels provoquent-ils ? C'est la connotation, justifiée par des éléments précis.
\item \textbf{Quel message ?} Formule en une phrase le sens global de l'image et sa fonction (informer, convaincre, valoriser).
\end{enumerate}
\end{methode}

\courssub{Les fonctions de l'image dans la communication professionnelle}

Dans le monde professionnel, l'image n'est jamais décorative : elle remplit des \cle{fonctions} précises, que l'on peut regrouper en trois grandes catégories.

\begin{definition}[Fonctions de l'image professionnelle]
Une image professionnelle peut :
\begin{itemize}
\item \textbf{informer} : donner une information rapidement et clairement (pictogramme de sécurité, plan d'accès, organigramme) ;
\item \textbf{convaincre} : pousser le lecteur à agir ou à adhérer (affiche de recrutement, publicité, campagne de sensibilisation) ;
\item \textbf{valoriser} : donner une image positive d'une entreprise, d'un produit ou d'une personne (logo, plaquette d'entreprise, photo de couverture d'un rapport annuel).
\end{itemize}
Une même image peut cumuler plusieurs fonctions : une affiche de recrutement de stagiaires informe (postes proposés, contacts), convainc (rejoindre l'entreprise) et valorise (image moderne de l'entreprise).
\end{definition}

\begin{exemplebox}[Analyse d'un logo et de son slogan]
Prenons le logo imaginaire de l'entreprise camerounaise « \textbf{Nyambé BTP} », spécialisée dans le bâtiment : un triangle orange stylisé formant un toit, surmontant le nom de l'entreprise en lettres bleues, avec le slogan « Bâtir solide, bâtir durable ».
\begin{itemize}
\item \textbf{Dénotation} : un triangle orange (forme de toit), un texte bleu, un slogan en italique.
\item \textbf{Connotation} : le toit renvoie directement au métier du bâtiment ; le triangle, figure stable, suggère la solidité ; l'orange suggère l'énergie ; le bleu, la confiance.
\item \textbf{Fonction} : valoriser l'entreprise (image de fiabilité) et convaincre le client.
\item \textbf{Le slogan} : la répétition de « bâtir » et l'allitération en « b » donnent un rythme martelé qui évoque la régularité du travail bien fait ; « solide » et « durable » reprennent les valeurs suggérées par le triangle.
\end{itemize}
On observe que le slogan et le logo \emph{se renforcent mutuellement} : c'est l'articulation texte/image.
\end{exemplebox}

\begin{savaistu}[Le poids de l'image dans le recrutement]
Des études de communication montrent qu'une annonce de recrutement accompagnée d'une image soignée (logo visible, photo de l'équipe, mise en page claire) retient l'attention des lecteurs beaucoup plus longtemps qu'une annonce composée de texte seul, et reçoit nettement plus de candidatures. Au Cameroun, les grandes entreprises et les organisations internationales soignent particulièrement le visuel de leurs campagnes de recrutement de stagiaires, car l'affiche est souvent le \emph{premier contact} entre l'entreprise et le futur stagiaire. Toi aussi, quand tu rédigeras ta demande de stage, la présentation matérielle de ta lettre (propreté, mise en page, lisibilité) jouera le rôle d'une « image » : elle donnera immédiatement une impression de sérieux ou de négligence.
\end{savaistu}

\courssub{Analyse guidée d'une affiche de campagne de recrutement de stagiaires}

Appliquons maintenant la méthode des quatre questions à un document complet : l'affiche de la campagne « Stage d'été 2026 » de la société fictive « Cameroun Agro-Industries » (CAI), de Yaoundé.

\textbf{Description de l'affiche.} En haut, le logo de l'entreprise (une feuille verte dans un engrenage orange) et le titre en grosses lettres blanches sur bandeau vert : « STAGE D'ÉTÉ 2026 — TECHNICIENS DE DEMAIN ». Au centre, une photographie : trois jeunes stagiaires (deux garçons, une fille) en blouse blanche et lunettes de protection, souriants, autour d'une machine de conditionnement qu'un encadrant leur montre. En bas, un encadré : « Tu es en Première ou Terminale technique ? Rejoins-nous du 1\textsuperscript{er} juillet au 31 août. Envoie ta demande de stage à : recrutement@cai.cm ». Le fond est clair, lumineux, avec des touches de vert et d'orange.

\begin{exoresolu}[Analyse complète de l'affiche « Stage d'été 2026 »]
Énoncé : analysez cette affiche en suivant la méthode des quatre questions, en distinguant dénotation et connotation, et concluez sur ses fonctions.
\tcblower
\textbf{1. Que voit-on ? (dénotation)} On voit un bandeau vert avec un titre blanc, un logo associant une feuille et un engrenage, une photographie centrale de trois jeunes en blouse blanche avec un encadrant devant une machine, et un encadré d'informations pratiques (dates, adresse électronique). Les couleurs dominantes sont le vert, l'orange et le blanc ; la lumière est vive.

\textbf{2. Qui ?} L'émetteur est l'entreprise Cameroun Agro-Industries. Le destinataire est explicitement désigné : « Tu es en Première ou Terminale technique ? » — donc les élèves des lycées techniques, tutoyés pour créer la proximité.

\textbf{3. Quel effet ? (connotation)} Le logo associe nature (la feuille) et industrie (l'engrenage) : l'entreprise se présente comme moderne et respectueuse de l'agriculture. Les blouses blanches et les lunettes suggèrent l'hygiène, la sécurité et le sérieux du travail. Les sourires et la présence d'une jeune fille suggèrent une ambiance accueillante et ouverte à tous. Le tutoiement et le mot « demain » dans le titre valorisent le jeune : il est présenté comme l'avenir de l'entreprise. Le vert évoque la croissance, l'orange l'énergie, la lumière vive la transparence.

\textbf{4. Quel message ?} L'affiche dit : « Notre entreprise est moderne, accueillante et formatrice ; toi, jeune élève technicien, tu peux y construire ton avenir : postule. »

\textbf{Fonctions} : l'affiche \emph{informe} (dates, conditions, contact), \emph{convainc} (slogan, tutoiement, image positive) et \emph{valorise} l'entreprise (modernité, encadrement de qualité). L'articulation texte/image est efficace : la photographie rend crédible ce que le slogan promet.
\end{exoresolu}

\begin{exercice}[À toi d'analyser]
Ton lycée reçoit l'affiche d'une banque qui recrute des stagiaires : on y voit un jeune guichetier en costume, serrant la main d'un client, avec le slogan « La confiance, notre métier » et les coordonnées du service des stages. Analyse cette affiche avec les quatre questions de la méthode, en rédigeant un paragraphe de dénotation, un paragraphe de connotation et une conclusion sur les fonctions de l'image.
\end{exercice}

\begin{corrige}[Exercice — éléments de réponse]
\textbf{Dénotation} : un homme jeune en costume sombre serre la main d'un client ; couleurs sobres (bleu, gris) ; slogan en lettres blanches ; coordonnées en bas de l'affiche.
\textbf{Connotation} : le costume et la poignée de main suggèrent le sérieux, le respect du client et la confiance ; les couleurs sobres renforcent l'idée de rigueur et de sécurité, valeurs attendues d'une banque ; la jeunesse du guichetier permet l'identification des élèves.
\textbf{Fonctions} : informer (coordonnées, offre de stage), convaincre (slogan centré sur la confiance), valoriser (image de professionnalisme). Le slogan et la poignée de main photographiée se renforcent : l'image illustre littéralement le mot « confiance ».
\end{corrige}

\courssub{Complément utile au Bac : l'articulation texte/image}

\begin{aretenir}[Texte et image : trois relations possibles — utile pour le commentaire au Bac]
Au Probatoire et au Baccalauréat technique, tu peux rencontrer, dans un sujet de commentaire ou de réception de texte fonctionnel, un document mêlant texte et image (affiche, couverture de plaquette, annonce). Trois relations sont possibles :
\begin{enumerate}
\item \textbf{l'illustration} : l'image montre ce que le texte dit (la photo du stagiaire illustre l'annonce) ;
\item \textbf{le renforcement} : l'image ajoute une valeur émotionnelle ou symbolique au message du texte (le sourire, les couleurs) ;
\item \textbf{le complément} : le texte précise ce que l'image ne peut pas dire seule (dates, conditions, contacts).
\end{enumerate}
Dans ta copie, signale toujours cette articulation : c'est une remarque de niveau supérieur qui valorise ton analyse. Formule attendue : « L'image ne se contente pas d'illustrer le texte, elle le renforce en suggérant... ».
\end{aretenir}

\begin{attention}[Pièges de rédaction à l'examen]
\begin{itemize}
\item Ne rédige jamais une simple liste d'éléments (« il y a un logo, il y a du vert, il y a un slogan ») : relie chaque élément à l'effet qu'il produit.
\item N'affirme pas une connotation sans la justifier : écris « le vert \emph{suggère} la nature », pas « le vert \emph{est} la nature ».
\item N'oublie pas le destinataire : une image n'a de sens que pour un public visé ; demande-toi toujours à qui elle parle.
\item Dans une demande de stage, la « communication par l'image » se joue aussi dans ta propre lettre : écriture lisible, marges, disposition des blocs (coordonnées, objet, formules) donnent de toi une image avant même la lecture du contenu.
\end{itemize}
\end{attention}

En résumé, une image professionnelle est un message construit : elle se lit d'abord objectivement (dénotation), puis s'interprète (connotation) grâce à l'analyse du cadrage, des couleurs, de la lumière, des personnages et du texte associé. Cette compétence te servira doublement : pour analyser les documents professionnels à l'examen, et pour soigner la présentation de ta propre demande de stage, objet de la prochaine section.

\coursec{Rédaction de la demande de stage : introduction, corps, conclusion}

Dans la séquence précédente, nous avons analysé la \cle{communication par l'image} (photo, logo, schéma) comme support de persuasion immédiate. Pour la demande de stage, c'est le \cle{texte écrit}, structuré et argumenté, qui endosse ce rôle. Une lettre officielle n'est pas un empilement de phrases : c'est un édifice normé. Elle repose sur deux piliers indissociables : la \cle{structure externe} (mise en page, en-tête, objet, références) et la \cle{structure interne} (introduction, corps, conclusion). C'est la maîtrise de cette double architecture, garantie de \cle{cohésion} (liaisons linguistiques) et de \cle{cohérence} (logique des idées), que nous allons construire pas à pas.

\courssub{La structure externe : les normes de présentation}

Avant même de lire le contenu, le destinataire juge la forme. Le respect des codes de la lettre administrative conditionne la crédibilité du candidat.

\begin{definition}[Structure externe de la lettre officielle]
Ensemble des éléments matériels et codifiés qui encadrent le texte : en-tête (expéditeur, destinataire, date, lieu), objet, références éventuelles, corps de la lettre, formule de politesse, signature et mentions de pièces jointes.
\end{definition}

\begin{methode}[Mise en page type (normes administratives camerounaises)]
\begin{enumerate}
\item \textbf{En-tête expéditeur} (angle sup. gauche) : Nom, Prénom, Classe, Établissement, Adresse, Téléphone, Email.
\item \textbf{En-tête destinataire} (sous l'expéditeur, aligné à gauche) : Qualité + Nom de l'organisme (ex. : Monsieur le Directeur général de la CAMTEL), Adresse de l'organisme.
\item \textbf{Lieu et date} (aligné à droite, sous le destinataire) : Ex. : Douala, le 15 mai 2025.
\item \textbf{Objet} (aligné à gauche, en gras, souligné) : Précis et complet. Ex. : \textbf{Objet : Demande de stage académique (Génie électrique) — Juillet-Août 2025}.
\item \textbf{Références} (si réponse à une annonce) : Réf. : Votre annonce n°... du...
\item \textbf{Formule d'appel} (salutation) : « Monsieur le Directeur général, » (suivie d'une virgule, pas de point).
\item \textbf{Corps de la lettre} (structure interne : introduction, corps, conclusion).
\item \textbf{Formule de politesse} (closing) : Reprend \textbf{exactement} le titre du destinataire.
\item \textbf{Signature} : Manuscrite, précédée du nom tapé en majuscules.
\item \textbf{Pièces jointes} (P.J.) : Liste en bas à gauche (CV, attestation de scolarité, lettre de recommandation...).
\end{enumerate}
\end{methode}

\begin{attention}[Pièges de la mise en page]
\begin{itemize}
\item Ne jamais écrire « Objet : Demande de stage » seul : il faut le type, la spécialité, la période.
\item La date s'écrit en toutes lettres (jour en chiffres, mois en lettres, année en chiffres) : « Yaoundé, le 3 juin 2025 ».
\item L'objet ne prend pas de point final.
\item La formule d'appel ne contient pas le nom propre du destinataire (on écrit « Monsieur le Directeur général, », pas « Monsieur Dupont, »).
\end{itemize}
\end{attention}

\courssub{Analyser la situation d'écriture avant de rédiger}

Rédiger sans analyser la situation de communication, c'est construire sans plan. Le bon rédacteur identifie quatre paramètres invariants.

\begin{definition}[Situation d'écriture (ou situation de communication)]
Ensemble des contraintes qui déterminent les choix linguistiques et structurels : \textbf{qui} écrit (émetteur), \textbf{à qui} (destinataire), \textbf{pour quoi} (finalité), \textbf{dans quelles conditions} (registre, support, délai, longueur).
\end{definition}

Appliquons cette grille à la demande de stage en Première technique :

\begin{center}
\begin{tabularx}{\linewidth}{|l|X|}
\hline
\rowcolor{popblueL} \textbf{Paramètre} & \textbf{Configuration pour la demande de stage} \\
\hline
\textbf{Émetteur} & Élève de Première technique (spécialité identifiée), locuteur légitime, demandeur. \\
\hline
\textbf{Destinataire} & Décideur hiérarchique (Directeur général, Chef de service, DRH), inconnu, pouvoir d'accepter ou refuser. \\
\hline
\textbf{Finalité} & Obtenir un stage académique/professionnel $\rightarrow$ provoquer une décision favorable. \\
\hline
\textbf{Contraintes} & Registre soutenu obligatoire ; brièveté (1 page max) ; structure externe/internes rigides ; orthographe zéro faute ; ton respectueux mais ferme. \\
\hline
\end{tabularx}
\end{center}

\begin{attention}[Registre soutenu : zone rouge]
La familiarité est éliminatoire. Sont \textbf{interdits} : tutoiement, argot, abréviations SMS, formules toutes faites familières (« Salut », « Ça va ? », « Votre boîte », « Je kifferais », « Merci d'avance, tchao »).
\textbf{Exigence} : Vouvoiement systématique, vocabulaire précis (« solliciter » et non « demander », « bienveillance », « considération »), syntaxe complexe (subjonctif, conditionnel de politesse : « je souhaiterais », « je vous saurais gré »).
\end{attention}

\courssub{La structure interne : l'ossature en trois temps}

La structure interne organise la démonstration. Elle suit un mouvement dialectique : \textbf{Exposer} (Introduction) $\rightarrow$ \textbf{Argumenter} (Corps) $\rightarrow$ \textbf{Conclure/Engager} (Conclusion).

\begin{popfigure}
\begin{center}
\begin{tikzpicture}[
 bloc/.style={draw=popblue, thick, fill=popblueL, rounded corners, minimum width=4.8cm, minimum height=3.0cm, align=center, font=\small},
 titre/.style={font=\bfseries\small, text=popdark},
 fleche/.style={-{Stealth}, very thick, popmuted}]
\node[bloc, align=center] (intro) at (0,0){\textbf{1. INTRODUCTION}\\[4pt] \textit{« Accrocher et situer »}\\[2pt] — Présentation identitaire\\ — Objet explicite\\ — Précisions (type, durée, période, service)};
\node[bloc, fill=poporangeL, draw=poporange, align=center] (corps) at (6.5,0){\textbf{2. CORPS}\\[4pt] \textit{« Convaincre »}\\[2pt] — Formation \& compétences\\ — Motivations ciblées\\ — Disponibilités \& pièces jointes\\ \textit{Ordre : général $\rightarrow$ précis}};
\node[bloc, fill=popgreenL, draw=popgreen, align=center] (concl) at (13,0){\textbf{3. CONCLUSION}\\[4pt] \textit{« Laisser une trace positive »}\\[2pt] — Remerciement\\ — Espoir de suite favorable\\ — Formule de politesse \textbf{identique} au titre};
\draw[fleche] (intro) -- (corps);
\draw[fleche] (corps) -- (concl);
\end{tikzpicture}
\end{center}
\legende{Architecture fonctionnelle de la lettre de demande de stage : trois séquences obligatoires, chacune à finalité distincte. La flèche indique la progression logique.}
\end{popfigure}

\courssub{Rédiger l'introduction : se présenter et annoncer l'objet}

L'introduction est le « sas » de lecture. Le destinataire (souvent pressé) y décide de la suite. Elle doit être \textbf{concise} (3 à 4 phrases max), \textbf{exhaustive} sur l'identité et l'objet, \textbf{polie}.

\begin{methode}[Introduction en trois temps (modèle canonique)]
\begin{enumerate}
\item \textbf{Phrase 1 — Identité} : Nom, Prénom, Classe, Spécialité, Établissement. \\ \textit{Modèle :} « Je soussigné(e) \textit{Nom Prénom}, élève en classe de \textit{Première [Spécialité]} au \textit{Lycée technique de [Ville]}, ... »
\item \textbf{Phrase 2 — Objet} : Verbe de politesse + nature de la demande. \\ \textit{Modèle :} « ... viens très respectueusement solliciter un stage académique au sein de votre structure / de votre service [Nom du service]. »
\item \textbf{Phrase 3 — Cadre spatio-temporel} : Durée, dates prévisionnelles, service visé. \\ \textit{Modèle :} « Ce stage, d'une durée de \textit{X semaines}, se déroulerait du \textit{date début} au \textit{date fin}, de préférence au service \textit{[Nom]}. »
\end{enumerate}
\end{methode}

\begin{exemplebox}[Introduction aboutie — Stage CAMTEL Bamenda (Génie électrique)]
Je soussigné Ndi Collins, élève en classe de Première Génie électrique au Collège technique de Bamenda, viens très respectueusement auprès de votre haute bienveillance solliciter un stage académique au sein de Cameroon Telecommunications (CAMTEL). Ce stage, d'une durée de six semaines, se déroulerait du 1er juillet au 15 août prochains, de préférence dans le service des réseaux et de la maintenance.
\end{exemplebox}

\begin{aretenir}[Check-list introduction]
\begin{itemize}
\item[\checkmark] Identité complète (nom, classe, spécialité, établissement).
\item[\checkmark] Verbe de sollicitation au conditionnel ou indicatif poli (« sollicite », « souhaiterais »).
\item[\checkmark] Précisions chiffrées (durée) et temporelles (dates).
\item[\checkmark] Service ciblé (montre que l'élève s'est renseigné).
\item[\checkmark] Zéro familiarité, zéro faute.
\end{itemize}
\end{aretenir}

\courssub{Rédiger le corps : argumenter avec cohérence et cohésion}

Le corps est la \cle{preuve}. Il transforme une demande en candidature. Trois piliers argumentatifs, ordonnés du général au particulier, reliés par des \cle{connecteurs logiques} (garants de la cohésion) dans une progression logique (garante de la cohérence).

\begin{methode}[Construction du corps : 4 paragraphes-types]
\begin{enumerate}
\item \textbf{Paragraphe Formation (Général)} : Classe, spécialité, matières dominantes, niveau. \textit{Connecteurs :} « Actuellement », « En effet », « Ainsi ».
\item \textbf{Paragraphe Compétences (Précis)} : Savoir-faire techniques (TP, logiciels, machines, langues), stages antérieurs, projets. \textit{Connecteurs :} « De plus », « En outre », « Par ailleurs », « Notamment ».
\item \textbf{Paragraphe Motivations (Ciblé)} : Pourquoi \textbf{cette} entreprise ? (Réputation, équipements, domaine d'activité). Pourquoi \textbf{ce} stage ? (Lien avec projet professionnel). \textit{Connecteurs :} « C'est pourquoi », « Parce que », « Étant donné que », « Afin de ».
\item \textbf{Paragraphe Disponibilités / Praticité} : Dates confirmées, mobilité, pièces jointes (CV, attestation), demande d'entretien. \textit{Connecteurs :} « Enfin », « Par conséquent », « Je reste à votre disposition pour ».
\end{enumerate}
\end{methode}

\begin{aretenir}[Boîte à outils : Connecteurs logiques pour lettre officielle]
\begin{center}
\begin{tabularx}{\linewidth}{|l|X|}
\hline
\rowcolor{popblueL} \textbf{Fonction rhétorique} & \textbf{Connecteurs (registre soutenu)} \\
\hline
\textbf{Énumération / Ajout} & Premièrement, en premier lieu ; de surcroît, qui plus est, de surcroît ; enfin \\
\hline
\textbf{Explication / Cause} & En effet, du fait que, eu égard à, dans la mesure où \\
\hline
\textbf{Conséquence / Conclusion partielle} & Par conséquent, il en résulte que, c'est la raison pour laquelle \\
\hline
\textbf{But / Finalité} & Dans le dessein de, en vue de, dans l'optique de \\
\hline
\textbf{Concession / Nuance} & Certes, force est de constater que, toutefois, néanmoins \\
\hline
\textbf{Transition} & S'agissant de..., Quant à..., En ce qui concerne... \\
\hline
\end{tabularx}
\end{center}
\end{aretenir}

\begin{exemplebox}[Corps de lettre structuré — Exemple CAMTEL]
Actuellement en classe de Première Génie électrique, je suis un enseignement axé sur les installations électriques industrielles et l'électronique de puissance. \textbf{En effet}, les modules de technologie de l'équipement et de dessin technique industriel m'ont permis d'acquérir des bases solides en lecture de schémas unifilaires et multifilaires. \textbf{De surcroît}, les travaux pratiques de câblage d'armoires électriques et de dépannage de moteurs asynchrones, réalisés dans les ateliers de mon établissement, m'ont conféré une aisance pratique que je souhaite consolider.

\textbf{C'est pourquoi} je porte mon choix sur la CAMTEL, opérateur historique dont l'infrastructure de transmission par fibre optique et le réseau de commutation constituent un terrain d'apprentissage de premier plan pour un futur technicien en réseaux. \textbf{Dans cette perspective}, ce stage représenterait l'opportunité de confronter mes acquis théoriques aux réalités de la maintenance préventive et curative sur site.

\textbf{Enfin}, je suis disponible pour la totalité de la période indiquée et peux me rendre quotidiennement à votre direction régionale. \textbf{Par conséquent}, je joins à la présente mon curriculum vitae détaillé ainsi qu'une attestation de scolarité en cours de validité, et me tiens à votre entière disposition pour tout entretien que vous jugerez utile.
\end{exemplebox}

\begin{attention}[Cohésion vs Cohérence : ne pas confondre]
\begin{itemize}
\item \textbf{Cohésion} = \textit{Texture linguistique} : les mots se tiennent (connecteurs, reprises anaphoriques « ce stage », « cette entreprise », champs lexicaux). C'est visible \textit{entre} les phrases.
\item \textbf{Cohérence} = \textit{Architecture logique} : les idées s'enchaînent sans contradiction, du général au particulier, sans répétition ni trou. C'est visible \textit{à l'échelle du paragraphe et du texte}.
\item \textit{Une lettre peut être cohésive (beaux connecteurs) mais incohérente (arguments dans le désordre). L'idéal : les deux.}
\end{itemize}
\end{attention}

\courssub{Rédiger la conclusion : remercier, espérer, saluer}

La conclusion est la dernière trace laissée dans l'esprit du lecteur. Elle suit un rituel immuable en trois temps, sans improvisation.

\begin{methode}[Conclusion en trois temps (rituel administratif)]
\begin{enumerate}
\item \textbf{Remerciement} : Formule de gratitude pour l'attention portée à la lecture. \textit{Ex. :} « Je vous remercie de l'attention que vous voudrez bien accorder à ma demande. »
\item \textbf{Expression de l'espoir} : Manifestation de la volonté d'obtenir une réponse positive. \textit{Ex. : } « Dans l'espoir d'une suite favorable, » ou « Dans l'attente d'une réponse que j'espère positive, ».
\item \textbf{Formule de politesse (Closing)} : \textbf{Règle d'or : Reprise \textit{verbatim} du titre de l'en-tête destinataire.} \textit{Ex. :} « Je vous prie d'agréer, \textbf{Monsieur le Directeur régional}, l'expression de ma haute considération. »
\end{enumerate}
\end{methode}

\begin{exemplebox}[Conclusion type — Stage CNIC Douala]
Je vous remercie sincèrement de l'attention que vous voudrez bien réserver à ma candidature. Dans l'espoir d'une suite favorable, je vous prie d'agréer, \textbf{Monsieur le Directeur général}, l'expression de ma haute considération.
\end{exemplebox}

\begin{attention}[Erreurs fatales sur la formule de politesse]
\begin{itemize}
\item \textbf{Désaccord titre/en-tête} : En-tête « Monsieur le Directeur général » $\neq$ Clôture « Monsieur le Directeur ». \textbf{Sanction :} Rejet immédiat perçu comme négligence.
\item \textbf{Mélange de formules} : « Je vous prie de veuillez agréer » (faute de syntaxe). Correct : « Je vous prie d'agréer » OU « Veuillez agréer ».
\item \textbf{Oubli de la virgule} après le titre dans la formule : « Je vous prie d'agréer Monsieur le Directeur général l'expression... » $\rightarrow$ Illisible.
\item \textbf{Adjectif possessif} : « Ma haute considération » (pas « la haute considération »).
\end{itemize}
\end{attention}

\courssub{La relecture finale : protocole de vérification (Zéro faute)}

Une lettre technique avec des fautes signale un technicien imprécis. La relecture suit une grille en quatre passes.

\begin{methode}[Protocole de relecture en 4 passes]
\begin{enumerate}
\item \textbf{Passage 1 — Orthographe lexicale} : Mots pièges (sollicitez, bienveillance, considération, agréer, pièce jointe, curriculum vitæ, attestation). Vérifier les doubles consonnes, accents.
\item \textbf{Passage 2 — Accords grammaticaux} : Sujet/Verbe (je me permets, je joins), Participes passés (la demande que j'ai \textbf{faite}, les pièces \textbf{jointes}), Adjectifs (une suite \textbf{favorable}, ma \textbf{haute} considération).
\item \textbf{Passage 3 — Ponctuation \& Typographie} : Virgule après formule d'appel, majuscules en début de phrase et noms propres (Cameroun, CAMTEL, CNIC), guillemets français « » si citation, espaces insécables avant : ; ? ! et après « ».
\item \textbf{Passage 4 — Registre \& Cohérence globale} : Remplacer tout terme familier (bosser $\rightarrow$ travailler/exercer ; super $\rightarrow$ remarquable/pertinent ; stage $\rightarrow$ stage académique/professionnel). Vérifier l'enchaînement Introduction $\rightarrow$ Corps $\rightarrow$ Conclusion.
\end{enumerate}
\end{methode}

\begin{exoresolu}[Correction guidée : du familier à l'officiel]
\textbf{Énoncé.} Corrigez et restructurez cet extrait pour en faire une introduction et une conclusion conformes : \\
« Salut Monsieur, moi c'est Etoundi je suis en première au lycée technique de Yaoundé. Je voudrais faire un stage chez vous parce que votre boîte elle est connue. Merci d'avance et à plus. »

\tcblower
\textbf{Solution détaillée.}

\textbf{Étape 1 — Nettoyage (suppression familier) :} « Salut », « moi c'est », « je voudrais », « chez vous », « votre boîte », « elle est connue », « merci d'avance », « à plus ».

\textbf{Étape 2 — Reconstruction Introduction (3 phrases) :}
« Je soussigné \textbf{Etoundi Jean}, élève en classe de \textbf{Première technique} au \textbf{Lycée technique de Yaoundé}, viens très respectueusement \textbf{solliciter un stage académique} au sein de votre entreprise. \textbf{Ce stage, d'une durée de quatre semaines, se déroulerait au cours du mois de juillet prochain, idéalement au service maintenance.} »

\textbf{Étape 3 — Reconstruction Conclusion (3 temps) :}
« \textbf{Je vous remercie} de l'attention que vous voudrez bien porter à ma demande. \textbf{Dans l'espoir d'une suite favorable}, je vous prie d'agréer, \textbf{Monsieur le Directeur}, l'expression de ma \textbf{haute considération}. »

\textbf{Analyse :} Le titre « Monsieur le Directeur » (supposé dans l'en-tête) est repris \textit{à l'identique}. Le registre est soutenu (« soussigné », « solliciter », « bienveillance » implicite par « respectueusement », « haute considération »).
\end{exoresolu}

\begin{exercice}[Production autonome : Demande de stage au CNIC]
\textbf{Consigne.} Vous êtes \textbf{Kamga Patrice}, élève en \textbf{Première Chaudronnerie-Soudure} au \textbf{Lycée technique de Douala}. Rédigez la \textbf{lettre complète} (structure externe + introduction + corps + conclusion) adressée à \textbf{Monsieur le Directeur général des Chantiers Navals et Industriels du Cameroun (CNIC)}, BP 123 Douala, pour un \textbf{stage de cinq semaines} pendant les grandes vacances (1er juillet – 5 août), au \textbf{service construction navale}.

\textbf{Contraintes :} Respect strict de la mise en page (en-têtes, objet, date, P.J.) ; 3 arguments dans le corps reliés par connecteurs variés ; formule de politesse exacte ; registre soutenu ; tenue de page soignée.
\end{exercice}

\begin{corrige}[Exercice : Demande de stage au CNIC — Grille d'évaluation]
\textbf{Structure externe (4 pts) :} Expéditeur complet / Destinataire exact (M. le Directeur général CNIC, BP 123 Douala) / Lieu/Date (Douala, le [date]) / Objet précis (Demande de stage académique — Chaudronnerie — Juillet-Août 2025) / Formule d'appel correcte.

\textbf{Introduction (4 pts) :} Identité (Kamga Patrice, 1ère Chaudronnerie, LT Douala) / Objet (stage académique CNIC) / Précisions (5 sem., 1er juil.–5 aoû, service construction navale).

\textbf{Corps — Argumentation \& Cohérence (8 pts) :}
\begin{itemize}
\item \textit{Arg. 1 Formation :} Spécialité Chaudronnerie-Soudure, lecture de plans, technologie de soudage (connecteur : « Actuellement », « En effet »).
\item \textit{Arg. 2 Compétences :} Maîtrise soudure à l'arc/TIG/MIG, traçage/découpe tôles, normes sécurité (connecteur : « De surcroît », « Notamment »).
\item \textit{Arg. 3 Motivation :} CNIC = référence navale nationale, chantiers complexes, encadrement qualifié (connecteur : « C'est pourquoi », « Eu égard à »).
\item \textit{Arg. 4 Praticité :} Disponibilité dates exactes, mobilité, P.J. (CV, attestation), demande entretien (connecteur : « Enfin », « Par conséquent »).
\end{itemize}

\textbf{Conclusion (3 pts) :} Remerciement / Espoir / Formule exacte : « Je vous prie d'agréer, \textbf{Monsieur le Directeur général}, l'expression de ma haute considération. »

\textbf{Langue \& Registre (6 pts) :} Zéro faute orthographe/grammaire / Registre soutenu constant / Connecteurs variés et pertinents / Ponctuation rigoureuse / Présentation soignée.
\textbf{Total : 25 points.}
\end{corrige}

\begin{aretenir}[Synthèse opérationnelle : La lettre de demande de stage en 10 points-clés]
\begin{enumerate}
\item \textbf{Analyse préalable} : Qui ? À qui ? Pourquoi ? Contraintes ? (Registre soutenu).
\item \textbf{Structure externe} : En-têtes alignés, Objet précis, Date en lettres, Appel codifié.
\item \textbf{Introduction (3 phrases)} : Identité $\rightarrow$ Objet $\rightarrow$ Cadre (durée, dates, service).
\item \textbf{Corps (4 paragraphes)} : Formation $\rightarrow$ Compétences $\rightarrow$ Motivations ciblées $\rightarrow$ Disponibilités/P.J.
\item \textbf{Connecteurs} : Variés, adaptés à la fonction (cause, conséquence, but, ajout, concession).
\item \textbf{Cohésion} : Reprises, champs lexicaux, chaines de référence.
\item \textbf{Cohérence} : Progression logique (général $\rightarrow$ particulier), pas de redite, pas de trou.
\item \textbf{Conclusion (3 temps)} : Remerciement $\rightarrow$ Espoir $\rightarrow$ Formule de politesse \textbf{calquée sur l'en-tête}.
\item \textbf{Relecture 4 passes} : Lexicale / Grammaticale / Typographique / Registre.
\item \textbf{Signature + P.J.} : Manuscrit + Nom majuscules + Liste pièces (CV, Attestation, etc.).
\end{enumerate}
\end{aretenir}

\coursec{Production, confrontation, amélioration et remédiation}

Après avoir maîtrisé la rédaction de chaque partie de la lettre — introduction, corps et conclusion — nous entrons dans la phase décisive : produire une lettre complète, la confronter au regard d'un pair, l'améliorer, puis mesurer nos acquis pour combler les lacunes résiduelles. Cette démarche, au cœur de la pédagogie par projet, transforme l'élève en auteur réflexif, capable d'autoévaluation et de remédiation autonome.

\courssub{Rédaction individuelle complète d'une demande de stage}

Chaque élève reçoit une situation d'écriture authentique : une entreprise identifiée (nom, secteur, localisation), un poste visé (intitulé, service), une durée (dates précises) et son propre profil (formation, spécialité, compétences techniques). La consigne est explicite : « Rédigez une lettre de demande de stage adressée au directeur de l'entreprise, en respectant la structure externe et interne, le registre formel et les conventions de politesse. »

\begin{methode}[Rédiger sa lettre complète en quatre temps]
\begin{enumerate}
\item \textbf{Analyser la situation} : identifier le destinataire, l'objet, la durée, les arguments pertinents (compétences techniques, motivation, disponibilité).
\item \textbf{Construire le squelette} : noter sur brouillon les trois parties (introduction, corps, conclusion) avec les connecteurs logiques (\cle{d'abord}, \cle{ensuite}, \cle{enfin}, \cle{par ailleurs}, \cle{de plus}).
\item \textbf{Rédiger au net} : soigner la présentation (marges, alignement, saut de ligne entre paragraphes), inscrire la date, l'objet, les formules d'appel et de politesse adaptées.
\item \textbf{Relire à voix basse} : vérifier la cohérence (chaque paragraphe répond à une intention), la cohésion (pronoms, connecteurs, champs lexicaux), l'orthographe grammaticale et lexicale.
\end{enumerate}
\end{methode}

\begin{exemplebox}[Situation d'écriture type]
\textbf{Contexte :} Vous êtes élève en Première Technicien en Maintenance des Systèmes Automatisés (MSA) au Lycée Technique de Bafoussam. Vous sollicitez un stage de 8 semaines (du 2 juillet au 26 août 2025) à la SONARA (Société Nationale de Raffinage), service Maintenance Instrumentation, BP 123 Limbé.
\textbf{Consigne :} Rédigez la lettre complète sur une feuille blanche, en 20 minutes.
\end{exemplebox}

\courssub{Confrontation : lecture croisée et grille d'évaluation}

La confrontation en binôme est un levier puissant : l'élève devient correcteur, ce qui affine son propre regard critique. Chaque binôme dispose d'une grille critériée commune, transparente et notée sur 20 points.

\begin{popfigure}
\legende{Grille d'évaluation critériée de la lettre de demande de stage (20 points)}
\begin{center}
\begin{tabularx}{\linewidth}{|l|X|c|}\hline
\rowcolor{popblueL}
\textbf{Critère} & \textbf{Indicateurs observables} & \textbf{Points} \\ \hline
Structure externe & En-tête complet (expéditeur, destinataire, date, lieu), Objet clair, Formules d'appel et de politesse conformes, Signature & 4 \\ \hline
Structure interne & Introduction (accroche, objet, source), Corps (2--3 paragraphes argumentés : formation, compétences, motivation), Conclusion (disponibilité, demande d'entretien, formule de politesse finale) & 6 \\ \hline
Cohérence & Progression logique, pas de répétition inutile, chaque paragraphe a une fonction identifiée, ton constant & 4 \\ \hline
Cohésion & Connecteurs variés et pertinents, chaînes de référence (pronoms, synonymes), champs lexicaux professionnels & 3 \\ \hline
Langue & Registre formel soutenu, syntaxe correcte, orthographe grammaticale et lexicale, ponctuation rigoureuse & 3 \\ \hline
\end{tabularx}
\end{center}
\end{popfigure}

\begin{methode}[Lecture croisée en binôme : mode d'emploi]
\begin{enumerate}
\item \textbf{Échange} : l'élève A lit la lettre de B à voix haute ; B note les impressions globales.
\item \textbf{Grille} : A remplit la grille pour B (coche les indicateurs, note sur 20, écrit 2 forces et 2 axes d'amélioration en bas de page).
\item \textbf{Retour} : A restitue oralement son évaluation à B (3 min) ; puis inversion des rôles.
\item \textbf{Synthèse} : chaque élève reporte sur sa propre lettre les remarques pertinentes (surlignage, notes marginales).
\end{enumerate}
\end{methode}

\begin{attention}[Piège fréquent : la notation bienveillante mais floue]
Éviter les commentaires vagues (« c'est bien », « il faut améliorer »). Exiger des repérages précis : « Paragraphe 2 : connecteur "mais" mal employé, remplacer par "cependant" », « Formule de politesse finale : "Veuillez agréer" au lieu de "Cordialement" ». La grille force cette précision.
\end{attention}

\courssub{Amélioration : réécriture de la lettre}

Fort des retours, l'élève réécrit sa lettre. C'est le moment où la prise de conscience devient compétence. La réécriture n'est pas une simple correction : c'est une refonte guidée par les critères.

\begin{exemplebox}[Lettre d'élève corrigée : avant / après amélioration]
\textbf{Avant (extrait du corps) :}
« Je suis en première MSA. J'aime bien la maintenance. J'ai fait des stages à l'école. Je suis motivé. Je veux apprendre chez vous. »
\textbf{Remarques de la grille :} registre familier, pas de connecteurs, arguments génériques, pas de vocabulaire technique.
\textbf{Après (réécriture) :}
« Actuellement élève en Première Technicien Maintenance des Systèmes Automatisés, j'ai acquis des compétences solides en diagnostic de pannes sur automates programmables (Siemens S7-1200) et en lecture de schémas électriques normalisés NF C 15-100. \textbf{Par ailleurs}, mes travaux pratiques sur banc de test m'ont permis de développer rigueur et autonomie. \textbf{C'est pourquoi} je souhaite vivement mettre ces acquis au service de votre service Maintenance Instrumentation. »
\textbf{Gains :} vocabulaire technique précis, connecteurs logiques (\cle{par ailleurs}, \cle{c'est pourquoi}), registre formel, argumentation structurée.
\end{exemplebox}

\begin{savaistu}[Lecture à voix haute : le meilleur détecteur de ruptures]
Relire sa lettre à voix haute (ou la faire lire par un camarade) permet d'entendre les phrases trop longues, les répétitions, les changements de ton brutaux. L'oreille perçoit ce que l'œil laisse passer : une rupture de cohérence se manifeste souvent par une hésitation à la lecture. Essayez : vous serez surpris du nombre de corrections spontanées qui surgissent.
\end{savaistu}

\courssub{Étude du Texte fonctionnel N°4 : modèle de référence pour l'autoévaluation}

Le Texte fonctionnel N°4 est une lettre de demande de stage rédigée par un professionnel (ou un ancien élève admis en stage), conforme aux normes administratives et rhétoriques. Son étude guidée permet à chaque élève de se positionner par rapport à un standard d'excellence.

\begin{exoresolu}[Analyse comparative : ma lettre vs Texte fonctionnel N°4]
\textbf{Extrait du Texte fonctionnel N°4 (introduction) :}
« Monsieur le Directeur, Actuellement élève en Première Technicien Électrotechnique au Lycée Technique de Douala-Bassa, je me permets de solliciter un stage de perfectionnement d'une durée de huit semaines, du 1er juillet au 31 août 2025, au sein de votre service Maintenance Haute Tension. »
\textbf{Points de repère pour l'autoévaluation :}
\begin{itemize}
\item Formule d'appel : « Monsieur le Directeur, » (pas « Monsieur, » seul).
\item Présentation de soi : classe, spécialité, établissement — en une phrase fluide.
\item Objet précis : stage de perfectionnement, dates exactes, service ciblé.
\item Ton : respectueux, confiant, sans familiarité.
\end{itemize}
\textbf{Exercice :} Sur votre lettre, soulignez au crayon vert les éléments conformes, au rouge les écarts. Comptez : plus de vert que de rouge ? Vous êtes sur la bonne voie.
\end{exoresolu}

\courssub{Compte rendu de la séquence : bilan des acquis et difficultés}

Le compte rendu collectif (10 min) structure la méta-cognition. Le professeur anime un tableau de synthèse au tableau, alimenté par les élèves.

\begin{center}
\begin{tabularx}{\linewidth}{|l|X|X|}\hline
\rowcolor{popblueL}
\textbf{Dimension} & \textbf{Acquis majoritaires} & \textbf{Difficultés récurrentes} \\ \hline
Structure externe & En-tête, date, objet, formules d'appel & Formule de politesse finale inadaptée (« Cordialement »), signature oubliée \\ \hline
Structure interne & Introduction présente, conclusion présente & Corps souvent réduit à un seul paragraphe, arguments peu techniques \\ \hline
Cohérence & Progression globale comprise & Transitions entre paragraphes faibles ou absentes \\ \hline
Cohésion & Utilisation de « par ailleurs », « de plus » & Surutilisation de « aussi », confusion « donc » / « par conséquent » \\ \hline
Langue & Registre globalement formel & Fautes d'accord (participes passés, adjectifs), imprécision terminologique (« stage » seul au lieu de « stage de perfectionnement ») \\ \hline
\end{tabularx}
\end{center}

\begin{popfigure}
\legende{Répartition fictive des erreurs fréquentes observées dans la classe (effectif 42 élèves)}
\begin{tikzpicture}
\begin{axis}[
    ybar,
    bar width=0.6cm,
    width=\linewidth,
    height=9cm,
    ymin=0, ymax=30,
    ylabel={Nombre d'élèves concernés},
    symbolic x coords={Politesse, Signature, Connecteurs, VocabTech, Accords, StructureCorps},
    xtick=data,
    xticklabels={Formule politesse finale, Signature/Date, Connecteurs, Vocabulaire technique, Accords participes, Structure corps (1 seul paragraphe)},
    x tick label style={rotate=30, anchor=east, font=\small},
    ymajorgrids=true,
    grid style=dashed,
    nodes near coords,
    nodes near coords align={vertical},
    every node near coord/.append style={font=\small, popdark},
    title style={font=\bfseries, popdark},
    title={Erreurs les plus fréquentes dans les premières versions}
]
\addplot[popblue, fill=popblue!60] coordinates {
    (Politesse, 28)
    (Signature, 22)
    (Connecteurs, 19)
    (VocabTech, 17)
    (Accords, 15)
    (StructureCorps, 14)
};
\end{axis}
\end{tikzpicture}
\end{popfigure}

\courssub{Remédiation : exercices ciblés sur les points faibles}

La remédiation n'est pas une punition : c'est un entraînement ciblé, différencié selon le profil de chaque élève (repéré via la grille et le diagramme). Elle se déploie en trois ateliers tournants (15 min chacun).

\begin{methode}[Construire son plan de remédiation personnel]
\begin{enumerate}
\item \textbf{Repérer} : sur ta grille d'autoévaluation, coche les 2 critères où ta note est la plus basse.
\item \textbf{Choisir} : sélectionne l'atelier correspondant (voir ci-dessous).
\item \textbf{Pratiquer} : réalise les 3 exercices de l'atelier en 15 min, corrige avec la fiche réponse.
\item \textbf{Valider} : réécris le passage faible de ta lettre en appliquant la règle travaillée.
\item \textbf{Archiver} : colle la fiche d'exercice corrigée dans ton cahier de remédiation ; note la date et la règle-clé.
\end{enumerate}
\end{methode}

\begin{exercice}[Atelier 1 : Formules de politesse et signature]
\textbf{Exercice 1.} Complétez les formules d'appel et de politesse finale adaptées à chaque destinataire :
\begin{enumerate}
\item Directeur général : \textit{Monsieur le Directeur général, \dots \quad \dots \quad Veuillez agréer, Monsieur le Directeur général, l'expression de ma considération distinguée.}
\item Chef de service Maintenance : \textit{Monsieur le Chef de service, \dots \quad \dots \quad Je vous prie d'agréer, Monsieur le Chef de service, l'assurance de ma parfaite considération.}
\item Directrice des Ressources Humaines : \textit{Madame la Directrice, \dots \quad \dots \quad Veuillez agréer, Madame la Directrice, l'expression de mes salutations respectueuses.}
\end{enumerate}
\textbf{Exercice 2.} Parmi ces formules, laquelle est \textbf{inadaptée} à une lettre officielle de demande de stage ? Justifiez.
\begin{itemize}
\item Cordialement,
\item Bien à vous,
\item Veuillez agréer, Monsieur, mes salutations distinguées.
\item Respectueusement,
\end{itemize}
\textbf{Exercice 3.} Réécrivez la conclusion de cette lettre en ajoutant la date, le lieu et la signature manuscrite (mentionnez « Signature » entre crochets).
\end{exercice}

\begin{corrige}[Exercice Atelier 1]
\textbf{Exercice 2.} « Cordialement » et « Bien à vous » sont des formules de courrier électronique ou de correspondance semi-formelle, non admises dans une lettre administrative officielle. « Respectueusement » s'emploie seul en bas de page, sans la formule développée, mais reste moins formel que « Veuillez agréer... ». La seule pleinement correcte est la troisième.
\textbf{Exercice 3.} Fait à Douala, le 15 mai 2025 \quad [Signature]
\end{corrige}

\begin{exercice}[Atelier 2 : Connecteurs logiques et cohésion]
\textbf{Exercice 1.} Reliez les deux phrases par le connecteur le plus pertinent (choisissez dans la liste : \textit{par ailleurs, de plus, c'est pourquoi, en effet, cependant, par conséquent}).
\begin{enumerate}
\item Je maîtrise la lecture de schémas hydrauliques. \dots\dots\dots\dots\dots, j'ai réalisé le câblage d'une armoire de commande sur banc d'essai.
\item Le stage me permettra de découvrir l'environnement industriel. \dots\dots\dots\dots\dots, il favorisera mon insertion professionnelle future.
\item Je n'ai pas encore utilisé d'automate Schneider. \dots\dots\dots\dots\dots, ma formation sur Siemens me donne des bases transférables.
\end{enumerate}
\textbf{Exercice 2.} Remplacez les répétitions soulignées par un pronom ou un synonyme :
« Mon stage \textbf{de maintenance} m'a permis d'acquérir des compétences en \textbf{maintenance} préventive. Lors de ce \textbf{stage}, j'ai aussi participé à la \textbf{maintenance} curative. »
\end{exercice}

\begin{corrige}[Exercice Atelier 2]
\textbf{Exercice 1.} 1. \textit{De plus} (addition d'arguments) 2. \textit{Par conséquent} (conséquence) 3. \textit{Cependant} (concession/opposition).
\textbf{Exercice 2.} « Mon stage de maintenance m'a permis d'acquérir des compétences en maintenance préventive. Lors de \textbf{celui-ci}, j'ai aussi participé à la \textbf{curative}. » (ou « à cette dernière »).
\end{corrige}

\begin{exercice}[Atelier 3 : Vocabulaire technique et structure du corps]
\textbf{Exercice 1.} Associez chaque compétence technique à l'intitulé de stage correspondant :
\begin{center}
\begin{tabular}{|l|l|}\hline
\textbf{Compétence} & \textbf{Stage visé} \\ \hline
Programmation d'automates (Siemens, Schneider) & A. Maintenance des systèmes automatisés \\ \hline
Lecture de plans d'implantation électrique HT/BT & B. Électrotechnique / Distribution d'énergie \\ \hline
Soudure TIG/MIG sur acier inoxydable & C. Chaudronnerie / Soudure \\ \hline
Diagnostic de pannes sur variateurs de vitesse & A. Maintenance des systèmes automatisés \\ \hline
\end{tabular}
\end{center}
\textbf{Exercice 2.} Rédigez un paragraphe unique (5--7 lignes) pour le corps de la lettre, en valorisant \textbf{deux} compétences techniques de votre spécialité, avec au moins \textbf{trois} connecteurs différents. Utilisez le vocabulaire précis de votre filière.
\end{exercice}

\begin{corrige}[Exercice Atelier 3]
\textbf{Exercice 1.} Correspondances : 1--A, 2--B, 3--C, 4--A.
\textbf{Exercice 2.} (Exemple pour MSA) « \textbf{En premier lieu}, mes travaux pratiques sur automate Siemens S7-1200 m'ont permis de programmer des séquences Grafcet pour la gestion d'une chaîne de convoyage. \textbf{Par ailleurs}, j'ai effectué le diagnostic de pannes sur variateurs de fréquence Altivar 312, en exploitant les codes défaut et les courbes de courant. \textbf{Ces expériences}, conjuguées à ma rigueur dans le respect des consignes de sécurité (cadenassage, vérification d'absence de tension), me permettront d'être opérationnel rapidement au sein de votre service. »
\end{corrige}

\begin{aretenir}[Bilan de la séquence « Demande de stage »]
\begin{enumerate}
\item La lettre officielle de demande de stage est un texte \textbf{fonctionnel} à structure codifiée (externe + interne) : tout écart la rend irrecevable.
\item La \textbf{cohérence} (logique des idées) et la \textbf{cohésion} (liaisons linguistiques) sont évaluées conjointement : une lettre bien structurée mais mal connectée perd en efficacité.
\item Le \textbf{registre formel soutenu} s'impose : vocabulaire technique précis, syntaxe complexe maîtrisée, formules de politesse ritualisées.
\item La \textbf{réécriture guidée par une grille critériée} est le levier principal de progression : elle objectifie l'évaluation et oriente la remédiation.
\item L'autoévaluation par comparaison à un \textbf{modèle de référence} (Texte fonctionnel N°4) développe l'autonomie de l'élève rédacteur.
\end{enumerate}
\end{aretenir}

\begin{attention}[Erreur administrative éliminatoire]
Une lettre sans \textbf{date} (lieu + date du jour) ou sans \textbf{signature manuscrite} (prénom + nom) est \textbf{irrecevable} par tout service administratif ou RH. Même un contenu parfait ne compensera pas cette omission. Vérifiez systématiquement ces deux éléments avant de remettre votre copie ou d'envoyer le courrier.
\end{attention}

\newpage

\coursec{Activité d'intégration}

\courssub{Objectifs de l'activité}
\noindent Cette activité d'intégration vise à mobiliser l'ensemble des savoirs et savoir-faire travaillés au cours de la leçon n°10. À son terme, l'élève doit être capable de :
\begin{itemize}
    \item \textbf{Analyser} une situation de communication professionnelle réelle (affiche d'appel à candidatures) en identifiant l'émetteur, le destinataire, l'objet, les contraintes formelles et le canal.
    \item \textbf{Produire} une lettre officielle de demande de stage complète, respectant scrupuleusement la \textbf{structure externe} (en-têtes, références, objet, formules de politesse, signature) et la \textbf{structure interne} (introduction, corps développé, conclusion).
    \item \textbf{Assurer} la \cle{cohésion} (connecteurs logiques, reprises anaphoriques, champs lexicaux) et la \cle{cohérence} (progression logique des idées, adaptation au destinataire, ton administratif) du texte.
    \item \textbf{Exploiter} les informations d'un document iconique (affiche : texte + image) pour étayer sa candidature.
    \item \textbf{S'auto-évaluer} et \textbf{co-évaluer} une production écrite à l'aide d'une grille critériée, puis \textbf{réviser} son texte en fonction des retours (remédiation).
\end{itemize}

\courssub{Situation}
\noindent \textbf{Contexte authentique :} La Commune d'\textbf{Arrondissement de Ngaoundéré 2\textsuperscript{e}} (Région de l'Adamaoua) lance, comme chaque année, son programme « \emph{Vacances Utiles - Jeunesse Technique} ». Ce programme offre des stages d'immersion professionnelle d'une durée de \textbf{quatre (4) semaines} (du 1\textsuperscript{er} au 31 juillet 2026) au sein des services techniques municipaux, à destination des élèves des lycées techniques publics et privés de l'arrondissement.

\medskip
\noindent \textbf{Document 1 : L'affiche officielle (support texte/image)}
\noindent L'affiche, placardée au tableau d'affichage de votre établissement (Lycée Technique de Ngaoundéré), présente les caractéristiques suivantes (description textuelle pour l'exercice) :

\begin{popfigure}
\begin{tikzpicture}[scale=0.9, every node/.style={font=\small}]
% Fond
\fill[popcream] (0,0) rectangle (14,10);
% Bandeau haut
\fill[popblue] (0,8.5) rectangle (14,10);
\node[white, align=center, font=\bfseries\Large] at (7,9.25) {COMMUNE D'ARRONDISSEMENT DE NGAOUNDÉRÉ 2\textsuperscript{e}};
\node[white, align=center, font=\bfseries] at (7,8.75) {PROGRAMME « VACANCES UTILES - JEUNESSE TECHNIQUE 2026 »};
% Logo (simulé)
\draw[popblue, thick] (1,6.5) circle (1.2);
\node[popblue] at (1,6.5) {\Huge $\star$};
% Corps de l'affiche
\node[align=left, text width=6cm, anchor=north west] at (2.8, 7.5) {
    \textbf{Services d'accueil :}\\
    \textbullet\ Service Urbanisme \& Topographie (CAO/DAO, levés)\\
    \textbullet\ Service Informatique \& Maintenance (Réseau, Hardware, Développement web)\\
    \textbullet\ Ateliers Municipaux : Mécanique auto, Électricité industrielle, Soudure, Menuiserie aluminium\\
};
\node[align=left, text width=5.5cm, anchor=north west] at (2.8, 5.5) {
    \textbf{Conditions d'éligibilité :}\\
    \textbullet\ Être élève de Première ou Terminale Technique\\
    \textbullet\ Âge : 16 à 20 ans\\
    \textbullet\ Aptitude physique (visite médicale fournie)\\
    \textbullet\ Dossier complet déposé avant le \textbf{15 juin 2026}
};
% Encadré Dossier
\draw[poporange, thick, rounded corners] (9.5, 7.5) rectangle (13.5, 4.5);
\node[poporange, align=center, font=\bfseries] at (11.5, 7.2) {DOSSIER DE CANDIDATURE};
\node[align=left, text width=3.5cm] at (11.5, 6.3) {
    1. Lettre de motivation adressée au Maire\\
    2. CV actualisé (photo + contacts)\\
    3. Copie bulletins (2 derniers trimestres)\\
    4. Attestation de scolarité 2025-2026\\
    5. Certificat médical de moins de 3 mois\\
    6. Photocopie CNI ou Acte de naissance
};
% Bas de l'affiche
\fill[popgreen] (0,0) rectangle (14,1.2);
\node[white, align=center, font=\bfseries] at (7, 0.6) {Dépôt : Secrétariat de la Mairie (Bureau 12) -- 7h30 - 15h30 -- Tél : 22 23 45 67};
\node[white, align=center, font=\small] at (7, 0.2) {Partenaires : MINJEC -- MINESEC -- ONG « Jeunesse Bâtisseuse »};
\end{tikzpicture}
\legende{Affiche simplifiée de l'appel à candidatures « Vacances Utiles 2026 » -- Mairie de Ngaoundéré 2\textsuperscript{e}.}
\end{popfigure}

\medskip
\noindent \textbf{Document 2 : Fiche de renseignements de l'élève (Votre profil)}
\noindent Vous êtes \textbf{TCHINDA Kevin}, élève en \textbf{Première Technicien Électricien (1re TEE)} au Lycée Technique de Ngaoundéré. Vous avez 17 ans (né le 12/03/2009 à Ngaoundéré). Vous résidez au quartier \textbf{Mardock}, BP 145 Ngaoundéré. Tél : 6 98 76 54 32. Email : kevin.tchinda@email.cm.
\noindent Vos motivations : Vous souhaitez mettre en pratique les modules \emph{Installation électrique bâtiment} et \emph{Dépannage maintenance} vus en classe. Vous visez le service « Ateliers Municipaux - Électricité industrielle » pour découvrir la maintenance préventive sur les groupes électrogènes et l'éclairage public. Votre moyenne générale au 2\textsuperscript{e} trimestre est de \textbf{14,5/20} (mention Assez Bien). Vous êtes titulaire du BEPC session 2023.

\medskip
\noindent \textbf{Document 3 : Modèle de lettre officielle étudié en classe (Rappel structure)}
\noindent Le modèle canonique présenté par l'enseignant respecte la norme administrative camerounaise :
\begin{enumerate}
    \item \textbf{Expéditeur} (Nom, Prénom, Adresse, Tél, Email) -- en haut à gauche.
    \item \textbf{Destinataire} (Titre, Fonction, Structure, Adresse) -- en haut à droite, sous l'expéditeur.
    \item \textbf{Lieu et Date} -- à droite, sous le destinataire.
    \item \textbf{Objet} : clair, concis, en gras, souligné -- aligné à gauche.
    \item \textbf{Référence(s)} (le cas échéant) -- sous l'objet.
    \item \textbf{Formule d'appel} (Monsieur le Maire,).
    \item \textbf{Corps de la lettre} : Introduction (source info, objet précis), Développement (profil, motivations, compétences, adéquation), Conclusion (disponibilité, demande d'entretien, formule de politesse développée).
    \item \textbf{Formule de politesse finale} (longue, administrative).
    \item \textbf{Signature} (manuscrite) précédée du nom tapé.
    \item \textbf{Mentions finales} : Pièces jointes (P.J. : n pièces), Copie à (si nécessaire).
\end{enumerate}

\courssub{Consignes}
\noindent Travaillez individuellement dans un premier temps, puis en binôme pour la phase d'évaluation. Répondez aux questions ci-dessous dans l'ordre, en justifiant chaque choix.

\begin{enumerate}
    \item \textbf{Analyse de la situation de communication (15 min)} \\
    À partir de l'affiche (Doc 1) et de votre profil (Doc 2), complétez le tableau d'analyse ci-dessous :
    \begin{center}
    \begin{tabularx}{\linewidth}{|l|X|}\hline
    \rowcolor{popblueL} \textbf{Composante} & \textbf{Votre réponse justifiée} \\ \hline
    Émetteur (Qui écrit ?) & \\ \hline
    Destinataire (À qui ? Titre exact) & \\ \hline
    Objet de la lettre (Quoi ?) & \\ \hline
    Finalité / But visé (Pourquoi ?) & \\ \hline
    Contraintes formelles (Délai, pièces, format) & \\ \hline
    Canal de transmission & \\ \hline
    Registre de langue / Ton & \\ \hline
    \end{tabularx}
    \end{center}

    \item \textbf{Analyse de l'affiche : Texte et Image (10 min)} \\
    L'affiche associe texte et éléments graphiques (logo, codes couleurs, mise en page).
    \begin{enumerate}
        \item Identifiez \textbf{deux} informations visuelles (couleur, forme, disposition) qui attirent l'attention sur les \emph{services d'accueil} et la \emph{date limite}.
        \item Quelle est la fonction du \textbf{logo} (l'étoile dans le cercle bleu) et de la \textbf{bande verte} en bas ? Justifiez par rapport à la communication institutionnelle.
    \end{enumerate}

    \item \textbf{Rédaction de la demande de stage -- Version 1 (Brouillon) (30 min)} \\
    Rédigez votre lettre complète sur feuille double. Respectez impérativement :
    \begin{itemize}
        \item La structure externe (10 éléments du modèle Doc 3).
        \item La structure interne en \textbf{trois paragraphes distincts} (Introduction, Corps, Conclusion) séparés par une ligne blanche.
        \item L'usage de \textbf{connecteurs logiques} variés (au moins 5 différents) pour assurer la cohésion.
        \item L'adaptation du vocabulaire technique (électricité, maintenance, sécurité) au service visé.
        \item Les formules de politesse administrative exactes (appel et finale).
    \end{itemize}

    \item \textbf{Évaluation par les pairs -- Grille critériée (20 min)} \\
    Échangez votre brouillon avec votre binôme. Évaluez la lettre de votre camarade en cochant le niveau atteint pour chaque critère. \textbf{Annotez} la marge pour signaler les forces (✓) et les faiblesses (✗ / ?).
    
    \begin{center}
    \begin{tabularx}{\linewidth}{|l|X|c|c|c|}\hline
    \rowcolor{popblueL} \textbf{Critères} & \textbf{Indicateurs observables} & \textbf{Non acquis (0)} & \textbf{En cours (1)} & \textbf{Acquis (2)} \\ \hline
    \textbf{Présentation \& Structure externe} & En-têtes complets, Objet clair, Date/Lieu, Formules exactes, Signature, P.J. & & & \\ \hline
    \textbf{Structure interne} & 3 paragraphes distincts (Intro/Corps/Concl), Progression logique & & & \\ \hline
    \textbf{Contenu \& Adaptation} & Profil élève cité, Service visé nommé, Compétences techniques liées aux modules, Motivation explicite & & & \\ \hline
    \textbf{Cohésion} & Connecteurs variés (au moins 5), Reprises (pronoms, synonymes), Ponctuation correcte & & & \\ \hline
    \textbf{Cohérence} & Ton administratif soutenu, Pas de familier, Argumentation crédible, Respect destinataire & & & \\ \hline
    \textbf{Langue (Orthographe/Grammaire)} & Accords, Temps (présent/futur), Ponctuation, Majuscules (Maire, Monsieur) & & & \\ \hline
    \end{tabularx}
    \end{center}
    \textbf{Note sur 12 :} \hrulefill \quad \textbf{Appréciation globale :} \hrulefill

    \item \textbf{Rédaction de la version finale -- Remédiation (15 min)} \\
    Intégrez les corrections de votre binôme et vos propres relectures. Produisez la \textbf{version définitive} (propre, sans ratures) qui sera remise à l'enseignant pour notation et affichage éventuel au « Tableau d'honneur ».
\end{enumerate}

\courssub{Ressources à mobiliser}
\noindent Rappel synthétique des outils linguistiques et méthodologiques vus dans la leçon.

\begin{aretenir}[Structure externe de la lettre administrative (Norme Cameroun)]
\textbf{Expéditeur} \\
Prénom NOM \\
Adresse complète \\
Tél / Email \\
\\
\textbf{Destinataire} \\
À Monsieur le Maire de l'Arrondissement de [Nom] \\
BP [Numéro] [Ville] \\
\\
\textbf{[Lieu], le [Jour] [Mois] [Année]} \\
\\
\textbf{Objet :} \textbf{\underline{Demande de stage de vacances au service [Nom du service]}} \\
\textbf{Réf. :} Votre avis d'appel à candidatures n°... / Affiche du [Date] \\
\\
\textbf{Monsieur le Maire,} \\
\\
\textbf{[Corps de la lettre : 3 paragraphes]} \\
\\
\textbf{En vous remerciant de l'attention que vous voudrez bien réserver à ma demande, je vous prie d'agréer, Monsieur le Maire, l'expression de ma considération distinguée.} \\
\\
\textbf{Signature} \\
Prénom NOM \\
\\
\textbf{P.J. :} 6 pièces (Lettre, CV, Bulletins, Attestation, Certificat médical, Pièce d'identité)
\end{aretenir}

\begin{aretenir}[Connecteurs logiques pour la cohésion (Exemples classés)]
\begin{center}
\begin{tabular}{|l|l|}\hline
\rowcolor{popblueL} \textbf{Fonction} & \textbf{Connecteurs (Registre administratif)} \\ \hline
Addition / Enchaînement & \textbf{En outre,} / \textbf{Par ailleurs,} / \textbf{De surcroît,} / \textbf{Également} \\ \hline
Cause / Explication & \textbf{En effet,} / \textbf{Du fait que,} / \textbf{Dans la mesure où,} / \textbf{Puisque} \\ \hline
Conséquence / Conclusion & \textbf{Par conséquent,} / \textbf{Aussi,} / \textbf{En définitive,} / \textbf{C'est pourquoi} \\ \hline
But / Finalité & \textbf{Afin de,} / \textbf{Dans le but de,} / \textbf{En vue de} + infinitif \\ \hline
Concession / Opposition & \textbf{Toutefois,} / \textbf{Néanmoins,} / \textbf{Cependant,} / \textbf{Malgré} \\ \hline
Illustration / Précision & \textbf{Namément,} / \textbf{Notamment,} / \textbf{En particulier,} / \textbf{C'est-à-dire} \\ \hline
\end{tabular}
\end{center}
\end{aretenir}

\begin{aretenir}[Formules de politesse administrative (Usage obligatoire)]
\begin{itemize}
    \item \textbf{Appel :} \textbf{Monsieur le Maire,} (seul, suivi d'une virgule, pas de « Cher »).
    \item \textbf{Finale (type long) :} \textbf{En vous remerciant de l'attention que vous voudrez bien réserver à ma demande, je vous prie d'agréer, Monsieur le Maire, l'expression de ma considération distinguée.}
    \item \textit{Variante (si réponse attendue) :} \textbf{Dans l'attente d'une réponse favorable, je vous prie d'agréer, Monsieur le Maire, l'assurance de ma très haute considération.}
\end{itemize}
\end{aretenir}

\begin{methode}[Rédaction du corps de la lettre : 3 étapes / 3 paragraphes]
\begin{enumerate}
    \item \textbf{Introduction (1 paragraphe, 3-4 lignes) :} \textbf{Source} (Où j'ai vu l'info : affiche mairie) + \textbf{Objet précis} (Quel stage, quel service, quelle période) + \textbf{Qualité} (Qui je suis : élève 1re TEE, établissement).
    \item \textbf{Développement / Corps (1 ou 2 paragraphes, 6-8 lignes) :} \textbf{Profil scolaire} (Série, modules forts, moyenne) $\rightarrow$ \textbf{Motivation technique} (Pourquoi ce service ? lien modules/activités service) $\rightarrow$ \textbf{Qualités personnelles} (Rigueur, sécurité, esprit d'équipe, disponibilité). \textit{Utiliser connecteurs : En effet, Par ailleurs, Notamment, Afin de.}
    \item \textbf{Conclusion (1 paragraphe, 3-4 lignes) :} \textbf{Disponibilité} (Dates, entretien) + \textbf{Demande formelle} (Bien vouloir accepter ma candidature) + \textbf{Formule de politesse finale}.
\end{enumerate}
\end{methode}

\courssub{Démarche et corrigé commenté}
\noindent Voici une proposition de résolution complète, étape par étape, servant de modèle d'excellence pour l'activité.

\begin{exoresolu}[Résolution guidée : De l'analyse à la version finale]
\textbf{Étape 1 : Analyse de la situation de communication (Consigne 1)} \\
\noindent Le tableau complété par l'élève (Kevin) :

\begin{center}
\begin{tabularx}{\linewidth}{|l|X|}\hline
\rowcolor{popblueL} \textbf{Composante} & \textbf{Réponse justifiée} \\ \hline
Émetteur & \textbf{TCHINDA Kevin}, élève de 1\textsuperscript{re} TEE, Lycée Technique de Ngaoundéré. (Justification : Doc 2, « Votre profil »). \\ \hline
Destinataire & \textbf{Monsieur le Maire de l'Arrondissement de Ngaoundéré 2\textsuperscript{e}}. (Justification : Affiche Doc 1, en-tête « COMMUNE D'ARRONDISSEMENT DE NGAOUNDÉRÉ 2\textsuperscript{e} » ; usage administratif : on écrit au décideur final). \\ \hline
Objet & \textbf{Demande de stage de vacances (1\textsuperscript{er} - 31 juillet 2026) au service Ateliers Municipaux - Électricité industrielle.} (Justification : Affiche précise services, dates ; profil élève vise électricité). \\ \hline
Finalité & Obtenir une convention de stage pour acquérir une expérience professionnelle en maintenance électrique industrielle (groupes électrogènes, éclairage public) et valider des compétences pratiques du cursus TEE. \\ \hline
Contraintes formelles & \textbullet\ Délai : Dépôt avant le \textbf{15 juin 2026} (Affiche).\\
& \textbullet\ Pièces : 6 pièces obligatoires listées dans l'encadré orange (Lettre, CV, Bulletins, Attestation, Certificat médical, CNI).\\
& \textbullet\ Format : Lettre officielle manuscrite ou dactylographiée, signée, structure externe/internes respectées.\\
& \textbullet\ Lieu dépôt : Secrétariat Mairie, Bureau 12 (7h30-15h30). \\ \hline
Canal & \textbf{Ecrit} (voie postale ou dépôt physique au Bureau 12). L'affiche est le support d'information initial (visuel). \\ \hline
Registre / Ton & \textbf{Administratif / Formel / Soutenu.} Vocabulaire technique (électricité), politesse codifiée, phrases complètes, pas d'argot, pas de familier. \\ \hline
\end{tabularx}
\end{center}

\textbf{Étape 2 : Analyse Texte/Image de l'affiche (Consigne 2)} \\
\begin{enumerate}
    \item \textbf{Services d'accueil :} Utilisation d'une \textbf{liste à puces} (textbullet) alignée à gauche dans une zone blanche, précédée d'un titre en \textbf{gras} sur fond crème. L'\textbf{icône « étoile » dans un cercle bleu} (logo stylisé) en marge gauche agit comme un \emph{point d'ancrage visuel} (ancrage spatial) qui guide l'œil vers le bloc texte principal. \textbf{Date limite :} Inscription dans un \textbf{encadré orange à coins arrondis} (contraste fort avec le fond crème), police \textbf{gras}, mention « \textbf{Avant le 15 juin 2026} » en majuscules/gras. Le code couleur \textbf{orange} (alerte/attention) signale l'urgence/délai.
    \item \textbf{Logo (étoile bleue) :} Fonction \textbf{identitaire} (symbole de la mairie/de l'excellence) et \textbf{attractive} (couleur institutionnelle bleue = confiance, institution). \textbf{Bande verte (bas) :} Fonction \textbf{informationnelle pratique} (coordonnées, horaires, partenaires). Le \textbf{vert} (couleur de l'espoir, de l'action, de la jeunesse) renforce le message positif « Vacances Utiles ». Elle ancre l'affiche dans le réel (téléphone, bureau).
\end{enumerate}

\textbf{Étape 3 : Rédaction Version 1 (Brouillon) -- Commentaire méthodologique} \\
\noindent L'élève applique la \textbf{Méthode des 3 paragraphes}. Voici le brouillon annoté (les annotations entre crochets \textcolor{popblue}{[...]} ne figurent pas dans la lettre finale).

\medskip
\noindent \textcolor{popmuted}{\textit{[En-têtes placés correctement : Expéditeur gauche, Destinataire droite, Lieu/Date droite, Objet gras souligné, Référence, Appel]}}

\medskip
\noindent \textbf{Introduction} \\
\textcolor{popmuted}{[Source + Objet + Qualité]} \\
Suite à l'affiche publiée par vos services annonçant le programme « Vacances Utiles 2026 », \textbf{par la présente}, je me permets de vous adresser ma candidature pour un stage de vacances d'une durée de quatre semaines, du 1\textsuperscript{er} au 31 juillet 2026, au sein du \textbf{service Ateliers Municipaux - Électricité industrielle}. Actuellement élève en \textbf{Première Technicien Électricien (1\textsuperscript{re} TEE)} au Lycée Technique de Ngaoundéré, je souhaite mettre à profit mes connaissances théoriques dans un cadre professionnel réel.

\medskip
\noindent \textbf{Corps de la lettre} \\
\textcolor{popmuted}{[Profil + Modules + Motivation technique + Qualités -- Connecteurs : En effet, Par ailleurs, Notamment, Afin de]} \\
\textbf{En effet,} ma formation en 1\textsuperscript{re} TEE m'a permis d'acquérir des compétences solides dans les modules d'\emph{Installation électrique bâtiment} et de \emph{Dépannage maintenance}, sanctionnés par une moyenne de \textbf{14,5/20} au deuxième trimestre. \textbf{Par ailleurs,} l'étude des normes de sécurité (NF C 15-100) et la pratique du câblage en atelier scolaire m'ont familiarisé avec la rigueur indispensable à ce métier. \\
\textbf{Notamment,} le service Électricité industrielle de votre mairie retient toute mon attention \textbf{car} il assure la maintenance préventive et curative des groupes électrogènes et de l'éclairage public. \textbf{Afin de} compléter ma formation, je désire vivement observer et participer, sous la tutelle de vos techniciens, aux opérations de diagnostic de pannes, de remplacement de composants et de vérification des protections différentielles sur site réel. \textbf{De surcroît,} mon stage d'initiation en classe de Seconde chez un électricien artisan (juillet 2024) m'a déjà confronté aux réalités du terrain : respect des consignes de sécurité, port des EPI, relation client.

\medskip
\noindent \textbf{Conclusion} \\
\textcolor{popmuted}{[Dispo + Demande + Formule finale]} \\
\textbf{C'est pourquoi,} fort de ma motivation et de ma disponibilité totale pendant la période indiquée, je sollicite votre bienveillance pour intégrer votre équipe. Je me tiens à votre entière disposition pour tout entretien que vous jugerez utile à l'appui de ma candidature. \\
\textcolor{popmuted}{[Formule finale standard]} \\
En vous remerciant de l'attention que vous voudrez bien réserver à ma demande, je vous prie d'agréer, Monsieur le Maire, l'expression de ma considération distinguée.

\medskip
\noindent \textcolor{popmuted}{[Signature + P.J.]}

\textbf{Étape 4 : Évaluation par les pairs -- Simulation (Consigne 4)} \\
\noindent Le binôme (ex: \textbf{NGONO Grace}, 1re TEE) remplit la grille pour Kevin :

\begin{center}
\begin{tabularx}{\linewidth}{|l|X|c|c|c|}\hline
\rowcolor{popblueL} \textbf{Critères} & \textbf{Indicateurs observables} & \textbf{0} & \textbf{1} & \textbf{2} \\ \hline
\textbf{Présentation \& Structure externe} & En-têtes complets, Objet clair, Date/Lieu, Formules exactes, Signature, P.J. & & & \cellcolor{popgreenL}\textbf{X} \\ \hline
\textbf{Structure interne} & 3 paragraphes distincts, Progression logique (Source $\rightarrow$ Preuves $\rightarrow$ Demande) & & & \cellcolor{popgreenL}\textbf{X} \\ \hline
\textbf{Contenu \& Adaptation} & Profil cité (1re TEE, 14,5), Service nommé (Élec. indus.), Modules cités (Install, Dépannage), Motivation claire & & & \cellcolor{popgreenL}\textbf{X} \\ \hline
\textbf{Cohésion} & Connecteurs : \textit{Par la présente, En effet, Par ailleurs, Notamment, Car, Afin de, De surcroît, C'est pourquoi} (8 connecteurs variés). Reprises : \textit{ma formation / elle / ce service / mon stage}. Ponctuation correcte. & & & \cellcolor{popgreenL}\textbf{X} \\ \hline
\textbf{Cohérence} & Ton soutenu (\textit{me permets, sollicite, bienveillance}). Vocabulaire technique précis (\textit{NF C 15-100, EPI, différentielles, groupes électrogènes}). Argumentation crédible (lien modules $\rightarrow$ service). & & \cellcolor{poporangeL}\textbf{X} & \\ \hline
\textbf{Langue} & Accords corrects. \textit{Attention :} « \textit{je me permets} » (accord) OK. « \textit{à votre entière disposition} » OK. Pas de fautes relevées. & & & \cellcolor{popgreenL}\textbf{X} \\ \hline
\end{tabularx}
\end{center}
\noindent \textbf{Note sur 12 :} \textbf{11/12} (Perte d'1 point sur Cohérence : la phrase « \textit{car il assure la maintenance...} » gagne à être reformulée en « \textit{dans la mesure où il assure...} » pour un style plus administratif). \\
\textbf{Appréciation :} \textit{Excellente lettre, structure parfaite, vocabulaire technique pertinent. Juste une légère lourdeur syntaxique au milieu du 2\textsuperscript{e} paragraphe. À afficher.}

\textbf{Étape 5 : Version finale -- Remédiation (Consigne 5)} \\
\noindent Kevin intègre la remarque (remplace « car » par « dans la mesure où »), vérifie l'orthographe des pièces jointes (6 pièces), s'assure que la date est le jour du dépôt (ex: \textbf{Ngaoundéré, le 10 juin 2026}) et produit la version propre ci-dessous (version finale à remettre).

\begin{center}
\fbox{\begin{minipage}{0.92\linewidth}
\textbf{TCHINDA Kevin} \\
Quartier Mardock, BP 145 Ngaoundéré \\
Tél : 6 98 76 54 32 \quad Email : kevin.tchinda@email.cm \\
\\
À Monsieur le Maire de l'Arrondissement de Ngaoundéré 2\textsuperscript{e} \\
BP 286 Ngaoundéré \\
\\
\textbf{Ngaoundéré, le 10 juin 2026} \\
\\
\textbf{Objet : \underline{Demande de stage de vacances (1\textsuperscript{er} - 31 juillet 2026) au service Ateliers Municipaux - Électricité industrielle}} \\
\textbf{Réf. :} Affiche programme « Vacances Utiles 2026 » affichée au Lycée Technique de Ngaoundéré. \\
\\
\textbf{Monsieur le Maire,} \\
\\
Suite à l'affiche publiée par vos services annonçant le programme « Vacances Utiles 2026 », \textbf{par la présente}, je me permets de vous adresser ma candidature pour un stage de vacances d'une durée de quatre semaines, du 1\textsuperscript{er} au 31 juillet 2026, au sein du \textbf{service Ateliers Municipaux - Électricité industrielle}. Actuellement élève en \textbf{Première Technicien Électricien (1\textsuperscript{re} TEE)} au Lycée Technique de Ngaoundéré, je souhaite mettre à profit mes connaissances théoriques dans un cadre professionnel réel. \\
\\
\textbf{En effet,} ma formation en 1\textsuperscript{re} TEE m'a permis d'acquérir des compétences solides dans les modules d'\emph{Installation électrique bâtiment} et de \emph{Dépannage maintenance}, sanctionnés par une moyenne de \textbf{14,5/20} au deuxième trimestre. \textbf{Par ailleurs,} l'étude des normes de sécurité (NF C 15-100) et la pratique du câblage en atelier scolaire m'ont familiarisé avec la rigueur indispensable à ce métier. \\
\textbf{Notamment,} le service Électricité industrielle de votre mairie retient toute mon attention \textbf{dans la mesure où} il assure la maintenance préventive et curative des groupes électrogènes et de l'éclairage public. \textbf{Afin de} compléter ma formation, je désire vivement observer et participer, sous la tutelle de vos techniciens, aux opérations de diagnostic de pannes, de remplacement de composants et de vérification des protections différentielles sur site réel. \textbf{De surcroît,} mon stage d'initiation en classe de Seconde chez un électricien artisan (juillet 2024) m'a déjà confronté aux réalités du terrain : respect des consignes de sécurité, port des EPI, relation client. \\
\\
\textbf{C'est pourquoi,} fort de ma motivation et de ma disponibilité totale pendant la période indiquée, je sollicite votre bienveillance pour intégrer votre équipe. Je me tiens à votre entière disposition pour tout entretien que vous jugerez utile à l'appui de ma candidature. \\
\\
En vous remerciant de l'attention que vous voudrez bien réserver à ma demande, je vous prie d'agréer, Monsieur le Maire, l'expression de ma considération distinguée. \\
\\
\textbf{Signature} \\
TCHINDA Kevin \\
\\
\textbf{P.J. :} 6 pièces (1. Lettre de motivation, 2. CV, 3. Bulletins 1\textsuperscript{er} et 2\textsuperscript{e} trimestres, 4. Attestation de scolarité 2025-2026, 5. Certificat médical du 05/06/2026, 6. Photocopie CNI).
\end{minipage}}
\end{center}
\end{exoresolu}


\courssub{Bilan de l'activité}
\noindent Cette activité d'intégration a permis de mettre en évidence les acquis majeurs suivants, conformément aux attendus du Probatoire technique :

\begin{enumerate}
    \item \textbf{Maîtrise de la situation de communication administrative :} L'élève a démontré sa capacité à décoder un appel à candidatures réel (affiche mixte texte/image), à en extraire les contraintes formelles (délai, pièces, destinataire exact) et à s'y conformer. L'analyse de l'image (codes couleurs, mise en page, logo) a renforcé la compétence « Communication par l'image ».
    \item \textbf{Respect rigoureux de la norme graphique et structurelle :} La distinction nette entre structure externe (norme administrative camerounaise : en-têtes, objet, formules codifiées) et structure interne (organisation tripartite Introduction/Corps/Conclusion) est acquise. L'élève sait que la forme \emph{est} du fond dans l'écrit professionnel.
    \item \textbf{Mobilisation de la cohésion et de la cohérence :} L'usage volontaire et varié des connecteurs logiques (\textit{en effet, par ailleurs, notamment, dans la mesure où, afin de, de surcroît, c'est pourquoi}) assure la fluidité syntaxique (cohésion). L'adaptation du registre (soutenu, technique, respectueux) et la progression argumentative (Profil $\rightarrow$ Compétences $\rightarrow$ Adéquation Service $\rightarrow$ Demande) garantissent la cohérence énonciative.
    \item \textbf{Compétence technique langagière :} L'injection de vocabulaire de spécialité (\textit{NF C 15-100, groupes électrogènes, protections différentielles, EPI, maintenance préventive/curative}) ancre la lettre dans la réalité de la série TEE, prouvant que l'élève ne récite pas un modèle vide mais projette son savoir-faire technique.
    \item \textbf{Démarche réflexive et remédiation :} L'évaluation par les pairs via une grille critériée explicite a objectivé la correction. Le passage de la version 1 (brouillon annoté) à la version finale (propre, corrigée) illustre concrètement le processus « rédaction $\rightarrow$ révision $\rightarrow$ publication » attendu en examen (Probatoire) et en milieu professionnel.
\end{enumerate}

\begin{attention}[Pièges évités grâce à l'activité]
\begin{itemize}
    \item \textbf{Confusion Maire / Directeur :} Écrire au Directeur des services techniques au lieu du Maire (décideur final).
    \item \textbf{Objet vague :} « Demande de stage » sans préciser le service, les dates, le programme.
    \item \textbf{Formules de politesse « scolaires » :} « Monsieur le Directeur, je vous salue » ou formules courtes type mail. L'exigence est la formule administrative longue.
    \item \textbf{Oubli des P.J. :} La liste numérotée des pièces jointes est obligatoire pour le traitement du dossier par le secrétariat.
    \item \textbf{Cohérence interne rompue :} Parler de mécanique alors qu'on postule en électricité, ou citer des modules inexistants en 1re TEE.
\end{itemize}
\end{attention}

\begin{savaistu}[Ancrage camerounais : Le stage « Vacances Utiles »]
Au Cameroun, le programme \textbf{« Vacances Utiles »}, piloté par le \textbf{MINJEC} (Ministère de la Jeunesse et de l'Éducation Civique) en partenariat avec le \textbf{MINESEC} et les collectivités territoriales décentralisées (CTD : Communes, Communautés urbaines), est un dispositif majeur d'insertion professionnelle des jeunes. Les mairies (comme Ngaoundéré 2\textsuperscript{e}, Douala 1\textsuperscript{er}, Yaoundé 3\textsuperscript{e}, etc.) publient annuellement des appels à candidatures (affiches, radio, réseaux sociaux) pour des stages dans leurs services techniques (voirie, éclairage, informatique, urbanisme, hygiène). La lettre de demande de stage est le \textbf{premier acte professionnel} de l'élève technicien. Sa qualité conditionne souvent l'obtention de la convention de stage, document indispensable pour valider l'année scolaire (module « Stage / PFMP » dans certaines spécialités) et construire son CV. Maîtriser cet écrit, c'est maîtriser sa première porte d'entrée dans le monde du travail.
\end{savaistu}

\newpage

\coursec{Méthodes et exercices résolus}

\coursec{Leçon 10 : La demande de stage}

\courssub{Méthode 1 : Assurer la cohésion et la cohérence d'un texte}

\begin{methode}[Rendre un texte cohérent et cohésif]
La \cle{cohérence} concerne le sens (logique des idées) ; la \cle{cohésion} concerne la forme (liens grammaticaux et lexicaux entre les phrases).
\begin{enumerate}
\item \textbf{Vérifier la cohérence} : une seule idée directrice ; les phrases s'enchaînent logiquement (cause, conséquence, opposition) ; pas de contradiction ni de hors-sujet.
\item \textbf{Assurer la cohésion grammaticale} : reprendre les noms par des \cle{pronoms} (il, elle, celui-ci, y, en) et accorder correctement ; maintenir la concordance des temps.
\item \textbf{Assurer la cohésion lexicale} : utiliser la reprise nominale, les synonymes, les champs lexicaux pour éviter les répétitions maladroites.
\item \textbf{Relier les phrases} par des \cle{connecteurs logiques} adaptés : addition (de plus, en outre), cause (car, parce que), conséquence (donc, c'est pourquoi), opposition (mais, cependant), chronologie (d'abord, ensuite, enfin).
\item \textbf{Relire} en vérifiant que chaque connecteur exprime bien la relation voulue.
\end{enumerate}
\end{methode}

\begin{exoresolu}[Remettre en cohérence un paragraphe]
\textbf{Énoncé.} Les phrases suivantes, extraites du brouillon d'un élève, ont été mélangées. Remettez-les dans l'ordre logique et améliorez la cohésion à l'aide de reprises et de connecteurs.
\begin{itemize}
\item[a.] Cette formation me passionne depuis le collège.
\item[b.] Je suis élève en classe de Première F3 au lycée technique de Douala.
\item[c.] C'est pourquoi je sollicite un stage dans votre entreprise.
\item[d.] J'y prépare un baccalauréat technique en électrotechnique.
\end{itemize}
\tcblower
\textbf{Solution détaillée.}
\begin{enumerate}
\item \textbf{Repérage de l'idée directrice} : il s'agit d'une présentation personnelle menant à une demande. Le texte doit aller du général (identité) au particulier (formation, motivation, demande).
\item \textbf{Ordonnancement logique} : b (identité) $\rightarrow$ d (précision sur la formation, introduite par « y ») $\rightarrow$ a (motivation) $\rightarrow$ c (conséquence : la demande). L'indice « y » de la phrase d suppose que b précède ; « C'est pourquoi » de la phrase c marque une conclusion, donc c ferme le paragraphe.
\item \textbf{Amélioration de la cohésion} : on relie a à la formation par une reprise nominale et un connecteur.
\end{enumerate}
\textbf{Texte réécrit :} « Je suis élève en classe de Première F3 au lycée technique de Douala. J'y prépare un baccalauréat technique en électrotechnique, \emph{une filière qui} me passionne depuis le collège. C'est pourquoi je sollicite un stage dans votre entreprise. »
\textbf{Conclusion :} l'ordre b--d--a--c, renforcé par la reprise pronominale « y », la reprise nominale « une filière » et le connecteur de conséquence « C'est pourquoi », rend le paragraphe cohérent et cohésif.
\end{exoresolu}

\courssub{Méthode 2 : Construire la structure externe de la lettre officielle}

\begin{methode}[Disposer les éléments externes d'une demande de stage]
La \cle{structure externe} est la présentation matérielle de la lettre. Elle obéit à un ordre fixe.
\begin{enumerate}
\item \textbf{En-tête (en haut à gauche)} : nom et prénom de l'expéditeur, classe, établissement, adresse, téléphone.
\item \textbf{Lieu et date (en haut à droite)} : « Douala, le 15 octobre 2026 ».
\item \textbf{Coordonnées du destinataire (à droite, sous la date)} : fonction (Monsieur le Directeur des ressources humaines), nom de l'entreprise, ville.
\item \textbf{Objet de la lettre (à gauche)} : « Objet : demande de stage académique » (une seule ligne, sans verbe conjugué).
\item \textbf{Formule d'appel} : « Monsieur le Directeur, » (reprendre le titre du destinataire, jamais « Cher Monsieur » dans une lettre officielle).
\item \textbf{Corps de la lettre} : paragraphes alignés avec alinéas.
\item \textbf{Formule de politesse} : « Veuillez agréer, Monsieur le Directeur, l'expression de ma considération distinguée. »
\item \textbf{Signature (en bas à droite)} : nom et prénom, paraphe.
\end{enumerate}
\end{methode}

\begin{exoresolu}[Compléter la structure externe d'une lettre]
\textbf{Énoncé.} Amina Etoundi, élève de Première G2 au lycée technique de Yaoundé, écrit au Directeur général de la SODECOTON à Garoua pour demander un stage. Reconstituez, dans l'ordre, les éléments externes de sa lettre (sans le corps).
\tcblower
\textbf{Solution détaillée.}
\begin{enumerate}
\item \textbf{En-tête à gauche} : Etoundi Amina -- Classe de Première G2 -- Lycée technique de Yaoundé -- BP 0000 Yaoundé -- Tél. : 6XX XX XX XX.
\item \textbf{Lieu et date à droite} : « Yaoundé, le 2 novembre 2026 » (la ville est celle de l'expéditeur, pas du destinataire).
\item \textbf{Destinataire à droite} : « Monsieur le Directeur général de la SODECOTON, Garoua ».
\item \textbf{Objet à gauche} : « Objet : demande de stage académique. »
\item \textbf{Appel} : « Monsieur le Directeur général, » (reprise exacte du titre).
\item \textbf{Formule de politesse} : « Veuillez agréer, Monsieur le Directeur général, l'expression de ma haute considération. »
\item \textbf{Signature à droite} : « Etoundi Amina ».
\end{enumerate}
\textbf{Conclusion :} la lettre présente les huit éléments externes dans l'ordre réglementaire ; elle est donc conforme aux normes de la correspondance officielle exigées au Probatoire technique.
\end{exoresolu}

\begin{aretenir}[Les 8 éléments de la structure externe]
\begin{center}
\begin{tabularx}{\linewidth}{|l|X|}\hline
\textbf{Position} & \textbf{Élément} \\ \hline
Haut gauche & En-tête expéditeur (nom, classe, établissement, contact) \\ \hline
Haut droit & Lieu, date \\ \hline
Sous la date (droit) & Destinataire (fonction, entreprise, ville) \\ \hline
Gauche & Objet (une ligne, nominale) \\ \hline
Gauche & Formule d'appel (titre seul) \\ \hline
Centre & Corps (paragraphes à alinéas) \\ \hline
Bas gauche & Formule de politesse développée (reprise du titre) \\ \hline
Bas droit & Signature (nom, prénom) \\ \hline
\end{tabularx}
\end{center}
\end{aretenir}

\courssub{Méthode 3 : Rédiger l'introduction et la conclusion de la demande}

\begin{methode}[Introduction et conclusion efficaces]
\begin{enumerate}
\item \textbf{Introduction (un paragraphe court)} : se présenter (nom, classe, établissement, filière) et annoncer clairement l'objet de la lettre : « J'ai l'honneur de solliciter auprès de votre structure un stage académique d'une durée de... ».
\item \textbf{Employer les formules protocolaires} : « J'ai l'honneur de... », « Je me permets de... », « Je viens très respectueusement... ».
\item \textbf{Conclusion} : remercier le destinataire de l'attention portée à la demande, exprimer l'espoir d'une suite favorable (« Dans l'attente d'une réponse favorable... »), puis enchaîner avec la formule de politesse.
\item \textbf{Accorder} la formule de politesse avec l'appel (même titre) et avec le genre de l'expéditeur si nécessaire.
\end{enumerate}
\end{methode}

\begin{exoresolu}[Rédiger introduction et conclusion]
\textbf{Énoncé.} Rédigez l'introduction et la conclusion de la demande de stage de Jean Kamga, élève de Première F4, qui sollicite un stage de quatre semaines à la Cameroon Tea Estates (CTE) de Buea.
\tcblower
\textbf{Solution détaillée.}
\begin{enumerate}
\item \textbf{Introduction} : présentation + annonce de l'objet, en deux phrases maximum.\\
« Je me nomme Kamga Jean, élève en classe de Première F4 au lycée technique de Buea, filière génie mécanique. J'ai l'honneur de solliciter auprès de votre entreprise un stage académique d'une durée de quatre semaines, du $1^{\text{er}}$ au 30 juin 2027. »
\item \textbf{Conclusion} : remerciement + attente + formule de politesse.\\
« Dans l'attente d'une suite que j'ose espérer favorable, je vous prie d'agréer, Monsieur le Directeur, l'expression de ma considération distinguée. »
\end{enumerate}
\textbf{Conclusion :} l'introduction identifie le demandeur et l'objet ; la conclusion remercie et formule la politesse attendue. Les deux paragraphes sont courts, corrects et protocolaires, comme l'exige la lettre officielle.
\end{exoresolu}

\courssub{Méthode 4 : Rédiger le corps de la lettre}

\begin{methode}[Construire le corps en trois paragraphes]
\begin{enumerate}
\item \textbf{Paragraphe 1 -- la formation} : filière suivie, compétences acquises (cours, travaux pratiques, logiciels ou machines maîtrisés).
\item \textbf{Paragraphe 2 -- la motivation} : pourquoi cette entreprise précisément (réputation, domaine d'activité, adéquation avec la filière) ; ce que le stage apportera au demandeur.
\item \textbf{Paragraphe 3 -- les disponibilités et pièces jointes} : période et durée du stage, documents joints (CV, attestation de scolarité), engagement à respecter le règlement de l'entreprise.
\item \textbf{Style} : phrases courtes, ton respectueux, vocabulaire précis, connecteurs logiques (d'abord, en outre, enfin), aucune faute d'orthographe.
\end{enumerate}
\end{methode}

\begin{exoresolu}[Rédiger le corps d'une demande de stage]
\textbf{Énoncé.} Rédigez le corps (trois paragraphes) de la demande de stage de Marie Ngo Bell, élève de Première G1 (secrétariat-bureautique), auprès de la Camtel de Douala, pour un stage d'un mois.
\tcblower
\textbf{Solution détaillée.}
\begin{enumerate}
\item \textbf{Formation} : « Actuellement en classe de Première G1, option secrétariat-bureautique, au lycée technique de Douala, j'ai acquis de solides compétences en dactylographie, en traitement de texte et en classement des dossiers. Je maîtrise également les logiciels Word et Excel. »
\item \textbf{Motivation} : « Votre entreprise, leader des télécommunications au Cameroun, offre un cadre idéal pour mettre en pratique ces connaissances. Effectuer mon stage au sein de vos services administratifs me permettrait de découvrir le fonctionnement réel d'un secrétariat professionnel. »
\item \textbf{Disponibilités} : « Disponible du $1^{\text{er}}$ au 31 juillet 2027, je m'engage à respecter scrupuleusement le règlement intérieur de votre structure. Vous trouverez ci-joints mon curriculum vitae et une attestation de scolarité. »
\end{enumerate}
\textbf{Conclusion :} le corps suit le plan formation -- motivation -- disponibilités ; chaque paragraphe développe une seule idée, reliée aux autres par des connecteurs, ce qui garantit clarté et cohérence.
\end{exoresolu}

\courssub{Méthode 5 : Lire un texte fonctionnel et un texte explicatif}

\begin{methode}[Analyser un texte fonctionnel ou explicatif]
\begin{enumerate}
\item \textbf{Identifier la nature du texte} : un \cle{texte fonctionnel} (lettre, affiche, mode d'emploi, formulaire, convention) vise une action ; un \cle{texte explicatif} vise à faire comprendre (définitions, causes, étapes d'un processus).
\item \textbf{Repérer le contexte} : qui écrit ? à qui ? dans quel but ?
\item \textbf{Dégager les informations essentielles} : pour une lettre, l'objet et la demande ; pour un texte explicatif, la notion expliquée et les étapes de l'explication.
\item \textbf{Repérer les marques de l'explication} : connecteurs (car, en effet, c'est-à-dire, par conséquent), définitions, exemples, comparaisons, vocabulaire technique défini.
\item \textbf{Répondre aux questions} en citant brièvement le texte pour justifier chaque réponse.
\end{enumerate}
\end{methode}

\begin{exoresolu}[Questions sur un texte explicatif]
\textbf{Énoncé.} Lisez l'extrait suivant, puis répondez aux questions.\\
« Le stage académique est une période d'immersion de l'élève en milieu professionnel. Il permet en effet de confronter les connaissances théoriques acquises en classe à la pratique de l'atelier ou du bureau. C'est pourquoi il constitue une étape obligatoire de la formation technique : il prépare le futur technicien aux exigences réelles du monde du travail. »
\begin{itemize}
\item[a.] Quelle est la nature de ce texte ? Justifiez.
\item[b.] Quelle définition l'auteur donne-t-il du stage ?
\item[c.] Relevez deux connecteurs logiques et précisez la relation qu'ils expriment.
\end{itemize}
\tcblower
\textbf{Solution détaillée.}
\begin{enumerate}
\item[a.] \textbf{Nature du texte} : c'est un texte \emph{explicatif}. Justification : l'auteur définit une notion (« Le stage académique est... »), puis en explique l'utilité et la nécessité ; il ne raconte pas, il ne convainc pas, il fait comprendre.
\item[b.] \textbf{Définition} : « une période d'immersion de l'élève en milieu professionnel ». C'est une définition par le verbe « être », typique du texte explicatif.
\item[c.] \textbf{Connecteurs} : « en effet » exprime la \emph{confirmation/explication} (il justifie l'intérêt du stage) ; « c'est pourquoi » exprime la \emph{conséquence} (le stage est utile, donc il est obligatoire).
\end{enumerate}
\textbf{Conclusion :} le texte explique le rôle du stage grâce à une définition claire et à des connecteurs qui enchaînent logiquement les idées ; toutes les réponses sont appuyées par des citations du texte.
\end{exoresolu}

\begin{exoresolu}[Analyse d'un texte fonctionnel : offre de stage]
\textbf{Énoncé.} Lisez l'avis suivant, affiché au lycée technique de Ngaoundéré, puis répondez aux questions.\\
« \textbf{AVIS DE RECRUTEMENT DE STAGIAIRES} \\
La Société de Développement du Coton (SODECOTON) informe les élèves des classes de Première et Terminale des filières Génie Mécanique, Électrotechnique et Maintenance Industrielle qu'elle offre 15 places de stage académique pour la période du 1$^{\text{er}}$ juillet au 31 août 2027. \\
Les dossiers de candidature (lettre de motivation, CV, attestation de scolarité, photocopie CNI) sont à déposer au secrétariat de la Direction des Ressources Humaines, usine de Garoua, au plus tard le 15 mai 2027. \\
Pour tout renseignement, contacter le 222 00 00 00. »
\begin{itemize}
\item[a.] Identifiez la nature et la fonction de ce texte.
\item[b.] Quelles sont les filières concernées ?
\item[c.] Relevez la date limite de dépôt et la durée du stage.
\item[d.] Quels sont les quatre documents constitutifs du dossier ?
\end{itemize}
\tcblower
\textbf{Solution détaillée.}
\begin{enumerate}
\item[a.] \textbf{Nature} : texte \emph{fonctionnel} (avis de recrutement / annonce administrative). \textbf{Fonction} : \emph{informer} et \emph{inciter à l'action} (déposer un dossier).
\item[b.] \textbf{Filières} : Génie Mécanique, Électrotechnique, Maintenance Industrielle.
\item[c.] \textbf{Date limite} : 15 mai 2027. \textbf{Durée} : 1$^{\text{er}}$ juillet au 31 août 2027 (deux mois).
\item[d.] \textbf{Pièces du dossier} : lettre de motivation, CV, attestation de scolarité, photocopie de la CNI.
\end{enumerate}
\textbf{Conclusion :} la lecture active d'un texte fonctionnel permet d'extraire rapidement les informations opérationnelles (qui, quoi, quand, comment) nécessaires pour agir.
\end{exoresolu}

\courssub{Méthode 6 : Analyser la communication par l'image}

\begin{methode}[Lire une image (affiche, logo, photographie)]
\begin{enumerate}
\item \textbf{Décrire objectivement} : éléments visuels (personnages, objets, couleurs, disposition), texte accompagnant l'image (slogan, titre).
\item \textbf{Identifier le code de communication} : l'image \cle{dénote} (montre) et \cle{connote} (suggère) ; relever ce que les couleurs, les gestes, les symboles évoquent.
\item \textbf{Déterminer la fonction} : informer, sensibiliser, persuader (publicité), instruire.
\item \textbf{Dégager le message} : à qui s'adresse l'image ? Quelle réaction cherche-t-elle à provoquer ?
\item \textbf{Relier image et texte} : montrer comment le slogan ou la légende renforce le visuel.
\end{enumerate}
\end{methode}

\begin{exoresolu}[Analyse d'une affiche de sensibilisation]
\textbf{Énoncé.} Une affiche du ministère de la Santé publique montre une mère camerounaise souriante faisant vacciner son bébé, avec le slogan : « Un enfant vacciné, un avenir protégé. » Décrivez cette image, dégagez sa fonction et son message.
\tcblower
\textbf{Solution détaillée.}
\begin{enumerate}
\item \textbf{Description (dénotation)} : l'affiche représente une mère souriante tenant son bébé pendant qu'un personnel de santé administre un vaccin ; le slogan « Un enfant vacciné, un avenir protégé » figure en gros caractères.
\item \textbf{Connotation} : le sourire de la mère suggère la confiance et l'absence de danger ; le bébé symbolise l'avenir ; le personnel soignant rassure sur le sérieux de l'acte médical.
\item \textbf{Fonction} : l'image \emph{sensibilise} et \emph{persuade} : elle cherche à convaincre les parents de faire vacciner leurs enfants.
\item \textbf{Message et cible} : destinée aux parents, elle associe vaccination et protection de l'avenir ; le slogan, construit sur un parallélisme (« un enfant... un avenir... »), renforce l'idée que le vaccin est un investissement pour demain.
\end{enumerate}
\textbf{Conclusion :} par la combinaison d'une image rassurante et d'un slogan équilibré, l'affiche communique efficacement un message de santé publique : vacciner, c'est protéger.
\end{exoresolu}

\begin{savaistu}[Le stage en entreprise dans le système technique camerounais]
Au Cameroun, le stage académique (ou stage d'initiation) est une \cle{obligation pédagogique} pour l'obtention du Baccalauréat technique et du Probatoire. Il est régi par les textes du MINESEC (notamment la circulaire annuelle d'organisation des examens). Sa durée varie selon les spécialités (généralement 4 à 8 semaines). L'élève est encadré par un \cle{maître de stage} en entreprise et un \cle{professeur tuteur} au lycée. À l'issue du stage, un \cle{rapport de stage} (souvent soutenu oralement) est évalué et compte pour la certification. La convention de stage, signée par le proviseur, le directeur d'entreprise et l'élève (ou ses parents), couvre la responsabilité civile de l'élève.
\end{savaistu}

\courssub{Méthode 7 : Confronter, améliorer et remédier}

\begin{methode}[Améliorer une production après confrontation]
\begin{enumerate}
\item \textbf{Confronter} : échanger les lettres entre pairs et les comparer à la grille de correction (structure externe complète ? plan du corps respecté ? formules protocolaires ?).
\item \textbf{Repérer les écarts} : élément manquant (objet, date, signature), faute de grammaire ou d'orthographe, maladresse de style, connecteur inadapté.
\item \textbf{Corriger et justifier} : pour chaque correction, nommer la règle appliquée (accord, concordance des temps, norme de la lettre officielle).
\item \textbf{Réécrire} la version définitive au propre, en intégrant toutes les corrections.
\item \textbf{Remédier} : refaire un exercice ciblé sur le point faible identifié (accords, formules de politesse, connecteurs).
\end{enumerate}
\end{methode}

\begin{exoresolu}[Corriger un extrait de lettre fautive]
\textbf{Énoncé.} Voici l'extrait fautif d'une demande de stage. Relevez quatre erreurs, corrigez-les et justifiez chaque correction.\\
« Cher Monsieur le Directeur, je vous écris pour que vous me donnez un stage. Je suis disponible à tout moment et je peut commencer quand vous voulez. Salutations distinguées. »
\tcblower
\textbf{Solution détaillée.}
\begin{enumerate}
\item \textbf{« Cher Monsieur le Directeur »} $\rightarrow$ « Monsieur le Directeur ». Justification : dans une lettre officielle, l'appel ne comporte jamais « Cher », réservé aux lettres amicales.
\item \textbf{« pour que vous me donnez »} $\rightarrow$ « pour que vous me donniez ». Justification : la locution « pour que » exprimant le but est suivie du \emph{subjonctif}.
\item \textbf{« je peut »} $\rightarrow$ « je peux ». Justification : à la première personne du singulier du présent, le verbe \emph{pouvoir} s'écrit « je peux » (le « t » n'apparaît qu'à la troisième personne : « il peut »).
\item \textbf{« Salutations distinguées »} $\rightarrow$ « Veuillez agréer, Monsieur le Directeur, l'expression de ma considération distinguée. » Justification : la formule de politesse d'une lettre officielle est développée et reprend le titre du destinataire.
\end{enumerate}
\textbf{Version corrigée :} « Monsieur le Directeur, j'ai l'honneur de solliciter auprès de votre entreprise un stage académique. Disponible à tout moment, je peux commencer dès que vous le jugerez utile. Veuillez agréer, Monsieur le Directeur, l'expression de ma considération distinguée. »
\textbf{Conclusion :} la confrontation à la grille a permis de repérer une erreur de registre, une faute de mode, une faute de conjugaison et une formule de politesse non conforme ; la version réécrite respecte toutes les normes de la lettre officielle.
\end{exoresolu}

\begin{attention}[Les 5 erreurs qui coûtent des points au Bac]
\begin{enumerate}
\item \textbf{Oublier un élément de la structure externe} : objet, date, formule d'appel ou signature manquants font perdre des points de présentation. Vérifiez les huit éléments avant de rendre la copie.
\item \textbf{Écrire « Cher Monsieur » ou « Salutations » dans une lettre officielle} : ces formules appartiennent à la correspondance familière. Utilisez « Monsieur le Directeur, » et « Veuillez agréer... l'expression de ma considération distinguée ».
\item \textbf{Confondre cohésion et cohérence} : la cohérence porte sur la logique des idées, la cohésion sur les liens de langue (reprises, connecteurs). Ne répondez pas à côté de la question posée.
\item \textbf{Négliger la concordance des temps et les modes} : « pour que » exige le subjonctif ; le récit de la formation se fait au présent ou au passé composé, sans mélange incohérent.
\item \textbf{Décrire une image sans dégager son message} : en communication par l'image, la description seule ne suffit pas ; il faut toujours conclure sur la fonction (informer, persuader, sensibiliser) et le message visé.
\end{enumerate}
\end{attention}

\newpage

\coursec{Exercices}

\courssub{Je vérifie mes connaissances}

\begin{exercice}[QCM : la demande de stage]
Pour chaque question, entoure la bonne réponse.
\begin{enumerate}
\item Une demande de stage est :
\begin{enumerate}
\item une lettre amicale ;
\item une lettre officielle ;
\item un texte littéraire.
\end{enumerate}
\item Dans la structure externe d'une lettre officielle, on trouve en haut à gauche :
\begin{enumerate}
\item le nom et l'adresse de l'expéditeur ;
\item la formule de politesse ;
\item la signature.
\end{enumerate}
\item La formule de politesse finale s'adressant à un directeur d'entreprise commence par :
\begin{enumerate}
\item « Salut monsieur » ;
\item « Veuillez agréer, Monsieur le Directeur, ... » ;
\item « À bientôt ».
\end{enumerate}
\item La cohésion d'un texte est assurée principalement par :
\begin{enumerate}
\item les connecteurs logiques et les reprises anaphoriques ;
\item la longueur des phrases ;
\item le nombre de paragraphes.
\end{enumerate}
\item L'objet de la lettre (« Objet : demande de stage ») se place :
\begin{enumerate}
\item après la formule de politesse finale ;
\item avant la formule d'appel ;
\item dans la signature.
\end{enumerate}
\end{enumerate}
\end{exercice}

\begin{exercice}[Vrai ou faux, en justifiant]
Dis si chaque affirmation est vraie ou fausse, puis justifie ta réponse en une phrase.
\begin{enumerate}
\item Dans une demande de stage, on peut tutoyer le directeur si l'on est pressé.
\item L'objet de la lettre se place avant la formule d'appel (« Monsieur le Directeur, ... »).
\item Un texte cohérent est un texte dont les idées s'enchaînent logiquement.
\item La communication par l'image ne peut pas transmettre d'informations pratiques.
\end{enumerate}
\end{exercice}

\begin{exercice}[Définitions]
Définis en deux ou trois phrases chacune des notions suivantes, en donnant un exemple pour chacune :
\begin{enumerate}
\item la \cle{cohésion} du texte ;
\item la \cle{cohérence} du texte ;
\item la \cle{dénotation} et la \cle{connotation} dans une image ;
\item le \cle{texte explicatif}.
\end{enumerate}
\end{exercice}

\begin{exercice}[Reconnaître les éléments de la structure externe]
Voici des éléments extraits d'une lettre officielle :\\
a) « Garoua, le 12 mars 2026 » ; b) « Monsieur le Directeur Général de la SODECOTON » ; c) « Objet : demande de stage académique » ; d) « L'élève Mbarga Alain, classe de Première F3 » ; e) « Veuillez agréer, Monsieur le Directeur Général, l'expression de ma haute considération. »\\
Pour chaque élément, indique son nom (lieu et date, destinataire, objet, expéditeur, formule de politesse) et sa place dans la lettre (haut à gauche, haut à droite, avant le corps, après le corps).
\end{exercice}

\courssub{Je m'entraîne}

\begin{exercice}[Remettre en ordre le corps d'une lettre]
Voici, dans le désordre, les cinq phrases du corps d'une demande de stage. Remets-les dans l'ordre logique et justifie ton choix en relevant les connecteurs et les reprises anaphoriques.
\begin{enumerate}
\item « C'est pourquoi je sollicite un stage d'observation et de pratique dans votre entreprise pour la période du 1\textsuperscript{er} juillet au 31 août 2026. »
\item « En effet, votre société, spécialisée dans la transformation du coton, offre un cadre idéal pour appliquer les connaissances acquises en classe. »
\item « Je reste à votre disposition pour tout renseignement complémentaire et pour un entretien à votre convenance. »
\item « Actuellement élève en classe de Première F3 au Lycée technique de Garoua, je prépare le Baccalauréat technique en mécanique. »
\item « De plus, ce stage me permettra de découvrir le monde du travail et de confirmer mon projet professionnel. »
\end{enumerate}
\end{exercice}

\begin{exercice}[Corriger les fautes de registre]
Un élève a écrit à un directeur d'entreprise. Releve les cinq fautes de registre ou de politesse dans les passages suivants et propose pour chacune une correction adaptée à une lettre officielle.
\begin{enumerate}
\item « Salut le patron ! »
\item « Je veux absolument faire mon stage chez vous, c'est obligé. »
\item « Tu peux m'envoyer les documents par WhatsApp. »
\item « Réponds-moi vite, c'est urgent. »
\item « Bisous, et à la prochaine ! »
\end{enumerate}
\end{exercice}

\begin{exercice}[Analyser une affiche de recrutement de stagiaires]
Une entreprise agroalimentaire de Douala publie une affiche intitulée « DEVIENS STAGIAIRE CHEZ NOUS ! ». On y voit un jeune technicien souriant en tenue de travail, devant une machine moderne, avec le slogan « Apprends un métier, construis ton avenir » et, en bas, les mentions : « Stages de 2 mois --- juillet-août --- envoyer la demande avant le 15 juin au siège de l'entreprise, Bonabéri ».
\begin{enumerate}
\item Releve deux éléments de \cle{dénotation} (ce que l'image montre) et deux éléments de \cle{connotation} (ce qu'elle suggère).
\item Quel message l'affiche veut-elle faire passer ? À quel public s'adresse-t-elle ?
\item Dégage de cette affiche trois informations pratiques utiles à la rédaction de ta demande de stage (durée, date limite, destinataire/lieu).
\end{enumerate}
\end{exercice}

\begin{exercice}[Relever les procédés explicatifs]
Lis le texte fonctionnel suivant :\\
« Le stage académique est une période de formation pratique en entreprise, c'est-à-dire un temps pendant lequel l'élève applique sur le terrain les notions apprises en classe. Il se déroule en trois étapes : d'abord l'observation, ensuite la participation aux tâches simples, enfin la réalisation d'un travail sous contrôle. Par exemple, un élève de Première F3 peut d'abord regarder un tourneur-fraiseur, puis régler une machine, enfin usiner une pièce simple. Ainsi, le stage permet de vérifier l'adéquation entre la formation et le métier. »
\begin{enumerate}
\item Releve trois procédés explicatifs du texte (définition, reformulation, énumération, exemple, conclusion).
\item Explique le rôle de chacun dans la clarté du message.
\end{enumerate}
\end{exercice}

\courssub{Je me prépare au Bac}

\begin{exercice}[Sujet type Probatoire : la demande de stage complète]
\textbf{Situation.} Tu es élève en classe de Première technique au Lycée technique de Garoua. Tu lis l'annonce suivante : « La SODECOTON de Garoua (BP 105, Garoua) offre des stages académiques du 1\textsuperscript{er} juillet au 31 août 2026 aux élèves des classes de Première et Terminale techniques. Les demandes, adressées à Monsieur le Directeur Général, doivent parvenir avant le 30 avril 2026. »
\begin{enumerate}
\item Analyse la situation : qui écrit ? à qui ? pourquoi ? quelles informations doivent figurer dans la lettre ?
\item Rédige la demande de stage complète en respectant la structure externe (coordonnées de l'expéditeur, lieu et date, destinataire, objet, formule d'appel, formule de politesse, signature) et la structure interne (introduction : présentation et objet ; corps : motivations, qualités, période souhaitée ; conclusion : remerciements et disponibilité).
\item Soigne la cohésion de ta lettre en employant au moins quatre connecteurs logiques différents.
\end{enumerate}
\textbf{Barème indicatif :} présentation (structure externe) : 4 pts ; cohérence et organisation du contenu : 6 pts ; correction de la langue : 6 pts ; registre soutenu et formules de politesse : 4 pts.
\end{exercice}

\begin{methode}[Rappel : les étapes de la rédaction]
\begin{enumerate}
\item \textbf{Analyser la situation} : identifier l'expéditeur, le destinataire, l'objet, la période, la date limite.
\item \textbf{Rédiger l'introduction} : se présenter (nom, classe, établissement) et annoncer l'objet de la demande.
\item \textbf{Rédiger le corps} : exposer les motivations, les qualités, la période souhaitée, en liant les phrases par des connecteurs (d'abord, en effet, de plus, c'est pourquoi).
\item \textbf{Rédiger la conclusion} : remercier, se déclarer disponible, exprimer l'espoir d'une réponse favorable.
\item \textbf{Relire et corriger} : vérifier la structure externe, l'orthographe, le registre soutenu.
\end{enumerate}
\end{methode}

\begin{exercice}[Sujet type Bac : introduction et conclusion d'une demande de stage]
\textbf{Situation imposée.} Tu es élève en Première F4 (génie civil) au Lycée technique de Douala. L'entreprise CAMROUTES, chargée de l'entretien des routes de la ville, accepte des stagiaires pendant les grandes vacances. Tu rédiges une demande de stage adressée à Monsieur le Directeur des Ressources Humaines.
\begin{enumerate}
\item Rédige uniquement l'\textbf{introduction} de ta lettre : présente-toi (nom, classe, établissement) et annonce clairement l'objet de ta demande (5 à 8 lignes).
\item Rédige uniquement la \textbf{conclusion} : exprime ta disponibilité, remercie le destinataire et formule l'espoir d'une réponse favorable (4 à 6 lignes).
\item Échange ta production avec un camarade : à l'aide d'une grille (clarté de l'objet, registre soutenu, correction de la langue, connecteurs), évalue sa copie et propose deux améliorations.
\item Réécris ton introduction ou ta conclusion en tenant compte des remarques reçues.
\end{enumerate}
\end{exercice}

\begin{exercice}[Sujet type Bac : texte fonctionnel, cohésion et réécriture]
\textbf{Texte.} « Pour obtenir un stage, il faut d'abord rédiger une demande claire et correcte. En effet, la lettre est le premier contact entre l'élève et l'entreprise : elle doit donc donner une bonne image du candidat. D'une part, la présentation doit être soignée ; d'autre part, le contenu doit être précis : l'objet, la période du stage et les motivations de l'élève. Cependant, beaucoup de demandes sont refusées à cause des fautes d'orthographe ou d'un ton trop familier. Enfin, il est conseillé de joindre une attestation de scolarité. En conclusion, une demande bien rédigée augmente les chances d'être accepté. »
\begin{enumerate}
\item Releve dans le texte quatre connecteurs logiques et précise la relation qu'ils expriment (cause, addition, opposition, conclusion).
\item Releve deux reprises anaphoriques (pronoms ou mots de reprise) et indique ce qu'elles remplacent.
\item Le texte est-il cohérent ? Justifie ta réponse en montrant l'enchaînement des idées.
\item Un élève a écrit : « Je demande un stage. Mon école est loin. J'aime le football. Envoyez-moi une réponse. » Explique pourquoi ce texte manque de cohésion et de cohérence, puis réécris-le en un paragraphe correct de quatre phrases liées par des connecteurs.
\item À partir des conseils du texte, rédige le \textbf{corps} d'une demande de stage (10 à 12 lignes) pour un stage dans un garage automobile de ton quartier.
\end{enumerate}
\end{exercice}

\begin{attention}[Pièges fréquents à l'examen]
\begin{itemize}
\item Ne pas confondre expéditeur (en haut à gauche) et destinataire (en haut à droite).
\item Ne jamais tutoyer ni employer de formules familières (« Salut », « Bisous ») dans une lettre officielle.
\item Ne pas oublier l'objet, placé avant la formule d'appel.
\item Vérifier que chaque phrase du corps est liée à la précédente par un connecteur ou une reprise anaphorique : c'est la cohésion qui garantit la cohérence.
\end{itemize}
\end{attention}

\newpage

\coursec{Corrigés des exercices}

\coursec{Leçon 10 : La demande de stage — Corrigés détaillés}

\courssub{Corrigés de « Je vérifie mes connaissances »}

\begin{corrige}[Exercice 1 — QCM : la demande de stage]
\begin{enumerate}
\item \textbf{b)} Une demande de stage est une \cle{lettre officielle} : elle est adressée à un responsable d'entreprise ou d'administration, dans un cadre formel, avec un objet précis.
\item \textbf{a)} En haut à gauche figurent le nom, la classe et l'adresse de l'\cle{expéditeur} (celui qui écrit la lettre).
\item \textbf{b)} La formule de politesse soutenue commence par « Veuillez agréer, Monsieur le Directeur, ... ». Les formules « Salut » et « À bientôt » relèvent du registre familier, interdit ici.
\item \textbf{a)} La \cle{cohésion} repose sur les connecteurs logiques (d'abord, ensuite, en effet, c'est pourquoi) et les reprises anaphoriques (pronoms, mots de reprise).
\item \textbf{b)} L'objet se place \emph{avant} la formule d'appel, sous les coordonnées du destinataire : il annonce le sujet de la lettre.
\end{enumerate}
\end{corrige}

\begin{corrige}[Exercice 2 — Vrai ou faux, en justifiant]
\begin{enumerate}
\item \textbf{Faux.} Une lettre officielle exige le vouvoiement et le registre soutenu, quelles que soient les circonstances : l'urgence ne justifie jamais la familiarité.
\item \textbf{Vrai.} L'objet annonce le sujet de la lettre ; il se place avant la formule d'appel (« Monsieur le Directeur, ... »), sous les coordonnées du destinataire.
\item \textbf{Vrai.} La \cle{cohérence} est l'enchaînement logique des idées : chaque phrase prolonge la précédente sans contradiction ni hors-sujet.
\item \textbf{Faux.} Une image (affiche, schéma, photographie) peut transmettre des informations pratiques : dates, lieux, durées, consignes, comme le montre l'affiche de recrutement de stagiaires étudiée dans la leçon.
\end{enumerate}
\end{corrige}

\begin{corrige}[Exercice 3 — Définitions]
\begin{enumerate}
\item \textbf{La cohésion} est l'ensemble des liens formels qui relient les phrases d'un texte entre elles : connecteurs logiques, reprises anaphoriques (pronoms, mots de reprise), champs lexicaux. \emph{Exemple :} « Je prépare le Baccalauréat technique. \textbf{C'est pourquoi} je sollicite un stage dans votre entreprise. » Le connecteur « c'est pourquoi » lie les deux phrases.
\item \textbf{La cohérence} est l'enchaînement logique des idées : le texte forme un tout sensé, sans contradiction ni digression. \emph{Exemple :} dans une demande de stage, on se présente, on expose ses motivations, puis on formule la demande : l'ordre des idées est logique.
\item \textbf{La dénotation} désigne ce que l'image montre objectivement (les éléments visibles) ; \textbf{la connotation} désigne ce qu'elle suggère, les idées ou les sentiments qu'elle évoque. \emph{Exemple :} un technicien souriant devant une machine (dénotation) suggère la réussite et le dynamisme (connotation).
\item \textbf{Le texte explicatif} est un texte qui donne à comprendre une notion, un phénomène ou une démarche. Il utilise des procédés comme la définition, la reformulation (« c'est-à-dire »), l'énumération (« d'abord, ensuite, enfin »), l'exemple (« par exemple ») et la conclusion (« ainsi »). \emph{Exemple :} « Le stage académique est une période de formation pratique en entreprise, c'est-à-dire un temps pendant lequel l'élève applique les notions apprises en classe. »
\end{enumerate}
\end{corrige}

\begin{corrige}[Exercice 4 — Reconnaître les éléments de la structure externe]
\begin{center}
\begin{tabularx}{\linewidth}{|c|X|X|}
\hline
\rowcolor{popblueL} Élément & Nom & Place dans la lettre \\
\hline
a) « Garoua, le 12 mars 2026 » & Lieu et date & En haut à droite, sous les coordonnées du destinataire \\
\hline
b) « Monsieur le Directeur Général de la SODECOTON » & Destinataire & En haut à droite \\
\hline
c) « Objet : demande de stage académique » & Objet & Avant le corps, sous le destinataire \\
\hline
d) « L'élève Mbarga Alain, classe de Première F3 » & Expéditeur & En haut à gauche \\
\hline
e) « Veuillez agréer, ... ma haute considération. » & Formule de politesse & Après le corps, avant la signature \\
\hline
\end{tabularx}
\end{center}
\textbf{Point de vigilance :} ne pas inverser expéditeur (haut à gauche) et destinataire (haut à droite) ; c'est l'erreur la plus fréquente à l'examen.
\end{corrige}

\courssub{Corrigés de « Je m'entraîne »}

\begin{corrige}[Exercice 5 — Remettre en ordre le corps d'une lettre]
Ordre logique : \textbf{d -- b -- a -- e -- c}.
\begin{enumerate}
\item \textbf{d} ouvre le corps : l'élève se présente (classe, établissement, diplôme préparé). C'est l'information de départ.
\item \textbf{b} justifie le choix de l'entreprise : le connecteur « \textbf{En effet} » introduit une explication qui prolonge la présentation.
\item \textbf{a} formule la demande : « \textbf{C'est pourquoi} » exprime la conséquence logique de ce qui précède.
\item \textbf{e} ajoute un argument : « \textbf{De plus} » marque l'addition ; « \textbf{ce stage} » reprend anaphoriquement « un stage » de la phrase précédente.
\item \textbf{c} conclut le corps par la formule de disponibilité, qui précède naturellement la formule de politesse finale.
\end{enumerate}
La cohésion est donc assurée par trois connecteurs (en effet, c'est pourquoi, de plus) et une reprise anaphorique (ce stage), ce qui garantit la cohérence de l'ensemble.
\end{corrige}

\begin{corrige}[Exercice 6 — Corriger les fautes de registre]
\begin{enumerate}
\item « Salut le patron ! » : salutation familière et irrespectueuse $\rightarrow$ « \textbf{Monsieur le Directeur,} » (formule d'appel).
\item « Je veux absolument..., c'est obligé » : ton impératif et affirmatif $\rightarrow$ « \textbf{Je souhaiterais effectuer mon stage au sein de votre entreprise.} » (conditionnel de politesse).
\item « Tu peux m'envoyer les documents par WhatsApp » : tutoiement et moyen informel $\rightarrow$ « \textbf{Je vous serais reconnaissant de bien vouloir me faire parvenir les documents nécessaires.} »
\item « Réponds-moi vite, c'est urgent » : impératif et familiarité $\rightarrow$ « \textbf{Dans l'attente de votre réponse, je vous prie d'agréer, Monsieur le Directeur, ...} »
\item « Bisous, et à la prochaine ! » : clôture familière $\rightarrow$ « \textbf{Veuillez agréer, Monsieur le Directeur, l'expression de ma haute considération.} »
\end{enumerate}
\textbf{Point de vigilance :} à l'examen, toute formule familière fait perdre des points sur le critère « registre soutenu et formules de politesse ».
\end{corrige}

\begin{corrige}[Exercice 7 — Analyser une affiche de recrutement de stagiaires]
\begin{enumerate}
\item \textbf{Dénotation} (ce que l'image montre) : un jeune technicien souriant en tenue de travail ; une machine moderne ; les mentions écrites (durée, date limite, lieu). \textbf{Connotation} (ce qu'elle suggère) : le sourire évoque la réussite et l'épanouissement au travail ; la machine moderne suggère une entreprise dynamique et bien équipée, où l'on apprend un vrai métier.
\item \textbf{Message} : l'affiche veut convaincre les jeunes qu'un stage dans cette entreprise permet d'apprendre un métier et de préparer son avenir (slogan : « Apprends un métier, construis ton avenir »). \textbf{Public visé} : les élèves des lycées techniques, en particulier ceux qui cherchent un stage académique.
\item \textbf{Trois informations pratiques utiles à la demande de stage} :
\begin{itemize}
\item la \textbf{durée} du stage : 2 mois, juillet-août (à reprendre comme période souhaitée) ;
\item la \textbf{date limite} d'envoi : le 15 juin (la lettre doit être écrite et postée avant cette date) ;
\item le \textbf{destinataire et le lieu} : le siège de l'entreprise, Bonabéri (coordonnées à placer en haut à droite de la lettre).
\end{itemize}
\end{enumerate}
\end{corrige}

\begin{corrige}[Exercice 8 — Relever les procédés explicatifs]
\begin{enumerate}
\item \textbf{Définition} : « Le stage académique est une période de formation pratique en entreprise ». Rôle : donner le sens précis du terme dès le début, pour que le lecteur sache de quoi l'on parle.
\item \textbf{Reformulation} : « c'est-à-dire un temps pendant lequel l'élève applique sur le terrain les notions apprises en classe ». Rôle : reprendre la définition en termes plus simples et éviter tout malentendu.
\item \textbf{Énumération} : « d'abord l'observation, ensuite la participation aux tâches simples, enfin la réalisation d'un travail sous contrôle ». Rôle : ordonner les étapes et rendre l'explication facile à suivre.
\item \textbf{Exemple} : « Par exemple, un élève de Première F3 peut d'abord regarder un tourneur-fraiseur... ». Rôle : rendre la notion concrète et proche du lecteur.
\item \textbf{Conclusion} : « Ainsi, le stage permet de vérifier l'adéquation entre la formation et le métier ». Rôle : tirer la leçon générale de l'explication et clore le raisonnement.
\end{enumerate}
\end{corrige}

\courssub{Corrigés de « Je me prépare au Probatoire »}

\begin{corrige}[Exercice 9 — Sujet type Probatoire : la demande de stage complète]
\textbf{1. Analyse de la situation.}
\begin{itemize}
\item \textbf{Qui écrit ?} Un élève de Première technique du Lycée technique de Garoua (l'expéditeur).
\item \textbf{À qui ?} Monsieur le Directeur Général de la SODECOTON, BP 105, Garoua (le destinataire).
\item \textbf{Pourquoi ?} Pour solliciter un stage académique du 1\textsuperscript{er} juillet au 31 août 2026.
\item \textbf{Informations obligatoires} : coordonnées de l'expéditeur, lieu et date (avant le 30 avril 2026), coordonnées du destinataire, objet, formule d'appel, présentation, motivations, période souhaitée, remerciements, formule de politesse, signature.
\end{itemize}

\textbf{2. Modèle de lettre.}

\begin{center}
\begin{tabularx}{\linewidth}{X X}
L'élève Mbarga Alain & Monsieur le Directeur Général \\
Classe de Première F3 & de la SODECOTON \\
Lycée technique de Garoua & BP 105, Garoua \\
BP 32, Garoua & \\
& Garoua, le 12 mars 2026 \\
\end{tabularx}
\end{center}

\textbf{Objet :} demande de stage académique.

Monsieur le Directeur Général,

Je me nomme Mbarga Alain, élève en classe de Première F3 au Lycée technique de Garoua, où je prépare le Baccalauréat technique. Je me permets de solliciter un stage académique au sein de votre entreprise pour la période du 1\textsuperscript{er} juillet au 31 août 2026.

En effet, votre société, spécialisée dans la transformation du coton, offre un cadre idéal pour appliquer sur le terrain les connaissances acquises en classe. De plus, ce stage me permettra de découvrir le monde du travail et de confirmer mon projet professionnel. Sérieux, ponctuel et désireux d'apprendre, je m'engage à respecter le règlement intérieur de votre entreprise. C'est pourquoi je vous serais très reconnaissant de bien vouloir examiner favorablement ma demande.

Dans l'attente de votre réponse, je reste à votre disposition pour tout renseignement complémentaire et pour un entretien à votre convenance.

Veuillez agréer, Monsieur le Directeur Général, l'expression de ma haute considération.

\begin{flushright}L'élève Mbarga Alain\end{flushright}

\textbf{3. Connecteurs employés} (au moins quatre) : « En effet » (explication), « De plus » (addition), « C'est pourquoi » (conséquence), « Dans l'attente de » (transition vers la conclusion).

\textbf{Barème indicatif (sur 20) :}
\begin{itemize}
\item Présentation et structure externe complète (expéditeur, destinataire, lieu/date, objet, formules, signature) : 4 pts ;
\item Cohérence et organisation du contenu (introduction, corps, conclusion ; connecteurs) : 6 pts ;
\item Correction de la langue (orthographe, grammaire, ponctuation) : 6 pts ;
\item Registre soutenu et formules de politesse : 4 pts.
\end{itemize}
\textbf{Points de vigilance :} date de rédaction antérieure au 30 avril 2026 ; période du stage conforme à l'annonce ; vouvoiement constant ; aucune formule familière.
\end{corrige}

\begin{corrige}[Exercice 10 — Sujet type Probatoire (entraînement Baccalauréat) : introduction et conclusion]
\textbf{1. Modèle d'introduction.}

Monsieur le Directeur des Ressources Humaines,

Je me nomme Ngo Bell Sandrine, élève en classe de Première F4 (génie civil) au Lycée technique de Douala, où je prépare le Baccalauréat technique. Informée que votre entreprise, chargée de l'entretien des routes de la ville, accueille des stagiaires pendant les grandes vacances, je me permets de solliciter un stage académique au sein de vos services.

\textit{Attendus :} nom et prénom, classe et spécialité, établissement, annonce claire de l'objet (demande de stage). Le ton est soutenu, les phrases sont liées.

\textbf{2. Modèle de conclusion.}

Dans l'attente d'une réponse que j'espère favorable, je reste à votre entière disposition pour tout renseignement complémentaire et pour un entretien à votre convenance. Je vous remercie par avance de l'attention que vous voudrez bien accorder à ma demande.

Veuillez agréer, Monsieur le Directeur des Ressources Humaines, l'expression de ma haute considération.

\textit{Attendus :} disponibilité, remerciement, espoir d'une réponse favorable, formule de politesse adaptée au titre exact du destinataire.

\textbf{3. Grille d'évaluation entre pairs (à appliquer à la copie du camarade) :}
\begin{center}
\begin{tabularx}{\linewidth}{|X|c|}
\hline
\rowcolor{popblueL} Critère & Note /4 \\
\hline
Clarté de l'objet de la demande & \\
\hline
Registre soutenu (vouvoiement, aucune familiarité) & \\
\hline
Correction de la langue (orthographe, grammaire) & \\
\hline
Présence de connecteurs et cohésion & \\
\hline
\end{tabularx}
\end{center}
Exemples d'améliorations proposables : remplacer « je veux » par « je souhaiterais » ; ajouter un connecteur (« en effet », « c'est pourquoi ») entre deux phrases juxtaposées.

\textbf{4. Réécriture :} reprendre sa production en intégrant les deux remarques reçues, puis la relire à voix haute pour vérifier l'enchaînement des idées.

\textbf{Points de vigilance :} la formule de politesse doit reprendre le titre exact du destinataire (« Monsieur le Directeur des Ressources Humaines ») ; l'introduction annonce l'objet sans raconter toute sa vie ; la conclusion ne contient aucune information nouvelle.
\end{corrige}

\begin{corrige}[Exercice 11 — Sujet type Probatoire (entraînement Baccalauréat) : texte fonctionnel, cohésion et réécriture]
\textbf{1. Connecteurs logiques.}
\begin{itemize}
\item « \textbf{En effet} » : relation de cause/explication ;
\item « \textbf{D'une part... d'autre part} » : relation d'addition/énumération ;
\item « \textbf{Cependant} » : relation d'opposition ;
\item « \textbf{En conclusion} » : relation de conclusion. (On peut ajouter « Enfin » : énumération ; « donc » : conséquence.)
\end{itemize}

\textbf{2. Reprises anaphoriques.}
\begin{itemize}
\item « \textbf{elle} » (dans « elle doit donc donner une bonne image ») reprend « \textbf{la lettre} » ;
\item « \textbf{beaucoup de demandes} » reprend l'idée de « \textbf{une demande claire et correcte} ».
\end{itemize}

\textbf{3. Cohérence du texte.} Le texte est cohérent : il part du conseil principal (rédiger une demande claire et correcte), l'explique (la lettre est le premier contact avec l'entreprise), le détaille (présentation soignée et contenu précis), signale les erreurs à éviter (fautes d'orthographe, ton familier), ajoute un conseil (joindre une attestation de scolarité) et conclut (une demande bien rédigée augmente les chances d'acceptation). Chaque phrase prolonge logiquement la précédente.

\textbf{4. Analyse et réécriture.} Le texte de l'élève manque de \textbf{cohésion} : les phrases sont juxtaposées, sans aucun connecteur ni reprise anaphorique. Il manque aussi de \textbf{cohérence} : les idées (« mon école est loin », « j'aime le football ») n'ont aucun rapport logique avec la demande de stage.\\
Réécriture possible : « Je sollicite un stage dans votre garage. En effet, je prépare le Baccalauréat technique et je souhaite appliquer mes connaissances en mécanique. De plus, ce stage me permettra de découvrir concrètement le métier de garagiste. C'est pourquoi je vous serais reconnaissant de bien vouloir m'adresser une réponse favorable. »

\textbf{5. Modèle de corps de demande de stage (garage automobile).}

Actuellement élève en classe de Première technique au Lycée technique de [Ville], je prépare le Baccalauréat technique en mécanique. C'est pourquoi je sollicite un stage d'observation et de pratique dans votre garage pour la période des grandes vacances. En effet, votre atelier, réputé pour la qualité de son travail, offre un cadre idéal pour appliquer sur le terrain les notions apprises en classe. D'une part, ce stage me permettra de participer aux tâches simples de l'entretien automobile ; d'autre part, il me fera découvrir les exigences réelles du métier. De plus, je suis sérieux, ponctuel et désireux d'apprendre. Enfin, je joins à ma demande une attestation de scolarité, comme il est conseillé. Je reste donc à votre disposition pour un entretien à votre convenance.

\textit{Connecteurs utilisés :} c'est pourquoi, en effet, d'une part... d'autre part, de plus, enfin, donc.

\textbf{Barème indicatif (sur 20) :} questions de compréhension (1 à 4) : 8 pts ; réécriture du paragraphe : 4 pts ; corps de la lettre (cohésion, contenu, langue) : 8 pts.

\textbf{Points de vigilance :} citer les connecteurs \emph{entre guillemets} et nommer précisément la relation logique ; dans la réécriture, supprimer toute idée hors sujet ; dans le corps de la lettre, respecter la période et le destinataire de la situation imposée.
\end{corrige}

\newpage

\coursec{Fiche bilan}

\begin{popfigure}
\begin{tikzpicture}[mindmap, grow cyclic, every node/.style={concept, font=\small},
  root concept/.append style={concept color=popblue, font=\bfseries\small, minimum size=2.6cm},
  level 1/.append style={level distance=3.4cm, sibling angle=51},
  level 1 concept/.append style={font=\footnotesize, minimum size=1.9cm}]
\node[root concept]{La demande de stage}
  child[concept color=poporange]{ node {Cohésion et cohérence : reprises, connecteurs, thème unique} }
  child[concept color=popgreen]{ node {Structure externe : en-têtes, objet, formules, signature} }
  child[concept color=poppurple]{ node {Structure interne : introduction, corps, conclusion} }
  child[concept color=poppink]{ node {Textes fonctionnels : lire, repérer les informations utiles} }
  child[concept color=popgold]{ node {Texte explicatif : expliquer, définir, illustrer} }
  child[concept color=popblueL]{ node {Communication par l'image : dénoter, connoter, interpréter} }
  child[concept color=popmuted]{ node {Rédaction : analyser, rédiger, confronter, améliorer} };
\end{tikzpicture}
\legende{Carte mentale de la leçon : les sept savoirs articulés autour de la rédaction d'une lettre officielle de demande de stage.}
\end{popfigure}

\begin{bilan}
\textbf{1. Cohésion et cohérence du texte}
\begin{itemize}
  \item \cle{Cohésion} : liens visibles entre les éléments du texte (reprises anaphoriques : pronoms, synonymes ; connecteurs : d'abord, ensuite, donc, cependant ; reprises du temps verbal).
  \item \cle{Cohérence} : unité de sens du texte (un seul thème, une progression logique, aucune contradiction).
  \item Un bon texte est à la fois \textbf{cohérent} (sens) et \textbf{cohésif} (forme).
\end{itemize}

\textbf{2. La lettre officielle de demande de stage}
\begin{center}
\begin{tabularx}{\linewidth}{|l|X|}\hline
\rowcolor{popblueL} \textbf{Partie} & \textbf{Contenu attendu} \\ \hline
En-tête & Coordonnées de l'expéditeur (en haut à gauche), lieu et date (à droite), coordonnées du destinataire (à droite) \\ \hline
Objet & Formule brève : \emph{Objet : demande de stage} \\ \hline
Formule d'appel & \emph{Monsieur le Directeur,} / \emph{Madame,} \\ \hline
Introduction & Se présenter brièvement et annoncer la demande \\ \hline
Corps & Motiver la demande : formation suivie, période et durée du stage, qualités, intérêt pour l'entreprise \\ \hline
Conclusion & Remercier et exprimer l'espoir d'une réponse favorable \\ \hline
Formule de politesse & \emph{Veuillez agréer, Monsieur le Directeur, l'expression de ma considération distinguée.} \\ \hline
Signature & Nom, prénom et signature en bas à droite \\ \hline
\end{tabularx}
\end{center}

\textbf{3. Texte explicatif et communication par l'image}
\begin{itemize}
  \item Le \cle{texte explicatif} rend compte d'un phénomène : définitions, exemples, comparaisons, connecteurs logiques (car, parce que, en effet, c'est-à-dire).
  \item Lire une image : \cle{dénotation} (ce qu'on voit) puis \cle{connotation} (ce que l'image suggère), cadrage, couleurs, message visé.
\end{itemize}

\textbf{4. Méthode express de rédaction}
\begin{enumerate}
  \item Analyser la situation : à qui ? pourquoi ? quand ? quelles informations ?
  \item Rédiger un plan en trois parties (introduction, corps, conclusion).
  \item Rédiger en veillant à la cohésion (connecteurs) et à la correction (orthographe, accords).
  \item Confronter sa lettre à celle d'un camarade, puis améliorer et recopier au propre.
\end{enumerate}
\end{bilan}

\begin{aretenir}[Checklist avant le Bac]
\begin{itemize}
  \item $\square$ Je sais distinguer cohésion (liens formels) et cohérence (unité de sens).
  \item $\square$ Je connais au moins cinq connecteurs logiques et je sais les employer.
  \item $\square$ Je place correctement les coordonnées de l'expéditeur, du destinataire, le lieu et la date.
  \item $\square$ Je rédige un objet bref et une formule d'appel adaptée au destinataire.
  \item $\square$ Mon introduction présente qui je suis et annonce clairement ma demande.
  \item $\square$ Mon corps de lettre motive la demande : formation, période, durée, qualités.
  \item $\square$ Ma conclusion remercie et formule l'espoir d'une suite favorable.
  \item $\square$ Je connais une formule de politesse administrative complète.
  \item $\square$ Je sais analyser un texte fonctionnel : nature, émetteur, destinataire, informations utiles.
  \item $\square$ Je sais décrire et interpréter une image (dénotation puis connotation).
  \item $\square$ Je relis ma lettre pour corriger orthographe, accords et ponctuation avant de la rendre.
\end{itemize}
\end{aretenir}

\findecours

\end{document}
